\documentclass[11pt,a4paper,oneside,openany]{book}

\usepackage[utf8]{inputenc}
\usepackage[T1]{fontenc}
\usepackage[english]{babel}
\usepackage{csquotes}
\usepackage[charter]{mathdesign}
\usepackage{amsmath}
\usepackage{lipsum}
\usepackage{amsthm}
\usepackage{mathtools}
\usepackage{physics}
\usepackage{graphicx}
\usepackage{enumitem}
\usepackage{tikz-cd}
\usepackage{geometry}

% Page header management
\usepackage{fancyhdr}
\pagestyle{fancy}
\fancyhf{} 
\fancyhead[L]{\slshape\nouppercase{\leftmark}} 
\fancyhead[R]{\thepage} 
\renewcommand{\headrulewidth}{0pt}

% Redefinition of the plain style for chapter starting pages
\fancypagestyle{plain}{
	\fancyhf{}
	\renewcommand{\headrulewidth}{0pt}
	\fancyfoot[R]{\thepage} 
}

\usepackage[backend=biber, sorting=none, style=alphabetic]{biblatex}
\addbibresource{bib.bib}
\usepackage[colorlinks=true,linkcolor=black,citecolor=black,urlcolor=black]{hyperref}

\usepackage{epigraph}
\setlength{\epigraphwidth}{0.5\textwidth}
\renewcommand{\epigraphflush}{flushright}
\renewcommand{\textflush}{flushright}

\newtheorem{theorem}{Theorem}[chapter]
\newtheorem{definition}{Definition}[section]
\newtheorem{lemma}[definition]{Lemma}
\newtheorem{proposition}[definition]{Proposition}
\newtheorem{corollary}[definition]{Corollary}
\theoremstyle{remark}
\newtheorem{remark}[definition]{Remark}
\newtheorem{example}[definition]{Example}

\let\mathcal\relax
\newcommand{\mathcal}[1]{\text{\usefont{OMS}{cmsy}{m}{n}#1}}
\DeclareMathOperator{\Spec}{Spec}
\DeclareMathOperator{\Ad}{Ad}
\DeclareMathOperator{\ad}{ad}
\newcommand{\poisson}[2]{\{#1,#2\}}
% ============================================================
% Comandi Matematici
% ============================================================

% --- Insiemi numerici ---
\newcommand{\R}{\mathbb{R}}
\newcommand{\C}{\mathbb{C}}
\newcommand{\N}{\mathbb{N}}
\newcommand{\K}{\mathbb{K}}
\newcommand{\A}{\mathfrak{A}}
\newcommand{\B}{\mathcal{B}}

% --- Spazio di Hilbert e spazio proiettivo ---
% ATTENZIONE: \H è già un comando di LaTeX (accento ungherese, es. \H{o}=ő).
% Il renewcommand lo sovrascrive: nella tesi non ti servirà mai quell'accento,
% quindi è sicuro, ma se preferisci non toccare comandi di sistema usa \Hil
% al posto di \H ovunque sotto.
\renewcommand{\H}{\mathcal{H}}
\newcommand{\PH}{\mathbb{P}\H}                    % P H, spazio proiettivo

% --- Algebre C* ---
\newcommand{\Cstar}{C^*}                          % da usare in math mode: $\Cstar$
\newcommand{\Csalg}{$C^*$-algebra}                % da usare nel testo: "è una \Csalg\ commutativa"
\newcommand{\Csalgs}{$C^*$-algebre}               % plurale
\newcommand{\BH}{\mathcal{B}(\H)}                           % operatori limitati su H
\newcommand{\Cz}[1]{C_0(#1)}                      % C_0(M)
\newcommand{\Cinf}[1]{C^\infty(#1)}               % C^infty(M)
\newcommand{\Cinfc}[1]{C^\infty_c(#1)}            % C^infty_c(M)
\newcommand{\Mat}[2]{M_{#1}(#2)}                  % M_n(C), es. \Mat{2}{\C}

% --- Algebre di Lie (fraktur generico) ---
\newcommand{\mf}[1]{\mathfrak{#1}}                % \mf{g}, \mf{su}(2), \mf{u}(n), \mf{X}(M)

\newcommand{\qcomm}[2]{\dfrac{1}{i\hbar}[#1,#2]}  % (1/ih)[A,B]  -- condizione di Dirac, Qubit IV

% --- Notazione bra-ket ---
\newcommand{\melem}[3]{\langle#1|#2|#3\rangle}    % elemento di matrice <psi|A|psi>
\newcommand{\expect}[2]{\langle#1\rangle_{#2}}    % <A>_psi

% --- Campo vettoriale hamiltoniano ---
\newcommand{\Ham}[1]{X_{#1}}                      % \Ham{f}, \Ham{f_H}, \Ham{f_A}

% --- Simboli e operatori vari ---
\newcommand{\Id}{\mathbb{1}}
\renewcommand{\Re}{\mathrm{Re}}                   % la tesi usa Re/Im in tondo, non ℜ/ℑ
\renewcommand{\Im}{\mathrm{Im}}
\newcommand{\LC}{\varepsilon_{ijk}}               % simbolo di Levi-Civita (indici fissi ijk)
\DeclareMathOperator{\Ran}{Ran}
\DeclareMathOperator{\Ker}{Ker}
\DeclareMathOperator{\Span}{span}
\DeclareMathOperator{\sing}{sing}
\newgeometry{margin=2cm}
\setlength{\headwidth}{\textwidth}

\begin{document}
	\setcounter{chapter}{5}

\chapter{Recovery of the Dirac Equation}
\label{ch:dirac}

This chapter addresses a question on states. The pure states of a system are the points of the projective Hilbert space $\mathbb{P}\mathcal{H}$ of Chapter 3, with its Kähler structure, and the spacetime symmetries are an action $\Phi$ of the Poincar\'e group $\mathcal{P}$ on $\mathbb{P}\mathcal{H}$ by Kähler isomorphisms. Nothing else is given, so that no momentum, no position, no spin and no wave equation is available at the outset. The task is to produce explicit models of the pair $(\mathbb{P}\mathcal{H},\Phi)$, \textit{i.e.}, to describe the states and the action by concrete objects, and to determine which features of these models are forced by the symmetry. Chapter~\ref{ch:analytical-hamiltonians} followed the same strategy for the Galilei group, and Theorem~\ref{thm:free-hamiltonian} showed there that covariance alone fixes the Hamiltonian of an elementary system up to an additive constant. Here the strategy is applied to $\mathcal{P}$, and it yields more: the mass shell, the energy as a function of the momentum, the spin as a label and the equation of motion are all outputs.

The chain consists of eight steps, and each of them replaces a model of the states by an equivalent one. The abstract space $\mathbb{P}\mathcal{H}$ is exchanged first for the rays of square integrable functions on the mass shell, the set of values that the four-momentum can take, with values in a spin space, and then for the rays of the solutions of a first order equation. Equivalence means a Kähler isomorphism intertwining the two actions, so that $\mathbb{P}\mathcal{H}$ never leaves the stage and every model describes the same states with the same transition probabilities. The first replacement uses the symmetry, irreducibility and positivity of the energy. The second is of a different nature, since it requires that the transformation law of the fields be local, \textit{i.e.}, given by matrices independent of the momentum, and compatible with space inversion. These are requirements on the description of the states and not properties of the system, and for this reason the second replacement rests on more hypotheses than the first. The hypotheses are collected in Definition~\ref{def:hypotheses}, the main result is Theorem~\ref{thm:main}, and Table~\ref{tab:steps} indicates where each step is carried out.

\begin{definition}[Hypotheses of the chapter]\label{def:hypotheses}
	Let $\mathcal{H}$ be a separable Hilbert space and let $\mathbb{P}\mathcal{H}$ be its projective space with the Kähler structure of Chapter 3. The hypotheses on the pair $(\mathbb{P}\mathcal{H},\Phi)$ and on its description are the following.
	\begin{enumerate}
		\item[\textbf{H1}] \emph{Symmetry.} $\Phi$ is a homomorphism from $\mathcal{P}$ to the Kähler isomorphisms of $\mathbb{P}\mathcal{H}$ with continuous orbit maps, as per Definition~\ref{def:poincare-action}.
		\item[\textbf{H2}] \emph{Irreducibility.} $\Phi$ is irreducible, as per Definition~\ref{def:poincare-irreducible}.
		\item[\textbf{H3}] \emph{Positive energy.} The energy operator of the lift $\Pi$ of $\Phi$, as per Definition~\ref{def:energy-operator}, is bounded below, as per Definition~\ref{def:positive-energy}.
		\item[\textbf{H4}] \emph{Locality.} $(\mathbb{P}\mathcal{H},\Phi)$ admits a realization, as per Definition~\ref{def:realization}, by fields $\Psi(p)\in F(p)\subseteq W$ with $W$ a finite dimensional complex vector space, on which $\mathcal{P}$ acts through fixed matrices $D(g)$ not required to be unitary, with $F(\mathring{p})=\iota(V)$, as per Definition~\ref{def:local-realization}.
		\item[\textbf{H5}] \emph{Parity.} The realization of H4 carries a parity matrix $M$ with $M^2=\mathbb{1}$, $M D(g)M^{-1}=D(\theta(g))$ and $MF(\mathring{p})=F(\mathring{p})$, as per Definition~\ref{def:parity-realization}.
		\item[\textbf{H6}] \emph{Minimality.} $\dim W$ is minimal among the realizations satisfying H4 and H5, as per Definitions~\ref{def:local-realization} and~\ref{def:parity-realization}.
	\end{enumerate}
\end{definition}

\begin{remark}\label{rem:hypotheses-nature}
	Hypotheses H1 to H3 concern the pair $(\mathbb{P}\mathcal{H},\Phi)$, and they define an elementary relativistic system as per Definition~\ref{def:elementary-system}. Hypotheses H4 to H6 are requirements on a realization of the pair and not properties of the system, because by Remark~\ref{rem:ph-no-spin} the space $\mathbb{P}\mathcal{H}$ does not distinguish the spin, and the same states admit realizations by fields of different type. The following choices are conventions and not hypotheses: $c=1$, $\hbar$ explicit, signature $(+,-,-,-)$, intrinsic parity $\epsilon=+1$ as per Proposition~\ref{prop:rest-selection}, and the label $j=\tfrac12$ in item (B) of Theorem~\ref{thm:main}.
\end{remark}

Hypotheses H1 to H3 are the only input on the physical system, while H4 to H6 select the way in which its states are described. Item (A) of the following theorem uses the first group alone, and item (B) uses both groups.

\begin{theorem}[Recovery of the Dirac equation]\label{thm:main}
	Let $\Phi$ be an elementary relativistic system, as per Definition~\ref{def:elementary-system}, \textit{i.e.}, satisfying H1 to H3, with mass $m>0$ as per Proposition~\ref{prop:mass-shell}.
	\begin{enumerate}
		\item[\textbf{(A)}] There exist a unique $j\in\{0,\tfrac12,1,\dots\}$ and a Kähler isomorphism $\Xi:\mathbb{P}\mathcal{H}\to\mathbb{P}L^2(\mathcal{O}_m^+,\mu;\mathbb{C}^{2j+1})$ such that $\Xi\circ\Phi_\kappa=[U_{m,j}(\kappa)\,\cdot\,]\circ\Xi$ for all $\kappa\in\mathcal{P}$. In this realization $H=\omega(\mathbf{P})$ with $D(H)=D(|\mathbf{P}|)$.
		\item[\textbf{(B)}] Assume $j=\tfrac12$. Up to equivalence of realizations, as per Proposition~\ref{prop:realization-equivalence}, and up to the sign $\epsilon$ of the intrinsic parity, there is exactly one realization of $(\mathbb{P}\mathcal{H},\Phi)$ satisfying H4 to H6, namely the Dirac realization. In it the states are the rays of the solutions of
		\begin{equation*}
			(\gamma_\mu p^\mu-m)\,\Psi(p)=0, \qquad p\in\mathcal{O}_m^+,
		\end{equation*}
		with the scalar product $\langle\cdot,\cdot\rangle_D$ of Theorem~\ref{thm:covariant-norm}, the Poincar\'e group acts through $[\mathcal{U}_D(\kappa)\,\cdot\,]$ with the fixed matrices $D=D^{(1/2,0)}\oplus D^{(0,1/2)}$ of Definition~\ref{def:dirac-space}, and the Hamiltonian flow of $\langle H\rangle$ on these rays is the time evolution of the Dirac equation.
	\end{enumerate}
\end{theorem}

The proof of Theorem~\ref{thm:main} is the chain of eight steps collected in Table~\ref{tab:steps}, in which each step replaces a model of the states by an equivalent one. The columns record the section in which the step is carried out, the result that closes it and the hypotheses of Definition~\ref{def:hypotheses} that the step uses.

\begin{table}[ht]
	\centering
	\small
	\renewcommand{\arraystretch}{1.25}
	\begin{tabular}{llp{7.6cm}l}
		\hline
		\textbf{Step} & \textbf{Section} & \textbf{Result} & \textbf{Hypotheses} \\
		\hline
		S0 & 6.2 & Unique unitary lift $\Pi$ of $\Phi$, Proposition~\ref{prop:poincare-lift}, and realizations as unitary equivalences of lifts, Proposition~\ref{prop:realization-equivalence} & H1 \\
		S1 & 6.3 & Spectral measure of the translations concentrated on $\mathcal{O}_m^+$, Proposition~\ref{prop:mass-shell} & H1 to H3 \\
		S2 & 6.4 & Realization on $L^2(\mathcal{O}_m^+,\mu;\mathbb{C}^{2j+1})$ and $H=\omega(\mathbf{P})$, Theorems~\ref{thm:wigner-mackey} and~\ref{thm:relativistic-hamiltonian} & H1 to H3 \\
		S3 & 6.5 & No local realization with unitary matrices, Proposition~\ref{prop:no-local-unitary} & H1 to H4 \\
		S4 & 6.5 & Parity and minimality force $W=\mathbb{C}^4$, Proposition~\ref{prop:minimal-extension} & H1 to H6 \\
		S5 & 6.6 & Rest subspace $F(\mathring{p})$, Proposition~\ref{prop:rest-selection} & H1 to H6 \\
		S6 & 6.6 & Transport along the mass shell and Dirac equation, Theorem~\ref{thm:dirac} & H1 to H6 \\
		S7 & 6.7 & Covariant norm, Kähler isomorphism and Dirac flow, Theorem~\ref{thm:covariant-norm}, Propositions~\ref{prop:kahler-dirac} and~\ref{prop:dirac-flow} & H1 to H6 \\
		\hline
	\end{tabular}
	\caption{The chain of the chapter. Steps S0 to S2 lead from the abstract space $\mathbb{P}\mathcal{H}$ to the model on the mass shell, and steps S3 to S7 lead from the latter to the model by solutions of the Dirac equation.}
	\label{tab:steps}
\end{table}

% ---------------------------------------------------------------------
\paragraph{Registry of the Hypotheses}
\label{sec:dirac-summary}

The chain of the chapter replaces the abstract description of the states by two successive models, and each replacement needs its own requirements. The question answered in this section concerns the states and reads as follows: which hypotheses of Definition~\ref{def:hypotheses} enter at which step, \textit{i.e.}, what has to be assumed on the pair $(\mathbb{P}\mathcal{H},\Phi)$ and what on the way in which the states are described by fields. The data are the results of Sections 6.2 to 6.7. The outcome is Table~\ref{tab:hypotheses-steps}, which assigns to every hypothesis the steps and the results in which it is used. It shows that the passage to the mass shell rests on the symmetry, irreducibility and positivity of the energy alone, while the passage to the Dirac equation rests in addition on locality, parity and minimality. This section closes the chain, and it serves to make the statement of Theorem~\ref{thm:main} verifiable item by item. A remark records what the construction leaves out.

\begin{table}[ht]
	\centering
	\small
	\renewcommand{\arraystretch}{1.25}
	\begin{tabular}{llp{6.4cm}l}
		\hline
		\textbf{Hypothesis} & \textbf{Steps} & \textbf{Results in which it is used} & \textbf{Theorem~\ref{thm:main}} \\
		\hline
		H1 & S0 & Corollary~\ref{cor:poincare-lift}, Proposition~\ref{prop:poincare-lift} & (A), (B) \\
		H2 & S1, S2 & Proposition~\ref{prop:mass-shell}, Theorem~\ref{thm:wigner-mackey} & (A), (B) \\
		H3 & S1, S2 & Proposition~\ref{prop:mass-shell}, Theorem~\ref{thm:relativistic-hamiltonian} & (A), (B) \\
		H4 & S3 & Lemma~\ref{lem:local-realization}, Proposition~\ref{prop:no-local-unitary} & (B) \\
		H5 & S4, S5 & Propositions~\ref{prop:parity-problem}, \ref{prop:minimal-extension} and~\ref{prop:rest-selection} & (B) \\
		H6 & S4 to S6 & Proposition~\ref{prop:minimal-extension}, Theorem~\ref{thm:dirac} & (B) \\
		\hline
	\end{tabular}
	\caption{Hypotheses of Definition~\ref{def:hypotheses}, the steps of Table~\ref{tab:steps} in which they enter and the item of Theorem~\ref{thm:main} that depends on them.}
	\label{tab:hypotheses-steps}
\end{table}

\begin{remark}[What is not treated]\label{rem:not-treated}
	Three situations are left out. The massless case $m=0$, allowed by Proposition~\ref{prop:mass-shell}, requires a different analysis, since the stabilizer of a null momentum is the covering of the Euclidean group of the plane and not $SU(2)$, so that the induced representations of Section 6.4 do not apply and helicity replaces spin. Interactions are not introduced, and the local transformation law of Section 6.5 is only the prerequisite for a coupling to an external potential. No position observable is constructed, so that the density $\Psi^\dagger\Psi$ of Proposition~\ref{prop:continuity} is the time component of a conserved current and not the probability density of a position measurement. Solutions supported on $\mathcal{O}_m^-$ are excluded by H3, and spins other than $\tfrac12$ are covered by item (A) alone.
\end{remark}

\section{Mathematical preliminaries}
\label{sec:preliminaries}

The chapter is organized around the pure states of a relativistic system. These are the rays of $\mathbb{P}\mathcal{H}$, on which a change of inertial observer acts by Kähler isomorphisms. This leaves three issues open: which group of observers acts, which phases its implementation on vectors of $\mathcal{H}$ may carry, and which momenta and which transformation laws of the components are available to a state of given mass and spin. The present section answers the algebraic part of each issue before any state space is introduced. What is given is only linear algebra on $2\times2$ matrices and the Lie theory of Chapter~1. What is obtained is a collection of results on the Poincar\'e group, on the mass hyperboloid and on the finite dimensional representations of $SL(2,\mathbb{C})$, none of which constructs a model. The section is the preliminary step of the chain and precedes the lift of the action to unitary operators. Collecting these results here prevents a later proof from stopping to verify an identity between matrices, and each result is followed by a line recording the statement that uses it. The first subsection fixes the group and proves that its Lie algebra has no central extension, which is used for the lift of Proposition~\ref{prop:poincare-lift} and for the covariance of Lemma~\ref{lem:poincare-covariance}. The second describes the mass shell as a homogeneous space together with the induced representations, and it is used for Proposition~\ref{prop:mass-shell} and Theorem~\ref{thm:wigner-mackey}. The third classifies the finite dimensional representations of $SL(2,\mathbb{C})$ and introduces space inversion, and it is used for Proposition~\ref{prop:no-local-unitary} and Proposition~\ref{prop:minimal-extension}.

%--------------------------------------------------------------------
\subsection{The Poincar\'e group and its Lie algebra}
\label{sec:poincare-group}

The states of an elementary system are acted upon by the group of isometries of Minkowski spacetime, and two properties of this group are needed before any state space is considered. The first is an explicit description of its covering group acting on Hermitian matrices, in which boosts, rotations and space inversion are computed by matrix products. The second is that its Lie algebra admits no nontrivial central extension, so that a continuous action on rays lifts to vectors without enlarging the group. The subsection starts from the matrix model and ends with the vanishing of the second cohomology group of the Lie algebra. It serves the lift of the action, and through it the spectral analysis of the translations.

We write $\mathbb{M}:=\mathbb{R}^{1,3}$ for Minkowski space, with points $x=(x^0,\mathbf{x})$, bilinear form $x\cdot y:=x^0y^0-\mathbf{x}\cdot\mathbf{y}=\eta(x,y)$ for $\eta=\operatorname{diag}(1,-1,-1,-1)$, and $x^2:=x\cdot x$. The speed of light is set equal to one. The Lorentz group $O(1,3)$ is the Lie group of linear maps of $\mathbb{M}$ preserving this form, and $SO^\uparrow(1,3)$ denotes its connected component of the identity, \textit{i.e.}, the group of linear maps preserving orientation and the direction of time \cite[TODO-location]{TODO-BarutRaczka}. It is a closed subgroup of $GL(4,\mathbb{R})$, because connected components are closed.

\begin{definition}[Hermitian matrix model of $\mathbb{M}$]\label{def:hermitian-model}
	For $x\in\mathbb{M}$ we set
	\begin{equation}\label{eq:hermitian-map}
		\hat{x}:=x^0\,\mathbb{1}+\mathbf{x}\cdot\boldsymbol{\sigma}, \qquad \mathbf{x}\cdot\boldsymbol{\sigma}=\sum_{j=1}^3x^j\sigma_j,
	\end{equation}
	where $\sigma_j$ are the Pauli matrices of Example~1.2.5.
\end{definition}

\begin{remark}
	We also use the matrix $\tilde{x}:=x^0\,\mathbb{1}-\mathbf{x}\cdot\boldsymbol{\sigma}$, which is the matrix of Definition~\ref{def:hermitian-model} for the vector with reversed spatial part, \textit{i.e.}, for the image of $x$ under the space inversion of Lemma~\ref{lem:inversion}.
\end{remark}

The determinant of $\hat{x}$ is the Minkowski square, and this is the property that makes the model useful, since matrices of determinant one preserve it. The adjugate and the sign of the eigenvalues are recorded together with it.

\begin{lemma}[Algebra of the matrix model]\label{lem:hermitian-model}
	\begin{enumerate}
		\item The map $x\mapsto\hat{x}$ is a real linear bijection from $\mathbb{M}$ onto the real vector space of Hermitian $2\times2$ matrices, and $\det\hat{x}=x^2$.
		\item $\tilde{x}=\operatorname{adj}\hat{x}$, $\hat{x}\tilde{x}=\tilde{x}\hat{x}=x^2\,\mathbb{1}$ and $x\cdot y=\tfrac12\operatorname{tr}(\hat{x}\tilde{y})$.
		\item If $x^2>0$, then $\hat{x}$ is positive definite for $x^0>0$ and negative definite for $x^0<0$.
	\end{enumerate}
\end{lemma}

\begin{proof}
	Item 1. A Hermitian $2\times2$ matrix has the form $\begin{pmatrix}x^0+x^3&x^1-ix^2\\x^1+ix^2&x^0-x^3\end{pmatrix}$ with real $x^\mu$, and its determinant is $(x^0)^2-|\mathbf{x}|^2$.
	
	Item 2. For a $2\times2$ matrix $X$ the adjugate is $\operatorname{adj}X=(\operatorname{tr}X)\mathbb{1}-X$, hence $\operatorname{adj}\hat{x}=\tilde{x}$, and $X\operatorname{adj}X=\operatorname{adj}X\,X=\det X\,\mathbb{1}$. Since $\operatorname{tr}\sigma_j=0$ and $(\mathbf{x}\cdot\boldsymbol{\sigma})(\mathbf{y}\cdot\boldsymbol{\sigma})=\mathbf{x}\cdot\mathbf{y}\,\mathbb{1}+i(\mathbf{x}\times\mathbf{y})\cdot\boldsymbol{\sigma}$, one has $\operatorname{tr}(\hat{x}\tilde{y})=2(x^0y^0-\mathbf{x}\cdot\mathbf{y})$.
	
	Item 3. The eigenvalues of $\hat{x}$ are $x^0\pm|\mathbf{x}|$. Their product is $x^2>0$, so they have the same sign, which is the sign of their sum $2x^0$.
\end{proof}

\noindent Item 3 is used in Lemma~\ref{lem:hyperboloid-orbits} and Proposition~\ref{prop:orbit}, and item 2 in Lemma~\ref{lem:gamma} and Lemma~\ref{lem:inversion}.

The covering map of the next lemma acts on $\hat{x}$ by $g\hat{x}g^\dagger$, and surjectivity onto $SO^\uparrow(1,3)$ requires the connectedness of $SL(2,\mathbb{C})$. This is obtained from the polar decomposition, which also fixes the topology of the group.

\begin{lemma}[Polar decomposition]\label{lem:polar}
	The map $\Psi:\mathbb{R}^3\times SU(2)\to SL(2,\mathbb{C})$ defined by
	\begin{equation}
		\Psi(\mathbf{w},u):=\exp\!\big(\tfrac12\,\mathbf{w}\cdot\boldsymbol{\sigma}\big)\,u,
	\end{equation}
	is a diffeomorphism. In particular $SL(2,\mathbb{C})$ is connected and simply connected.
\end{lemma}

\begin{proof}
	Let $g\in SL(2,\mathbb{C})$. The matrix $gg^\dagger$ is Hermitian and positive definite with determinant one. It has a unique positive definite Hermitian square root $K:=(gg^\dagger)^{1/2}$, which depends smoothly on $g$ and satisfies $\det K=1$, because the eigenvalues of $K$ are the positive square roots of those of $gg^\dagger$. The matrix $u:=K^{-1}g$ satisfies $uu^\dagger=K^{-1}gg^\dagger K^{-1}=\mathbb{1}$ and $\det u=1$, so that $u\in SU(2)$ and $g=Ku$. If $g=K'u'$ is another decomposition of this type, then $gg^\dagger=K'^2$, and the uniqueness of the positive square root yields $K'=K$ and $u'=u$.
	
	It remains to parametrize the matrices $K$. A positive definite Hermitian matrix is the exponential of a unique Hermitian matrix $h=\log K$, which depends smoothly on $K$, and $\operatorname{tr}h=\log\det K=0$ because $\det e^h=e^{\operatorname{tr}h}$. Traceless Hermitian matrices are exactly the matrices $\tfrac12\mathbf{w}\cdot\boldsymbol{\sigma}$ with a unique $\mathbf{w}\in\mathbb{R}^3$, which shows that $\Psi$ is a smooth bijection with smooth inverse. The group $SU(2)$ is diffeomorphic to $S^3$, since it consists of the matrices with columns $(\alpha,\beta)^T$ and $(-\bar\beta,\bar\alpha)^T$ subject to $|\alpha|^2+|\beta|^2=1$. Hence $SL(2,\mathbb{C})\simeq\mathbb{R}^3\times S^3$ is connected and simply connected.
\end{proof}

\noindent See \cite[TODO-location]{TODO-Hal03} and \cite[Chapter 3, \S6]{TODO-BarutRaczka} for the polar and Cartan decompositions in general. Lemma~\ref{lem:polar} is used in Lemma~\ref{lem:sl2c-covering} for connectedness, in Proposition~\ref{prop:orbit} for the decomposition $g=A_pu$ and in Corollary~\ref{cor:poincare-topology} for simple connectedness.

With connectedness available, the covering map can be built from the action $\hat{x}\mapsto g\hat{x}g^\dagger$. The lemma also yields its differential, which is needed for the Lie algebra, and the covariance of the matrices $\tilde{x}$, which is needed for the gamma matrices and for space inversion.

\begin{lemma}[Covering of the Lorentz group]\label{lem:sl2c-covering}
	For $g\in SL(2,\mathbb{C})$ let $\Lambda(g):\mathbb{M}\to\mathbb{M}$ be the linear map defined by
	\begin{equation}\label{eq:lorentz-action}
		\widehat{\Lambda(g)x}=g\,\hat{x}\,g^\dagger.
	\end{equation}
	\begin{enumerate}
		\item $g\mapsto\Lambda(g)$ is a smooth surjective homomorphism $SL(2,\mathbb{C})\to SO^\uparrow(1,3)$ with kernel $\{\mathbb{1},-\mathbb{1}\}$, and hence a covering map.
		\item Its differential at the identity is the injective linear map $\rho:\mathfrak{sl}(2,\mathbb{C})\to\mathfrak{gl}(\mathbb{M})$ defined by
		\begin{equation}\label{eq:rho}
			\widehat{\rho(\xi)x}=\xi\,\hat{x}+\hat{x}\,\xi^\dagger,
		\end{equation}
		where $\mathfrak{sl}(2,\mathbb{C})$ is regarded as a real Lie algebra of dimension $6$.
		\item For all $g\in SL(2,\mathbb{C})$ and $x\in\mathbb{M}$,
		\begin{equation}\label{eq:tilde-covariance}
			\widetilde{\Lambda(g)x}=(g^\dagger)^{-1}\,\tilde{x}\,g^{-1}.
		\end{equation}
	\end{enumerate}
\end{lemma}

\begin{proof}
	Item 1. The matrix $g\hat{x}g^\dagger$ is Hermitian, so Equation~\eqref{eq:lorentz-action} defines a linear map $\Lambda(g)$ by item 1 of Lemma~\ref{lem:hermitian-model}. It preserves $x^2=\det\hat{x}$ because $\det(g\hat{x}g^\dagger)=|\det g|^2\det\hat{x}=\det\hat{x}$, and it preserves the bilinear form by polarization, hence $\Lambda(g)\in O(1,3)$. The relation $\Lambda(gg')=\Lambda(g)\Lambda(g')$ is immediate, and $g\mapsto\Lambda(g)$ is polynomial in the real and imaginary parts of the entries of $g$. By Lemma~\ref{lem:polar} the group $SL(2,\mathbb{C})$ is connected, hence its image lies in the connected component $SO^\uparrow(1,3)$.
	
	Item 2. Differentiating $g(t)\hat{x}g(t)^\dagger$ along $g(t)=\exp(t\xi)$ at $t=0$ yields Equation~\eqref{eq:rho}. If $\rho(\xi)=0$, then $x=(1,\mathbf{0})$ yields $\xi+\xi^\dagger=0$, and $x=(0,\mathbf{e}_j)$ then yields $\xi\sigma_j-\sigma_j\xi=0$ for $j=1,2,3$. The matrix $\xi$ commutes with $\mathbb{1},\sigma_1,\sigma_2,\sigma_3$, which span $M_2(\mathbb{C})$, so $\xi$ is a multiple of the identity, and $\operatorname{tr}\xi=0$ forces $\xi=0$. Hence $\rho$ is injective. The Lie algebra of $O(1,3)$ is $\{M\in M_4(\mathbb{R})\mid\eta M+M^T\eta=0\}$, where $\eta M$ is antisymmetric, so it has dimension $6=\dim\mathfrak{sl}(2,\mathbb{C})$. The image of $\rho$ lies in it, since $\Lambda(\exp t\xi)=\exp(t\rho(\xi))$ is in $O(1,3)$. Hence $\rho$ is a linear bijection onto this Lie algebra, and by the inverse function theorem $g\mapsto\Lambda(g)$ maps a neighbourhood of $\mathbb{1}$ diffeomorphically onto a neighbourhood of the identity of $SO^\uparrow(1,3)$. The image is a subgroup containing this neighbourhood, and a neighbourhood of the identity in a connected group generates the whole group, so $g\mapsto\Lambda(g)$ is surjective.
	
	If $\Lambda(g)=\mathrm{id}$, then $gg^\dagger=\mathbb{1}$ from $x=(1,\mathbf{0})$ and $g\sigma_jg^{-1}=\sigma_j$ for all $j$, so that $g$ commutes with $M_2(\mathbb{C})$ and $g=\lambda\mathbb{1}$ with $\lambda^2=\det g=1$. The kernel is therefore $\{\mathbb{1},-\mathbb{1}\}$. A surjective Lie group homomorphism with discrete kernel is a covering map \cite[TODO-location]{TODO-Lee13}.
	
	Item 3. The adjugate satisfies $\operatorname{adj}(XY)=\operatorname{adj}(Y)\operatorname{adj}(X)$ and, for $\det g=1$, $\operatorname{adj}g=g^{-1}$. Together with $\operatorname{adj}\hat{x}=\tilde{x}$ from item 2 of Lemma~\ref{lem:hermitian-model}, it descends that
	\begin{equation*}
		\widetilde{\Lambda(g)x}=\operatorname{adj}(g\hat{x}g^\dagger)=\operatorname{adj}(g^\dagger)\,\tilde{x}\,\operatorname{adj}(g)=(g^\dagger)^{-1}\tilde{x}\,g^{-1}.\qedhere
	\end{equation*}
\end{proof}

\noindent Item 1 is used in Definition~\ref{def:poincare-group} and, through its kernel, in Proposition~\ref{prop:central-sign}, item 2 in Proposition~\ref{prop:poincare-algebra}, and item 3 in Lemma~\ref{lem:gamma} and Lemma~\ref{lem:inversion}.

The covering map assembles the isometries of $\mathbb{M}$ and their universal covering into a single group. The construction follows the pattern of Definition~\ref{def:galilei-group}, with rotations acting on translations, and it differs only in the action of the Lorentz part.

\begin{definition}[Poincar\'e group]\label{def:poincare-group}
	Let $\Lambda(g)$ be as per Lemma~\ref{lem:sl2c-covering}. The \emph{Poincar\'e group} $\mathcal{P}$ is the Lie group with underlying manifold $\mathbb{M}\times SL(2,\mathbb{C})$ and product
	\begin{equation}
		(a,g)(a',g'):=\big(a+\Lambda(g)a',\ gg'\big).
	\end{equation}
\end{definition}

\begin{remark}\label{rem:poincare-group}
	The product is the one of the semidirect product $\mathbb{M}\rtimes SL(2,\mathbb{C})$ for the action $g\mapsto\Lambda(g)$ by automorphisms of the additive group of $\mathbb{M}$ \cite[Chapter 3, \S4]{TODO-BarutRaczka}. Bracketing a triple product in the two possible ways yields in both cases $(a+\Lambda(g)a'+\Lambda(gg')a'',\,gg'g'')$, since $\Lambda(gg')=\Lambda(g)\Lambda(g')$, so the product is associative. The identity is $(0,\mathbb{1})$ and the inverse of $(a,g)$ is $(-\Lambda(g^{-1})a,\,g^{-1})$. These operations are smooth, because they are polynomial in $a$ and in the real and imaginary parts of the entries of $g$, and $\mathcal{P}$ is a Lie group of dimension $10$. The group acts on $\mathbb{M}$ by $(a,g)\cdot x:=\Lambda(g)x+a$, which is an action in the sense of Definition~1.5.1, because $\big((a,g)(a',g')\big)\cdot x=\Lambda(gg')x+a+\Lambda(g)a'=(a,g)\cdot\big((a',g')\cdot x\big)$. The kernel of the action is $\{(0,\mathbb{1}),(0,-\mathbb{1})\}$, so that $\bar{\mathcal{P}}_0:=\mathcal{P}/\{(0,\pm\mathbb{1})\}\simeq\mathbb{M}\rtimes SO^\uparrow(1,3)$ is the group of proper orthochronous isometries of $\mathbb{M}$ acting effectively, and $q:\mathcal{P}\to\bar{\mathcal{P}}_0$ is a double covering.
\end{remark}

\noindent Definition~\ref{def:poincare-group} is used in Definition~\ref{def:poincare-action}, where $\mathcal{P}$ acts on $\mathbb{P}\mathcal{H}$.

Every element of $\mathcal{P}$ is a product of three kinds of transformations, and the notation for them is fixed before their relations are recorded. The boosts are parametrized by the rapidity vector $\mathbf{w}$ and not by a velocity, because by item 4 of the next lemma the velocity of $\beta_{\mathbf{w}}$ is $\tanh|\mathbf{w}|$.

\begin{definition}[Elementary transformations]\label{def:poincare-elementary}
	Let $\mathcal{P}$ be as per Definition~\ref{def:poincare-group}. For $a\in\mathbb{M}$, $\mathbf{w}\in\mathbb{R}^3$ and $u\in SU(2)$ we set
	\begin{equation}
		\tau_a:=(a,\mathbb{1}),\qquad \beta_{\mathbf{w}}:=\big(0,\exp(\tfrac12\mathbf{w}\cdot\boldsymbol{\sigma})\big),\qquad \rho_u:=(0,u),
	\end{equation}
	and we call them the translation, the boost and the rotation. The time translation is $\tau_{(s,\mathbf{0})}$ and the spatial translation is $\tau_{(0,\mathbf{a})}$.
\end{definition}

The relations below are the finite form of the brackets computed afterwards. Item 2 states how the Lorentz part acts on translations, and item 5 states that the three families exhaust the group, which fixes its topology.

\begin{lemma}[Relations between the elementary transformations]\label{lem:poincare-relations}
	Let $\mathcal{P}$ be as per Definition~\ref{def:poincare-group}, with elementary transformations as per Definition~\ref{def:poincare-elementary}. For all $a,a'\in\mathbb{M}$, $\mathbf{w}\in\mathbb{R}^3$, $s\in\mathbb{R}$ and $u,u'\in SU(2)$ the following hold.
	\begin{enumerate}
		\item $\tau_a\tau_{a'}=\tau_{a+a'}$ and $\rho_u\rho_{u'}=\rho_{uu'}$.
		\item $(0,g)\,\tau_a\,(0,g)^{-1}=\tau_{\Lambda(g)a}$ for every $g\in SL(2,\mathbb{C})$.
		\item $\rho_u\tau_a\rho_u^{-1}=\tau_{(a^0,R_u\mathbf{a})}$ for $a=(a^0,\mathbf{a})$, and $\rho_u\beta_{\mathbf{w}}\rho_u^{-1}=\beta_{R_u\mathbf{w}}$, where $R_u\in SO(3)$ is the rotation determined by $u\,(\mathbf{x}\cdot\boldsymbol{\sigma})\,u^{-1}=(R_u\mathbf{x})\cdot\boldsymbol{\sigma}$, \textit{i.e.}, the image of $u$ under the covering homomorphism $SU(2)\to SO(3)$ of Example~\ref{ex:so3-su2}. In particular $\Lambda(u)(x^0,\mathbf{x})=(x^0,R_u\mathbf{x})$.
		\item $\beta_{\mathbf{w}}\tau_{(s,\mathbf{0})}\beta_{\mathbf{w}}^{-1}=\tau_{(s\cosh|\mathbf{w}|,\ s\sinh|\mathbf{w}|\,\mathbf{w}/|\mathbf{w}|)}$ for $\mathbf{w}\neq\mathbf{0}$.
		\item Every element of $\mathcal{P}$ is written in a unique way as $\tau_a\,\beta_{\mathbf{w}}\,\rho_u$, and the map $(a,\mathbf{w},u)\mapsto\tau_a\beta_{\mathbf{w}}\rho_u$ is a diffeomorphism $\mathbb{M}\times\mathbb{R}^3\times SU(2)\to\mathcal{P}$.
	\end{enumerate}
\end{lemma}

\begin{proof}
	Item 1 follows from Definition~\ref{def:poincare-group}, since $\Lambda(\mathbb{1})=\mathrm{id}$ and $(0,u)(0,u')=(0,uu')$.
	
	Item 2. By Definition~\ref{def:poincare-group}, $(0,g)(a,\mathbb{1})=(\Lambda(g)a,g)=(\Lambda(g)a,\mathbb{1})(0,g)$.
	
	Item 3. For $u\in SU(2)$ the matrix $u(\mathbf{x}\cdot\boldsymbol{\sigma})u^{-1}$ is traceless and Hermitian, hence of the form $\mathbf{y}\cdot\boldsymbol{\sigma}$, and $u\mathbb{1}u^{-1}=\mathbb{1}$. This yields $\Lambda(u)(x^0,\mathbf{x})=(x^0,\mathbf{y})$. Since $\det(\mathbf{x}\cdot\boldsymbol{\sigma})=-|\mathbf{x}|^2$ is preserved by conjugation, $|\mathbf{y}|=|\mathbf{x}|$, and $R_u:\mathbf{x}\mapsto\mathbf{y}$ is a linear isometry of $\mathbb{R}^3$. It has determinant $+1$ because $SU(2)$ is connected. The first relation is item 2 with $g=u$. For the second, $u\exp(\tfrac12\mathbf{w}\cdot\boldsymbol{\sigma})u^{-1}=\exp(\tfrac12u(\mathbf{w}\cdot\boldsymbol{\sigma})u^{-1})=\exp(\tfrac12(R_u\mathbf{w})\cdot\boldsymbol{\sigma})$.
	
	Item 4. By item 2 the left-hand side is $\tau_{\Lambda(g)a}$ with $g=\exp(\tfrac12\mathbf{w}\cdot\boldsymbol{\sigma})$ and $a=(s,\mathbf{0})$. The matrix $g$ is Hermitian, so $\Lambda(g)a$ is represented by
	\begin{equation*}
		g\,(s\mathbb{1})\,g=s\,e^{\mathbf{w}\cdot\boldsymbol{\sigma}}=s\Big(\cosh|\mathbf{w}|\,\mathbb{1}+\sinh|\mathbf{w}|\,\tfrac{\mathbf{w}}{|\mathbf{w}|}\cdot\boldsymbol{\sigma}\Big),
	\end{equation*}
	which is the stated vector.
	
	Item 5. By Definition~\ref{def:poincare-group}, $\tau_a\beta_{\mathbf{w}}\rho_u=(a,\exp(\tfrac12\mathbf{w}\cdot\boldsymbol{\sigma})\,u)=(a,\Psi(\mathbf{w},u))$ with $\Psi$ as per Lemma~\ref{lem:polar}, which is a diffeomorphism.
\end{proof}

\noindent Items 1 and 2 are used in Proposition~\ref{prop:snag-translations} and Lemma~\ref{lem:poincare-covariance}, item 3 in Proposition~\ref{prop:induced-action}, where its second relation reappears for the standard boosts, item 4 in Proposition~\ref{prop:orbit}, where it gives meaning to the rapidity vector $\mathbf{w}_p$, and item 5 in Corollary~\ref{cor:poincare-topology}.

To decide whether the group has central extensions, and to identify the generators of its unitary representations, the commutation relations of its Lie algebra are needed. They are obtained by realizing $\mathcal{P}$ as a matrix group, which reduces the bracket to a commutator of block matrices.

\begin{proposition}[Lie algebra of the Poincar\'e group]\label{prop:poincare-algebra}
	Let $\rho$ be as in Equation~\eqref{eq:rho}. The Lie algebra $\mathfrak{p}$ of $\mathcal{P}$ is the real vector space $\mathfrak{sl}(2,\mathbb{C})\oplus\mathbb{M}$ with the bracket
	\begin{equation}\label{eq:poincare-bracket}
		[(\xi,x),(\eta,y)]=\big([\xi,\eta],\ \rho(\xi)y-\rho(\eta)x\big).
	\end{equation}
	Let $r_j:=(-\tfrac i2\sigma_j,0)$, $b_j:=(\tfrac12\sigma_j,0)$ for $j=1,2,3$ and $t_\mu:=(0,\mathbf{e}_\mu)$ for $\mu=0,\dots,3$, where $\mathbf{e}_\mu$ is the canonical basis of $\mathbb{M}$. Then these ten elements form a basis of $\mathfrak{p}$, and up to antisymmetry the only nonvanishing brackets between them are
	\begin{align}
		[r_j,r_k]&=\varepsilon_{jkl}\,r_l, & [r_j,b_k]&=\varepsilon_{jkl}\,b_l, & [b_j,b_k]&=-\varepsilon_{jkl}\,r_l,\label{eq:lorentz-algebra}\\
		[r_j,t_k]&=\varepsilon_{jkl}\,t_l, & [b_j,t_0]&=t_j, & [b_j,t_k]&=\delta_{jk}\,t_0,\label{eq:translation-action}
	\end{align}
	with $\varepsilon_{jkl}$ the Levi-Civita symbol and a sum over the repeated index $l$.
\end{proposition}

\begin{proof}
	Consider the smooth map $F:\mathcal{P}\to GL(5,\mathbb{R})$ given in block form by
	\begin{equation*}
		F(a,g)=\begin{pmatrix}\Lambda(g)&a\\0&1\end{pmatrix},
	\end{equation*}
	which is a homomorphism because $F(a,g)F(a',g')=\begin{pmatrix}\Lambda(g)\Lambda(g')&\Lambda(g)a'+a\\0&1\end{pmatrix}=F\big((a,g)(a',g')\big)$. Along the curve $t\mapsto(tx,\exp(t\xi))$ one has $\Lambda(\exp(t\xi))=\exp(t\rho(\xi))$, by the naturality property in Proposition~1.2.14 applied to the homomorphism of Lemma~\ref{lem:sl2c-covering}. The differential of $F$ at the identity is therefore
	\begin{equation*}
		dF(\xi,x)=\begin{pmatrix}\rho(\xi)&x\\0&0\end{pmatrix}.
	\end{equation*}
	This map is injective by Lemma~\ref{lem:sl2c-covering}, and it is a Lie algebra homomorphism into $\mathfrak{gl}(5,\mathbb{R})=M_5(\mathbb{R})$ with the commutator of Example~1.2.4, since the differential of a Lie group homomorphism preserves the brackets \cite[TODO-location]{TODO-Lee13}. A direct computation of the commutator of two such matrices yields
	\begin{equation*}
		\left[\begin{pmatrix}\rho(\xi)&x\\0&0\end{pmatrix},\begin{pmatrix}\rho(\eta)&y\\0&0\end{pmatrix}\right]=\begin{pmatrix}\rho([\xi,\eta])&\rho(\xi)y-\rho(\eta)x\\0&0\end{pmatrix},
	\end{equation*}
	where $[\rho(\xi),\rho(\eta)]=\rho([\xi,\eta])$ because $\rho$ is the differential of a homomorphism. Since $dF$ is an injective homomorphism, Equation~\eqref{eq:poincare-bracket} follows.
	
	The elements $r_j$ and $b_j$ form a basis of $\mathfrak{sl}(2,\mathbb{C})$ as a real vector space, because $\tfrac12\sigma_j$ and $\tfrac i2\sigma_j$ span the traceless complex matrices. For the brackets in $\mathfrak{sl}(2,\mathbb{C})$, $[\sigma_j,\sigma_k]=2i\varepsilon_{jkl}\sigma_l$ yields $[r_j,r_k]=-\tfrac14[\sigma_j,\sigma_k]=-\tfrac i2\varepsilon_{jkl}\sigma_l=\varepsilon_{jkl}r_l$, then $[r_j,b_k]=-\tfrac i4[\sigma_j,\sigma_k]=\tfrac12\varepsilon_{jkl}\sigma_l=\varepsilon_{jkl}b_l$ and $[b_j,b_k]=\tfrac14[\sigma_j,\sigma_k]=\tfrac i2\varepsilon_{jkl}\sigma_l=-\varepsilon_{jkl}r_l$. For the action on translations we evaluate $\rho$ on $x=(x^0,\mathbf{x})$. Since $r_j^\dagger=-r_j$, $\widehat{\rho(r_j)x}=[r_j,\hat{x}]=\varepsilon_{jkl}x^k\sigma_l$, so that $\rho(r_j)x=(0,\mathbf{e}_j\times\mathbf{x})$. Since $b_j^\dagger=b_j$, $\widehat{\rho(b_j)x}=\tfrac12\{\sigma_j,\hat{x}\}=x^0\sigma_j+x^j\mathbb{1}$, so that $\rho(b_j)x=(x^j,\,x^0\mathbf{e}_j)$. Equation~\eqref{eq:poincare-bracket} yields $[(\xi,0),(0,x)]=(0,\rho(\xi)x)$ and $[(0,x),(0,y)]=0$, and evaluating on $x=\mathbf{e}_0$ and $x=\mathbf{e}_k$ yields Equation~\eqref{eq:translation-action}.
\end{proof}

\noindent The brackets are used in Lemma~\ref{lem:chiral} and Lemma~\ref{lem:rotation-invariant}, and the map $\rho$ in Theorem~\ref{thm:poincare-cohomology}.

\begin{remark}
	The bracket $[b_j,b_k]=-\varepsilon_{jkl}r_l$ states that two boosts in different directions do not commute and that their commutator is a rotation. It is the infinitesimal form of the momentum dependent rotation of Proposition~\ref{prop:wigner-rotation}. The bracket $[b_j,t_k]=\delta_{jk}t_0$ is not central, and this is the fact that makes the Poincar\'e algebra perfect in Corollary~\ref{cor:poincare-structure}.
\end{remark}

Over the reals the Lorentz algebra shows no structure beyond Equation~\eqref{eq:lorentz-algebra}. After complexification, rotations and boosts combine into two commuting copies of $\mathfrak{su}(2)$. The decomposition is needed twice, to prove semisimplicity, and to label the finite dimensional representations of $SL(2,\mathbb{C})$ by a pair of spins. We write $\mathfrak{sl}(2,\mathbb{C})_{\mathbb{C}}:=\mathfrak{sl}(2,\mathbb{C})\otimes_{\mathbb{R}}\mathbb{C}$, and $r_j$ and $b_j$ also denote the corresponding elements $(\xi,0)$ of $\mathfrak{sl}(2,\mathbb{C})$.

\begin{lemma}[Chiral decomposition]\label{lem:chiral}
	In $\mathfrak{sl}(2,\mathbb{C})_{\mathbb{C}}$ let $A_k:=\frac12(r_k-ib_k)$ and $B_k:=\frac12(r_k+ib_k)$. Then
	\begin{equation}\label{eq:chiral-brackets}
		[A_j,A_k]=\sum_l\varepsilon_{jkl}A_l, \qquad [B_j,B_k]=\sum_l\varepsilon_{jkl}B_l, \qquad [A_j,B_k]=0,
	\end{equation}
	and $r_k=A_k+B_k$, $b_k=i(A_k-B_k)$. Hence $\mathfrak{sl}(2,\mathbb{C})_{\mathbb{C}}\cong\mathfrak{su}(2)_{\mathbb{C}}\oplus\mathfrak{su}(2)_{\mathbb{C}}$, the sum of the two commuting ideals spanned by the $A_k$ and by the $B_k$.
\end{lemma}

\begin{proof}
	By Equation~\eqref{eq:lorentz-algebra}, and using $[b_j,r_k]=\sum_l\varepsilon_{jkl}b_l$,
	\begin{align*}
		[A_j,A_k]&=\tfrac14\big([r_j,r_k]-i[r_j,b_k]-i[b_j,r_k]-[b_j,b_k]\big)=\tfrac14\sum_l\varepsilon_{jkl}\big(2r_l-2ib_l\big)=\sum_l\varepsilon_{jkl}A_l,\\
		[A_j,B_k]&=\tfrac14\big([r_j,r_k]+i[r_j,b_k]-i[b_j,r_k]+[b_j,b_k]\big)=\tfrac14\sum_l\varepsilon_{jkl}\big(r_l+ib_l-ib_l-r_l\big)=0,
	\end{align*}
	and the computation for $B$ is the complex conjugate of the first one. The expressions of $r_k$ and $b_k$ are solved from the definitions. Since $r_k,b_k$ form a basis of $\mathfrak{sl}(2,\mathbb{C})_{\mathbb{C}}$, so do $A_k,B_k$, which yields the decomposition into the two ideals.
\end{proof}

\noindent Lemma~\ref{lem:chiral} is used in Corollary~\ref{cor:poincare-structure}, Theorem~\ref{thm:finite-reps} and Proposition~\ref{prop:no-unitary}.

Two structural properties of $\mathfrak{p}$ follow from the brackets and from the chiral decomposition. They are the hypotheses on which the vanishing of the second cohomology and the uniqueness of the lift rest.

\begin{corollary}[Semisimplicity and perfectness]\label{cor:poincare-structure}
	\begin{enumerate}
		\item The real Lie algebra $\mathfrak{l}:=\mathfrak{sl}(2,\mathbb{C})$ is semisimple.
		\item $[\mathfrak{p},\mathfrak{p}]=\mathfrak{p}$, \textit{i.e.}, the Poincar\'e algebra is perfect.
	\end{enumerate}
\end{corollary}

\begin{proof}
	Item 1. Each ideal of Lemma~\ref{lem:chiral} is isomorphic to $\mathfrak{su}(2)_{\mathbb{C}}=\mathfrak{sl}(2,\mathbb{C})$ as a complex Lie algebra, which is simple. Hence the complexification of $\mathfrak{l}$ is semisimple, and a real Lie algebra is semisimple if and only if its complexification is \cite[TODO-location]{TODO-Hum10}.
	
	Item 2. By Equations~\eqref{eq:lorentz-algebra} and~\eqref{eq:translation-action}, $r_l=[r_j,r_k]$ and $b_l=[r_j,b_k]$ for $j,k,l$ a cyclic permutation of $1,2,3$, $t_j=[b_j,t_0]$ and $t_0=[b_1,t_1]$.
\end{proof}

\noindent Item 1 is used in Theorem~\ref{thm:poincare-cohomology}, and item 2 in Proposition~\ref{prop:poincare-lift}, where it forces the lift of the action to be unique.

The topology of $\mathcal{P}$ is the last ingredient needed to apply Bargmann's correspondence. It also identifies the double covering of Remark~\ref{rem:poincare-group} as the universal one, so that the sign ambiguity of the rays has a single topological origin.

\begin{corollary}[Topology of the Poincar\'e group]\label{cor:poincare-topology}
	The group $\mathcal{P}$ is diffeomorphic to $\mathbb{R}^7\times S^3$, hence it is a connected and simply connected Lie group of dimension $10$. The homomorphism $q:\mathcal{P}\to\bar{\mathcal{P}}_0$ of Remark~\ref{rem:poincare-group} is the universal covering of the proper orthochronous Poincar\'e group, and $\pi_1(SO^\uparrow(1,3))\cong\mathbb{Z}_2$.
\end{corollary}

\begin{proof}
	By item 5 of Lemma~\ref{lem:poincare-relations}, $\mathcal{P}\simeq\mathbb{M}\times\mathbb{R}^3\times SU(2)$, and $SU(2)\simeq S^3$ is simply connected. The map $q$ is a covering map with kernel $\{(0,\pm\mathbb{1})\}$ by Lemma~\ref{lem:sl2c-covering} and Remark~\ref{rem:poincare-group}, and its total space is simply connected, so it is the universal covering. In a universal covering the group of deck transformations is isomorphic to the fundamental group of the base, and the deck group of $SL(2,\mathbb{C})\to SO^\uparrow(1,3)$ is $\{\pm\mathbb{1}\}\cong\mathbb{Z}_2$.
\end{proof}

\noindent Corollary~\ref{cor:poincare-topology} is used in Corollary~\ref{cor:poincare-lift}, where simple connectedness allows Bargmann's correspondence of Proposition~\ref{prop:bargmann-correspondence} to be applied.

The remainder of the subsection proves that $H^2(\mathfrak{p},\mathbb{R})=0$. A $2$-cocycle of $\mathfrak{p}$ records the phase between the commutator of two generators and the generator of their bracket, and it is harmless when it is a coboundary, as per Proposition~\ref{prop:lie-degree-two}. The Poincar\'e algebra is the semidirect sum of the semisimple Lorentz algebra and the abelian algebra of translations, so a cocycle can be removed in three layers. The Lorentz and mixed layers are removed by Whitehead's lemmas, and the translation layer by invariance under rotations.

\begin{theorem}[Whitehead's lemmas]\label{thm:whitehead}
	Let $\mathfrak{g}$ be a finite dimensional semisimple real Lie algebra and let $\varrho:\mathfrak{g}\to\mathfrak{gl}(N)$ be a representation on a finite dimensional real vector space $N$.
	\begin{enumerate}
		\item If a linear map $\phi:\mathfrak{g}\to N$ satisfies $\phi([\xi,\eta])=\varrho(\xi)\phi(\eta)-\varrho(\eta)\phi(\xi)$ for all $\xi,\eta\in\mathfrak{g}$, then there is $\mathsf{c}\in N$ such that $\phi(\xi)=\varrho(\xi)\mathsf{c}$ for all $\xi\in\mathfrak{g}$.
		\item $H^2(\mathfrak{g},\mathbb{R})=0$, \textit{i.e.}, every $2$-cocycle $\psi$ of $\mathfrak{g}$ is of the form $\psi_\lambda$ for some $\lambda\in\mathfrak{g}^*$.
	\end{enumerate}
\end{theorem}

\noindent See \cite[TODO-location]{TODO-Hum10}, where the statements are proved over $\mathbb{C}$. They hold for real Lie algebras since the cohomologies involved are compatible with complexification. Item 2 was already used for $\mathfrak{su}(2)$ in Example~\ref{ex:su2-cohomology}. Theorem~\ref{thm:whitehead} is used in Steps 1 and 2 of the proof of Theorem~\ref{thm:poincare-cohomology}.

Item 1 is applied to the representation of $\mathfrak{l}$ on the dual of the translation algebra, which is identified next. Throughout, $\mathfrak{t}:=\mathbb{M}$, so that $\mathfrak{p}=\mathfrak{l}\oplus\mathfrak{t}$ as in Proposition~\ref{prop:poincare-algebra}.

\begin{remark}\label{rem:lorentz-semisimple}
	We identify $\xi\in\mathfrak{l}$ with $(\xi,0)$ and $x\in\mathfrak{t}$ with $(0,x)$, so that Equation~\eqref{eq:poincare-bracket} reads
	\begin{equation}\label{eq:semidirect-brackets}
		[\xi,\eta]=[\xi,\eta]_{\mathfrak{l}},\qquad [\xi,x]=\rho(\xi)x,\qquad [x,y]=0,
	\end{equation}
	for $\xi,\eta\in\mathfrak{l}$ and $x,y\in\mathfrak{t}$. The dual space $\mathfrak{t}^*$ carries the representation $\varrho(\xi)f:=-f\circ\rho(\xi)$. One checks that $[\varrho(\xi),\varrho(\eta)]f=-f\circ\rho([\xi,\eta])=\varrho([\xi,\eta])f$, so that $\varrho$ is a representation, and both $\rho$ and $\varrho$ are finite dimensional. Hence Theorem~\ref{thm:whitehead} applies to $\mathfrak{g}=\mathfrak{l}$ with $N=\mathfrak{t}^*$, by item 1 of Corollary~\ref{cor:poincare-structure}.
\end{remark}

After the first two layers are removed only an antisymmetric form on the translations survives, and it is invariant under the Lorentz algebra. Rotations alone suffice to annihilate it, and no semisimplicity is needed for this step. Intuitively, a form on $\mathbb{R}\oplus\mathbb{R}^3$ that is invariant under rotations would have to be built from an invariant vector or from the cross product, and $\mathbb{R}^3$ has no nonzero invariant vector.

\begin{lemma}[Rotation invariant forms]\label{lem:rotation-invariant}
	Let $\Xi:\mathfrak{t}\times\mathfrak{t}\to\mathbb{R}$ be a bilinear antisymmetric map such that $\Xi(\rho(r_j)x,y)+\Xi(x,\rho(r_j)y)=0$ for all $x,y\in\mathfrak{t}$ and $j=1,2,3$. Then $\Xi=0$.
\end{lemma}

\begin{proof}
	Every bilinear antisymmetric map on $\mathbb{M}$ is of the form
	\begin{equation*}
		\Xi(x,y)=x^0\,\mathbf{v}\cdot\mathbf{y}-y^0\,\mathbf{v}\cdot\mathbf{x}+\mathbf{k}\cdot(\mathbf{x}\times\mathbf{y}),
	\end{equation*}
	for unique $\mathbf{v},\mathbf{k}\in\mathbb{R}^3$, namely $v^j=\Xi(\mathbf{e}_0,\mathbf{e}_j)$ and $k^l=\tfrac12\varepsilon_{jkl}\,\Xi(\mathbf{e}_j,\mathbf{e}_k)$. By the proof of Proposition~\ref{prop:poincare-algebra}, $\rho(r_j)x=(0,\mathbf{e}_j\times\mathbf{x})$. For $\mathbf{n}=\sum_jn^j\mathbf{e}_j$ the linear combination of the invariance conditions reads $\Xi(\mathbf{n}\times\mathbf{x},y)+\Xi(x,\mathbf{n}\times\mathbf{y})=0$, where $\mathbf{n}\times\mathbf{x}$ stands for $(0,\mathbf{n}\times\mathbf{x})$. Taking $\mathbf{x}=\mathbf{0}$ yields $x^0\,\mathbf{v}\cdot(\mathbf{n}\times\mathbf{y})=0$ for all $x^0$, $\mathbf{y}$ and $\mathbf{n}$, so that $\mathbf{n}\cdot(\mathbf{y}\times\mathbf{v})=0$ for all $\mathbf{n}$ and $\mathbf{y}$. Hence $\mathbf{y}\times\mathbf{v}=\mathbf{0}$ for all $\mathbf{y}$, which entails $\mathbf{v}=\mathbf{0}$. With $\mathbf{v}=\mathbf{0}$ the condition becomes
	\begin{equation*}
		\mathbf{k}\cdot\big((\mathbf{n}\times\mathbf{x})\times\mathbf{y}+\mathbf{x}\times(\mathbf{n}\times\mathbf{y})\big)=\mathbf{k}\cdot\big(\mathbf{n}\times(\mathbf{x}\times\mathbf{y})\big)=\mathbf{n}\cdot\big((\mathbf{x}\times\mathbf{y})\times\mathbf{k}\big)=0,
	\end{equation*}
	where the first equality is the Jacobi identity of the cross product. Since $\mathbf{n}$ is arbitrary, $(\mathbf{x}\times\mathbf{y})\times\mathbf{k}=\mathbf{0}$ for all $\mathbf{x},\mathbf{y}$. If $\mathbf{k}\neq\mathbf{0}$, choosing $\mathbf{x}\times\mathbf{y}$ nonvanishing and orthogonal to $\mathbf{k}$ yields a contradiction. Hence $\mathbf{k}=\mathbf{0}$.
\end{proof}

\noindent Lemma~\ref{lem:rotation-invariant} is used in Step 3 of the proof of Theorem~\ref{thm:poincare-cohomology}.

All the ingredients are now available. The proof decomposes a cocycle along $\mathfrak{p}=\mathfrak{l}\oplus\mathfrak{t}$ and removes its Lorentz, mixed and translation components one at a time by subtracting coboundaries. Each subtraction is allowed because $\lambda\mapsto\psi_\lambda$ is linear and every $\psi_\lambda$ is a cocycle, so that the difference between a cocycle and a coboundary is again a cocycle.

\begin{theorem}[$H^2$ of the Poincar\'e algebra]\label{thm:poincare-cohomology}
	Let $\mathfrak{p}$ be the Lie algebra of $\mathcal{P}$. Then $H^2(\mathfrak{p},\mathbb{R})=0$, \textit{i.e.}, every $2$-cocycle $\psi$ of $\mathfrak{p}$ is of the form $\psi_\lambda$ for some $\lambda\in\mathfrak{p}^*$.
\end{theorem}

\begin{proof}
	Let $\psi$ be a $2$-cocycle of $\mathfrak{p}$.
	
	\emph{Step 1, the Lorentz sector.} The restriction of $\psi$ to $\mathfrak{l}\times\mathfrak{l}$ is a $2$-cocycle of $\mathfrak{l}$, because $\mathfrak{l}$ is a subalgebra. By item 1 of Corollary~\ref{cor:poincare-structure} and item 2 of Theorem~\ref{thm:whitehead} there is $\lambda_0\in\mathfrak{l}^*$ with $\psi(\xi,\eta)=\lambda_0([\xi,\eta])$ for $\xi,\eta\in\mathfrak{l}$. We extend $\lambda_0$ to $\mathfrak{p}$ by $\lambda_0|_{\mathfrak{t}}=0$. Equation~\eqref{eq:semidirect-brackets} yields $\psi_{\lambda_0}(\xi,x)=\lambda_0(\rho(\xi)x)=0$ and $\psi_{\lambda_0}(x,y)=0$, so that $\psi':=\psi-\psi_{\lambda_0}$ is a cocycle vanishing on $\mathfrak{l}\times\mathfrak{l}$ and satisfying $\psi'(\xi,x)=\psi(\xi,x)$ and $\psi'(x,y)=\psi(x,y)$.
	
	\emph{Step 2, the mixed sector.} Set $\phi(\xi)(x):=\psi'(\xi,x)$, which defines a linear map $\phi:\mathfrak{l}\to\mathfrak{t}^*$. The cocycle condition for $\psi'$ on $(\xi,\eta,x)$ with $\xi,\eta\in\mathfrak{l}$ and $x\in\mathfrak{t}$ reads
	\begin{equation*}
		\psi'([\xi,\eta],x)+\psi'([\eta,x],\xi)+\psi'([x,\xi],\eta)=0.
	\end{equation*}
	By Equation~\eqref{eq:semidirect-brackets} and the antisymmetry of $\psi'$, the first term is $\phi([\xi,\eta])(x)$, the second is $-\phi(\xi)(\rho(\eta)x)$ and the third is $\phi(\eta)(\rho(\xi)x)$. Therefore
	\begin{equation*}
		\phi([\xi,\eta])(x)=\phi(\xi)(\rho(\eta)x)-\phi(\eta)(\rho(\xi)x).
	\end{equation*}
	In terms of the representation $\varrho$ of Remark~\ref{rem:lorentz-semisimple}, $\phi(\xi)\circ\rho(\eta)=-\varrho(\eta)\phi(\xi)$, and this identity becomes $\phi([\xi,\eta])=\varrho(\xi)\phi(\eta)-\varrho(\eta)\phi(\xi)$. Item 1 of Theorem~\ref{thm:whitehead} yields $\mathsf{c}\in\mathfrak{t}^*$ with $\phi(\xi)=\varrho(\xi)\mathsf{c}=-\mathsf{c}\circ\rho(\xi)$, \textit{i.e.}, $\psi'(\xi,x)=-\mathsf{c}(\rho(\xi)x)$. Let $\lambda_1\in\mathfrak{p}^*$ be the extension of $\mathsf{c}$ with $\lambda_1|_{\mathfrak{l}}=0$. Then $\psi_{\lambda_1}(\xi,x)=\lambda_1(\rho(\xi)x)=\mathsf{c}(\rho(\xi)x)$, while $\psi_{\lambda_1}(\xi,\eta)=\lambda_1([\xi,\eta]_{\mathfrak{l}})=0$ and $\psi_{\lambda_1}(x,y)=0$. It descends that $\psi'':=\psi'+\psi_{\lambda_1}$ is a cocycle vanishing on $\mathfrak{l}\times\mathfrak{l}$ and on $\mathfrak{l}\times\mathfrak{t}$.
	
	\emph{Step 3, the translation sector.} The cocycle condition for $\psi''$ on $(\xi,x,y)$ with $\xi\in\mathfrak{l}$ and $x,y\in\mathfrak{t}$ reads $\psi''([\xi,x],y)+\psi''([x,y],\xi)+\psi''([y,\xi],x)=0$. The second term vanishes because $[x,y]=0$, and the third equals $\psi''(x,\rho(\xi)y)$ by antisymmetry, so that
	\begin{equation*}
		\psi''(\rho(\xi)x,y)+\psi''(x,\rho(\xi)y)=0.
	\end{equation*}
	In particular this holds for $\xi=r_j$, and Lemma~\ref{lem:rotation-invariant} yields $\psi''|_{\mathfrak{t}\times\mathfrak{t}}=0$. Since $\psi''$ vanishes on all three sectors, $\psi''=0$ by bilinearity and antisymmetry. Collecting the three steps, $\psi=\psi_{\lambda_0}-\psi_{\lambda_1}=\psi_{\lambda_0-\lambda_1}$.
\end{proof}

\noindent Theorem~\ref{thm:poincare-cohomology} is used in Corollary~\ref{cor:poincare-lift}, where it removes every multiplier of a projective unitary representation of $\mathcal{P}$.

The hypothesis of semisimplicity is used in Steps 1 and 2 and cannot be dropped, as the lowest dimensional analogue shows.

\begin{remark}[The case of one space dimension]\label{rem:one-plus-one}
	In one space dimension the Poincar\'e algebra has basis $k,t_0,t_1$ with $[k,t_0]=t_1$, $[k,t_1]=t_0$ and $[t_0,t_1]=0$, and the Lorentz algebra spanned by $k$ is abelian, hence not semisimple. The antisymmetric form $\psi$ with $\psi(t_0,t_1)=1$ and $\psi(k,\cdot)=0$ satisfies the cocycle condition of Proposition~\ref{prop:lie-degree-two}. Its left-hand side is alternating in its three arguments, so that it suffices to evaluate it on $(k,t_0,t_1)$, where it reads $\psi(t_1,t_1)-\psi(t_0,t_0)=0$. It is not a coboundary, because $\psi_\lambda(t_0,t_1)=\lambda([t_0,t_1])=0$ for every $\lambda$. Hence the Poincar\'e algebra in $1+1$ dimensions has a nontrivial central extension, and Steps 1 to 3 do not apply to it.
\end{remark}

\subsection{The Mass Shell and Induced Representations}
\label{subsec:mass-shell-induced}

The pure states of an elementary relativistic system of mass $m$ remain the points of $\mathbb{P}\mathcal{H}$. The wave functions of Definition~\ref{def:induced-rep} depend on the four-momentum, and the set of its possible values is the mass shell $\mathcal{O}_m^+$. This set is therefore not a set of states but the arena on which a state is written. The question of this subsection is how a Lorentz transformation acts on a state with spin once the state is written as a function on that arena. The data are the covering $\Lambda$ of Lemma~\ref{lem:sl2c-covering} and the polar decomposition of Lemma~\ref{lem:polar}. The outcome is a geometric picture. The mass shell is three-dimensional hyperbolic space, and above each of its points stands a copy of $SU(2)$. A state with spin is a choice of a vector above each point. A Lorentz transformation moves the point together with the frame used to read the vector, and the transported frame differs from the frame at the new point by the Wigner rotation. This step makes the spin label of Theorem~\ref{thm:wigner-mackey} operational. It supplies the transformation law behind Definition~\ref{def:induced-rep} and the obstruction used in Proposition~\ref{prop:no-local-unitary}.

\subsubsection*{The Mass Shell and the Standard Boosts}

The four-momentum of a particle takes values in the dual $M^*$ of the translation group $M=\mathbb{R}^{1,3}$, which is identified with $M$ through $\eta$. The covering $\Lambda$ of Lemma~\ref{lem:sl2c-covering} acts on $M^*$ through this identification and preserves $p\cdot p$. The momenta available to a massive particle of positive energy therefore form an invariant subset of $M^*$.

\begin{definition}[Mass shell]\label{def:mass-shell}
	For $m>0$ the \emph{forward mass shell} is
	\begin{equation*}
		\mathcal{O}_m^+:=\{p\in M^*\mid p\cdot p=m^2,\ p^0>0\},
	\end{equation*}
	with \emph{rest momentum} $\mathring{p}:=(m,\mathbf{0})\in\mathcal{O}_m^+$, and $\mathcal{O}_m^-:=-\mathcal{O}_m^+$ is the \emph{backward mass shell}. For $\mathbf{p}\in\mathbb{R}^3$ we set $\omega(\mathbf{p}):=\sqrt{|\mathbf{p}|^2+m^2}$.
\end{definition}

\begin{remark}
	The map $\mathbf{p}\mapsto(\omega(\mathbf{p}),\mathbf{p})$ is a diffeomorphism $\mathbb{R}^3\to\mathcal{O}_m^+$, and $\mathbf{p}$ is the coordinate used throughout. Differentiating $p\cdot p=m^2$ shows that every tangent vector $v$ at $p\in\mathcal{O}_m^+$ satisfies $p\cdot v=0$. Since $p$ is timelike, $v$ is spacelike, hence $-\eta$ restricted to the tangent space of $\mathcal{O}_m^+$ is positive definite. This Riemannian metric is invariant under $\Lambda(g)$, and with it $\mathcal{O}_m^+$ is the hyperboloid model of three-dimensional hyperbolic space of curvature $-1/m^2$ \cite[TODO-location]{TODO-Ratcliffe}. The Lorentz group acts on it by isometries, and the rest momentum is a base point with no distinguished geometric role.
\end{remark}

That the mass shell is a single orbit is what allows every momentum to be described starting from the rest momentum alone. The next lemma proves it. It also treats the hyperboloids of negative squared mass, which the proof of Proposition~\ref{prop:mass-shell} has to exclude by positivity of the energy.

\begin{lemma}[Orbits of the Lorentz group on the hyperboloids]\label{lem:hyperboloid-orbits}
	Let $m>0$ and $\nu>0$.
	\begin{enumerate}
		\item $\Lambda(g)\mathcal{O}_m^{\pm}=\mathcal{O}_m^{\pm}$ for every $g\in SL(2,\mathbb{C})$, and $\Lambda(SL(2,\mathbb{C}))$ acts transitively on each of $\mathcal{O}_m^+$ and $\mathcal{O}_m^-$.
		\item The hyperboloid $\{p\in M^*\mid p\cdot p=-\nu^2\}$ is a single orbit of $\Lambda(SL(2,\mathbb{C}))$.
	\end{enumerate}
\end{lemma}

\begin{proof}
	Here $p$ is regarded as an element of $M$ through $\eta$, so that $\hat{p}$ and $\tilde{p}$ are as per Definition~\ref{def:hermitian-model}.
	
	Item 1. Let $p\in\mathcal{O}_m^+$. By item 3 of Lemma~\ref{lem:hermitian-model}, $\hat{p}$ is positive definite, hence so is $g\hat{p}g^\dagger$, and its determinant is $m^2$. It descends that $\Lambda(g)p$ has square $m^2$ and time component $\tfrac12\operatorname{tr}(g\hat{p}g^\dagger)>0$, \textit{i.e.}, it lies in $\mathcal{O}_m^+$. For transitivity, $\hat{p}/m$ is positive definite of determinant one, and its positive definite square root $A$ belongs to $SL(2,\mathbb{C})$. Then $\widehat{\Lambda(A)\mathring{p}}=A\,(m\mathbb{1})\,A=\hat{p}$, so that $p=\Lambda(A)\mathring{p}$. The statement on $\mathcal{O}_m^-=-\mathcal{O}_m^+$ follows from the linearity of $\Lambda(g)$.
	
	Item 2. Let $p\cdot p=-\nu^2$. The matrix $\hat{p}$ is Hermitian of determinant $-\nu^2<0$, so its eigenvalues satisfy $\lambda_+>0>\lambda_-$ and $\lambda_+\lambda_-=-\nu^2$. By the spectral theorem $\hat{p}=V\operatorname{diag}(\lambda_+,\lambda_-)V^\dagger$ with $V$ unitary, and multiplying $V$ by a phase we may take $\det V=1$. The matrix $D:=\operatorname{diag}\big(\sqrt{\lambda_+/\nu},\sqrt{\nu/\lambda_+}\big)$ belongs to $SL(2,\mathbb{C})$, and $D\operatorname{diag}(\nu,-\nu)D=\operatorname{diag}(\lambda_+,\lambda_-)$ because $\lambda_-=-\nu^2/\lambda_+$. Hence $\hat{p}=g\,(\nu\sigma_3)\,g^\dagger$ with $g=VD$, and $\nu\sigma_3$ is the matrix of $q=(0,0,0,\nu)$. It descends that $p=\Lambda(g)q$.
\end{proof}

The lemma provides, for every $p\in\mathcal{O}_m^+$, some Lorentz transformation taking the rest momentum to $p$. It does not single one out, and a function on the mass shell with values in a fixed space can only be compared across momenta through such a choice. The pure boost, \textit{i.e.}, the Hermitian positive factor of the polar decomposition, is the canonical candidate, and the freedom left is exactly the stabilizer of $\mathring{p}$.

\begin{proposition}[Stabilizer and standard boosts]\label{prop:orbit}
	Let $m>0$.
	\begin{enumerate}
		\item The group $\Lambda(SL(2,\mathbb{C}))$ acts transitively on $\mathcal{O}_m^+$, and the stabilizer of $\mathring{p}$ in $SL(2,\mathbb{C})$ is $SU(2)$.
		\item For $p\in\mathcal{O}_m^+$ let $\mathbf{w}_p:=\operatorname{arsinh}(|\mathbf{p}|/m)\,\mathbf{p}/|\mathbf{p}|$ be the rapidity vector of $p$, with $\mathbf{w}_{\mathring{p}}:=\mathbf{0}$, and let $A_p:=\exp(\tfrac12\mathbf{w}_p\cdot\boldsymbol{\sigma})$. Then $A_p\in SL(2,\mathbb{C})$ is Hermitian and positive definite, $p\mapsto A_p$ is smooth, $A_{\mathring{p}}=\mathbb{1}$, $\Lambda(A_p)\mathring{p}=p$, and
		\begin{equation}\label{eq:standard-boost}
			A_p=\Big(\frac{\hat{p}}{m}\Big)^{1/2}=\frac{\hat{p}+m\mathbb{1}}{\sqrt{2m(p^0+m)}},\qquad A_p^2=\frac{\hat{p}}{m},\qquad A_p^{-2}=\frac{\tilde{p}}{m}.
		\end{equation}
		\item Every $g\in SL(2,\mathbb{C})$ is written in a unique way as $g=A_pu$ with $p\in\mathcal{O}_m^+$ and $u\in SU(2)$, namely $p=\Lambda(g)\mathring{p}$. The map $g\mapsto(p,u)$ is a diffeomorphism $SL(2,\mathbb{C})\to\mathcal{O}_m^+\times SU(2)$, and $gSU(2)\mapsto\Lambda(g)\mathring{p}$ identifies $SL(2,\mathbb{C})/SU(2)$ with $\mathcal{O}_m^+$.
	\end{enumerate}
\end{proposition}

\begin{proof}
	Item 1. Transitivity is item 1 of Lemma~\ref{lem:hyperboloid-orbits}. Moreover $\Lambda(g)\mathring{p}=\mathring{p}$ if and only if $g\,(m\mathbb{1})\,g^\dagger=m\mathbb{1}$, \textit{i.e.}, $gg^\dagger=\mathbb{1}$, which together with $\det g=1$ means $g\in SU(2)$.
	
	Item 2. Put $s:=|\mathbf{p}|/m$, so that $\sinh|\mathbf{w}_p|=s$ and $\cosh|\mathbf{w}_p|=\sqrt{1+s^2}=p^0/m$. It descends that $\hat{p}/m=\cosh|\mathbf{w}_p|\,\mathbb{1}+\sinh|\mathbf{w}_p|\,\tfrac{\mathbf{w}_p}{|\mathbf{w}_p|}\cdot\boldsymbol{\sigma}=\exp(\mathbf{w}_p\cdot\boldsymbol{\sigma})$, and $A_p$ is its positive square root. As the exponential of a traceless Hermitian matrix, $A_p$ is Hermitian, positive definite and of determinant one, and it is smooth in $p$ as the positive square root of a positive definite matrix depending smoothly on $p$. The Cayley--Hamilton identity $\hat{p}^2=2p^0\hat{p}-m^2\mathbb{1}$ yields $(\hat{p}+m\mathbb{1})^2=2(p^0+m)\hat{p}$, and the explicit formula follows by uniqueness of the positive square root. Moreover $\widehat{\Lambda(A_p)\mathring{p}}=A_p\,(m\mathbb{1})\,A_p=mA_p^2=\hat{p}$, and $A_p^{-2}=m\hat{p}^{-1}=\tilde{p}/m$ by Lemma~\ref{lem:hermitian-model}.
	
	Item 3. By Lemma~\ref{lem:polar}, every $g$ is uniquely $\exp(\tfrac12\mathbf{w}\cdot\boldsymbol{\sigma})\,u$. Let $p:=m(\cosh|\mathbf{w}|,\ \sinh|\mathbf{w}|\,\mathbf{w}/|\mathbf{w}|)\in\mathcal{O}_m^+$, so that $\mathbf{w}=\mathbf{w}_p$. Then $\hat{p}=m\,e^{\mathbf{w}\cdot\boldsymbol{\sigma}}=m\,gg^\dagger$, hence $p=\Lambda(g)\mathring{p}$ and $\exp(\tfrac12\mathbf{w}\cdot\boldsymbol{\sigma})=A_p$. The map $\mathbf{w}\mapsto p$ is a diffeomorphism $\mathbb{R}^3\to\mathcal{O}_m^+$ with inverse $p\mapsto\mathbf{w}_p$, so $g\mapsto(p,u)$ is the composition of the diffeomorphism of Lemma~\ref{lem:polar} with a diffeomorphism. The map $gSU(2)\mapsto\Lambda(g)\mathring{p}$ is constant on cosets and has the smooth inverse $p\mapsto A_pSU(2)$.
\end{proof}

\begin{remark}
	Item 3 exhibits $SL(2,\mathbb{C})$ as the manifold $\mathcal{O}_m^+\times SU(2)$, with the copy $\{A_pu\mid u\in SU(2)\}$ standing above $p$. The decomposition is not compatible with the product. The set of standard boosts is not a subgroup, since $A_qA_{p'}$ is Hermitian only if it equals $A_{p'}A_q$, \textit{i.e.}, only if $\hat{q}$ and $\hat{p}'$ commute, and this holds exactly for collinear momenta by Proposition~\ref{prop:wigner-rotation}. In the Galilean case the boosts do form a subgroup acting simply transitively on the momenta, and this is why no analogue of the construction below is needed there.
\end{remark}

A state will be a square integrable function on the mass shell, so a measure compatible with the Lorentz action is needed. The flat measure $d^3\mathbf{p}$ is not invariant, and the correction is fixed by the way $\omega(\mathbf{p})$ transforms under a boost.

\begin{lemma}[Invariant measure]\label{lem:invariant-measure}
	The measure $d\mu(p):=d^3\mathbf{p}\,/\,2\omega(\mathbf{p})$ on $\mathcal{O}_m^+\cong\mathbb{R}^3$ is invariant under $\Lambda(g)$ for every $g\in SL(2,\mathbb{C})$.
\end{lemma}

\begin{proof}
	By Lemma~\ref{lem:polar}, $g=\exp(\tfrac12\mathbf{w}\cdot\boldsymbol{\sigma})\,u$, and for $\mathbf{w}\neq\mathbf{0}$ there is $u_1\in SU(2)$ with $\mathbf{w}\cdot\boldsymbol{\sigma}=u_1\,(|\mathbf{w}|\sigma_3)\,u_1^{-1}$. Hence $\Lambda(g)$ is a product of rotations and of boosts $\Lambda(\exp(\tfrac12 s\sigma_3))$ along $\mathbf{e}_3$. A rotation acts by $\mathbf{p}\mapsto R_u\mathbf{p}$, with $R_u$ as per Lemma~\ref{lem:poincare-relations}, and preserves $d^3\mathbf{p}$ and $\omega$. For the boost, $x^0\pm x^3\mapsto e^{\pm s}(x^0\pm x^3)$ yields $p'^3=\omega(\mathbf{p})\sinh s+p^3\cosh s$, with $p'^1=p^1$ and $p'^2=p^2$. Then $\partial p'^3/\partial p^3=\cosh s+(p^3/\omega)\sinh s=p'^0/\omega(\mathbf{p})$, and $p'^0=\omega(\mathbf{p}')$ by Lemma~\ref{lem:hyperboloid-orbits}. It descends that $d^3\mathbf{p}'=\tfrac{\omega(\mathbf{p}')}{\omega(\mathbf{p})}\,d^3\mathbf{p}$, \textit{i.e.}, $d^3\mathbf{p}'/\omega(\mathbf{p}')=d^3\mathbf{p}/\omega(\mathbf{p})$.
\end{proof}

\subsubsection*{Fibres and Sections}

A state with spin assigns a vector to each momentum, but the vectors at different momenta live in spaces that are identified with a fixed space only after a frame is chosen. The geometric object recording this is a bundle over the mass shell, whose total space is $SL(2,\mathbb{C})$ itself.

\begin{definition}[Principal bundle of the mass shell]\label{def:principal-bundle}
	The \emph{principal bundle of the mass shell} is the projection
	\begin{equation*}
		\varpi:SL(2,\mathbb{C})\to\mathcal{O}_m^+,\qquad\varpi(g):=\Lambda(g)\mathring{p},
	\end{equation*}
	with total space $SL(2,\mathbb{C})$, base $\mathcal{O}_m^+$ as per Definition~\ref{def:mass-shell}, and structure action the right multiplication of $SU(2)$ on $SL(2,\mathbb{C})$.
\end{definition}

\begin{remark}
	By Proposition~\ref{prop:orbit}, $\varpi(gu)=\varpi(g)$ for $u\in SU(2)$, and the fibre $\varpi^{-1}(p)=A_pSU(2)$ is the set of Lorentz transformations taking the rest momentum to $p$. The structure action is free and transitive on each fibre. Moreover $\varpi(g'g)=\Lambda(g')\varpi(g)$, so the left action of $SL(2,\mathbb{C})$ on the total space covers its action on the base, and it commutes with the structure action. The geometric content is the following. The stabilizer $SO(3)$ of $\mathring{p}$ in $SO^\uparrow(1,3)$ acts simply transitively on the orthonormal frames of the tangent space at $\mathring{p}$. Since $SO^\uparrow(1,3)$ acts transitively by isometries on $\mathcal{O}_m^+$, it acts simply transitively on the oriented orthonormal frames of $\mathcal{O}_m^+$, and $\varpi$ is the double covering of this frame bundle induced by Lemma~\ref{lem:sl2c-covering}. A point of the fibre over $p$ is a frame at $p$ up to the sign ambiguity of the spin, and $SU(2)$ is the group of changes of frame.
\end{remark}

A state with spin should attach to $p$ a vector of a space on which changes of frame act through a representation. The spin space is a representation $(\pi,V)$ of $SU(2)$, and the bundle below attaches to each $p$ a copy of $V$ which depends on no choice of frame.

\begin{definition}[Associated bundle]\label{def:associated-bundle}
	Let $(\pi,V)$ be a finite-dimensional unitary representation of $SU(2)$, to be specialized in Definition~\ref{def:spin-data}. The \emph{associated bundle} is the quotient
	\begin{equation*}
		E:=SL(2,\mathbb{C})\times_{SU(2)}V:=\big(SL(2,\mathbb{C})\times V\big)/\!\sim,
	\end{equation*}
	where $\sim$ is the equivalence relation generated by $(gu,v)\sim(g,\pi(u)v)$ with $g\in SL(2,\mathbb{C})$, $u\in SU(2)$ and $v\in V$. The class of $(g,v)$ is written $[g,v]$, and the projection is $\varpi_E:E\to\mathcal{O}_m^+$, $\varpi_E([g,v]):=\varpi(g)$.
\end{definition}

\begin{remark}
	For fixed $g\in SL(2,\mathbb{C})$ the map $v\mapsto[g,v]$ is a bijection of $V$ onto the fibre $E_p:=\varpi_E^{-1}(p)$ with $p=\varpi(g)$. Injectivity holds because the structure action is free, and surjectivity because the fibre of $\varpi$ is a single orbit. Since the relation is linear in $v$ and $\pi$ is unitary, $E_p$ inherits a vector space structure and a scalar product that do not depend on $g$. The identity $(gu,\pi(u)^{-1}v)\sim(g,v)$ shows how components depend on the frame. A vector of $E_p$ is a pair formed by a frame at $p$ and components in $V$, modulo the simultaneous change of frame by $u$ and of components by $\pi(u)^{-1}$. The space $E_p$ is therefore a copy of $V$ attached to $p$ and identified with $V$ only after a frame is chosen.
\end{remark}

Choosing a vector in $E_p$ for every $p$ is the geometric meaning of a state with spin. This choice has a description on the total space which refers to no frame, and the equivariance condition below expresses it.

\begin{definition}[Equivariant functions]\label{def:equivariant}
	The space $\mathcal{E}_\pi$ is the space of Borel functions $f:SL(2,\mathbb{C})\to V$ such that
	\begin{equation}\label{eq:equivariance}
		f(gu)=\pi(u^{-1})\,f(g),\qquad g\in SL(2,\mathbb{C}),\ u\in SU(2).
	\end{equation}
\end{definition}

The condition is not an ad hoc requirement. It is exactly the condition under which a function on the total space records one vector of $E_p$ for each $p$, as the following theorem shows.

\begin{theorem}[Equivariant functions are sections]\label{thm:equivariant-sections}
	A \emph{section} of $E$ is a map $s:\mathcal{O}_m^+\to E$ with $\varpi_E\circ s=\mathrm{id}$. The assignment $s\mapsto f_s$, where $f_s(g)$ is the unique $v\in V$ with $s(\varpi(g))=[g,v]$, is a bijection from the sections of $E$ onto the functions $f:SL(2,\mathbb{C})\to V$ satisfying Equation~\eqref{eq:equivariance}, with inverse $f\mapsto s_f$, $s_f(\varpi(g)):=[g,f(g)]$. Declaring $s$ Borel when $f_s$ is Borel, the Borel sections correspond to $\mathcal{E}_\pi$.
\end{theorem}

\begin{proof}
	The function $f_s$ is well defined by the remark after Definition~\ref{def:associated-bundle}. Since $\varpi(gu)=\varpi(g)$,
	\begin{equation*}
		[g,f_s(g)]=s(\varpi(g))=s(\varpi(gu))=[gu,f_s(gu)]=[g,\pi(u)f_s(gu)],
	\end{equation*}
	and injectivity of $v\mapsto[g,v]$ yields $f_s(g)=\pi(u)f_s(gu)$, which is Equation~\eqref{eq:equivariance}. Conversely, if $f$ satisfies Equation~\eqref{eq:equivariance}, then $[gu,f(gu)]=[gu,\pi(u^{-1})f(g)]=[g,f(g)]$. The class $[g,f(g)]$ depends only on $\varpi(g)$, since the fibre of $\varpi$ is a single orbit, and $s_f$ is a section. The two assignments are inverse to each other by construction.
\end{proof}

The group acts on sections by transporting the vector at $\Lambda(g)^{-1}p$ to the point $p$. On $E$ this transport is $[g',v]\mapsto[gg',v]$, and the next proposition expresses it on equivariant functions, where it becomes a plain translation.

\begin{proposition}[Left translation]\label{prop:action-equivariant}
	For $g\in SL(2,\mathbb{C})$ and $f\in\mathcal{E}_\pi$ let $(\mathcal{L}(g)f)(g'):=f(g^{-1}g')$. Then $\mathcal{L}(g)f\in\mathcal{E}_\pi$, the map $\mathcal{L}$ is a representation of $SL(2,\mathbb{C})$ on $\mathcal{E}_\pi$, and under the bijection of Theorem~\ref{thm:equivariant-sections} the operator $\mathcal{L}(g)$ corresponds to the action on sections
	\begin{equation*}
		(g\cdot s)(p):=g\cdot s\big(\Lambda(g)^{-1}p\big),\qquad g\cdot[g',v]:=[gg',v].
	\end{equation*}
\end{proposition}

\begin{proof}
	The action $g\cdot[g',v]=[gg',v]$ is well defined since left multiplication commutes with right multiplication by $SU(2)$. Let $s$ correspond to $f$. Since $\Lambda(g)^{-1}\varpi(g')=\varpi(g^{-1}g')$,
	\begin{equation*}
		(g\cdot s)(\varpi(g'))=g\cdot[g^{-1}g',f(g^{-1}g')]=[g',f(g^{-1}g')],
	\end{equation*}
	so that $g\cdot s$ corresponds to $g'\mapsto f(g^{-1}g')$, which satisfies Equation~\eqref{eq:equivariance} again because left and right multiplications commute. Finally $(\mathcal{L}(g)\mathcal{L}(\tilde{g})f)(g')=f(\tilde{g}^{-1}g^{-1}g')=(\mathcal{L}(g\tilde{g})f)(g')$.
\end{proof}

\begin{remark}
	In this description no frame has been chosen, and the group acts by translation with no matrix. Matrices appear only when the vectors of $E_p$ are read as elements of $V$. This requires a frame above every momentum, and the whole content of the Wigner rotation lies in the dependence of the result on this choice.
\end{remark}

\subsubsection*{Frames and Gauge}

The wave functions of the chapter are functions of the momentum alone, not of a point of $SL(2,\mathbb{C})$. Passing from the first to the second requires one frame above each momentum, \textit{i.e.}, a section of $\varpi$. Such a choice is a choice of gauge.

\begin{definition}[Section of the principal bundle]\label{def:section-principal}
	A \emph{section} of $\varpi$ is a smooth map $s_0:\mathcal{O}_m^+\to SL(2,\mathbb{C})$ with $\varpi\circ s_0=\mathrm{id}$.
\end{definition}

\begin{remark}
	A section is the choice of a frame at each momentum. By Proposition~\ref{prop:orbit} the standard boosts $p\mapsto A_p$ form a section, characterized by taking at each $p$ the pure boost with no rotation. If $s_0$ and $s_0'$ are two sections, then $\varpi(s_0(p)^{-1}s_0'(p))=\Lambda(s_0(p))^{-1}p=\mathring{p}$, hence $k(p):=s_0(p)^{-1}s_0'(p)$ belongs to $SU(2)$, and $s_0'=s_0k$ with $k:\mathcal{O}_m^+\to SU(2)$ smooth. Every such product is a section. A change of gauge is thus a rotation of the frame depending on the momentum.
\end{remark}

For $s_0=A$ the following lemma is item 3 of Proposition~\ref{prop:orbit}. The standard boosts have no special status in it, since any section identifies the total space with the product of base and fibre.

\begin{lemma}[Trivialization]\label{lem:trivialization}
	Let $s_0$ be a section of $\varpi$. The map $\Theta_{s_0}:\mathcal{O}_m^+\times SU(2)\to SL(2,\mathbb{C})$, $\Theta_{s_0}(p,u):=s_0(p)u$, is a diffeomorphism with inverse $g\mapsto(\varpi(g),s_0(\varpi(g))^{-1}g)$. Moreover $(p,v)\mapsto[s_0(p),v]$ is a bijection of $\mathcal{O}_m^+\times V$ onto $E$, linear on each fibre.
\end{lemma}

\begin{proof}
	For $g\in SL(2,\mathbb{C})$ and $p=\varpi(g)$ one has $\varpi(s_0(p)^{-1}g)=\Lambda(s_0(p))^{-1}p=\mathring{p}$, so $s_0(p)^{-1}g\in SU(2)$ by Proposition~\ref{prop:orbit}. Both maps are smooth and inverse to each other. The second statement follows from the remark after Definition~\ref{def:associated-bundle}.
\end{proof}

The trivialization lets us trade an equivariant function on the group for a function on the mass shell. The price is that the result depends on the section, as the last statement of the proposition shows.

\begin{proposition}[Equivariant functions as functions on the base]\label{prop:equivariant-base}
	Let $s_0$ be a section of $\varpi$. The map $f\mapsto\psi_f$, $\psi_f(p):=f(s_0(p))$, is a bijection from $\mathcal{E}_\pi$ onto the Borel functions $\mathcal{O}_m^+\to V$, with inverse $\psi\mapsto f_\psi$, $f_\psi(s_0(p)u):=\pi(u^{-1})\psi(p)$. If $s_0'=s_0k$ is another section, then $\psi_f'(p)=\pi(k(p)^{-1})\psi_f(p)$.
\end{proposition}

\begin{proof}
	By Lemma~\ref{lem:trivialization} every $g$ is uniquely $s_0(p)u$, so $f_\psi$ is well defined and Borel. It satisfies Equation~\eqref{eq:equivariance}, since $f_\psi(s_0(p)uu')=\pi(u'^{-1})\pi(u^{-1})\psi(p)=\pi(u'^{-1})f_\psi(s_0(p)u)$. Conversely, Equation~\eqref{eq:equivariance} yields $f(s_0(p)u)=\pi(u^{-1})f(s_0(p))$, so $f=f_{\psi_f}$. For the last claim, $\psi_f'(p)=f(s_0(p)k(p))=\pi(k(p)^{-1})f(s_0(p))$.
\end{proof}

\begin{remark}
	The vector $\psi_f(p)$ is the vector of $E_p$ read in the frame $s_0(p)$, and a change of gauge rotates the components by $\pi(k(p)^{-1})$. Since $\pi$ is unitary, $\|f(g)\|$ depends only on $\varpi(g)$, and $\|\psi_f(p)\|$ does not depend on the gauge. The Hilbert space $L^2(\mathcal{O}_m^+,\mu;V)$ of Definition~\ref{def:induced-rep} is thus the space of square integrable sections of $E$, with a norm attached to the mass shell and not to a frame.
\end{remark}

\subsubsection*{The Wigner Rotation}

Consider the translation $\mathcal{L}(g)$ read in a gauge $s_0$. The value at $p$ of the transformed function comes from the value at $q:=\Lambda(g)^{-1}p$ of the original one. The frame at $q$ is $s_0(q)$, and the group carries it to $g\,s_0(q)$, a frame above $\Lambda(g)q=p$. The frames above $p$ form the single copy $s_0(p)SU(2)$, hence $g\,s_0(q)$ differs from the reference frame $s_0(p)$ by a unique rotation. This rotation is the transition function.

\begin{definition}[Transition function]\label{def:transition-function}
	Let $s_0$ be a section of $\varpi$. The \emph{transition function} of $s_0$ is the map $h:SL(2,\mathbb{C})\times\mathcal{O}_m^+\to SU(2)$ defined by
	\begin{equation*}
		g^{-1}s_0(p)=s_0\big(\Lambda(g)^{-1}p\big)\,h(g,p).
	\end{equation*}
\end{definition}

\begin{remark}
	The element $h(g,p)$ exists and is unique by Lemma~\ref{lem:trivialization}, because $\varpi(g^{-1}s_0(p))=\Lambda(g)^{-1}p$. The formula $h(g,p)=s_0(\Lambda(g)^{-1}p)^{-1}g^{-1}s_0(p)$ shows that $h$ is smooth. Equivalently, $g\,s_0(q)=s_0(p)\,h(g,p)^{-1}$ with $q=\Lambda(g)^{-1}p$. The group carries the point $q$ to the point $p$ and the reference frame at $q$ to a frame at $p$ which is not the reference frame at $p$. The two differ by the rotation $h(g,p)^{-1}$.
\end{remark}

The transition function is not an arbitrary family of rotations. It composes consistently along products of group elements, and a change of gauge modifies it in a controlled way. The two statements below are what allows the action of the group to be written on functions of the momentum.

\begin{lemma}[Cocycle identity and change of gauge]\label{lem:transition-cocycle}
	Let $h$ be the transition function of a section $s_0$.
	\begin{enumerate}
		\item $h(\mathbb{1},p)=\mathbb{1}$, and $h(gg',p)=h\big(g',\Lambda(g)^{-1}p\big)\,h(g,p)$ for all $g,g'\in SL(2,\mathbb{C})$.
		\item If $s_0'=s_0k$ is another section with transition function $h'$, then
		\begin{equation*}
			h'(g,p)=k\big(\Lambda(g)^{-1}p\big)^{-1}\,h(g,p)\,k(p).
		\end{equation*}
	\end{enumerate}
\end{lemma}

\begin{proof}
	Item 1. The identity $h(\mathbb{1},p)=\mathbb{1}$ is immediate. With $q=\Lambda(g)^{-1}p$, Definition~\ref{def:transition-function} applied twice yields
	\begin{equation*}
		(gg')^{-1}s_0(p)=g'^{-1}s_0(q)\,h(g,p)=s_0\big(\Lambda(gg')^{-1}p\big)\,h(g',q)\,h(g,p),
	\end{equation*}
	and uniqueness in Definition~\ref{def:transition-function} yields the claim.
	
	Item 2. Since $s_0(q)=s_0'(q)k(q)^{-1}$,
	\begin{equation*}
		g^{-1}s_0'(p)=g^{-1}s_0(p)\,k(p)=s_0'(q)\,k(q)^{-1}h(g,p)k(p).\qedhere
	\end{equation*}
\end{proof}

\begin{remark}
	Item 1 says that $h$ is a cocycle for the action of $SL(2,\mathbb{C})$ on $\mathcal{O}_m^+$ with values in $SU(2)$, and item 2 says that a change of gauge replaces it by a cohomologous one. This is the non-abelian counterpart of the multiplier of a projective representation and of its rephasing in Chapter~\ref{ch:analytical-hamiltonians}, Proposition~\ref{prop:rephasing}, with rotations of the frame in place of phases. In both cases the cocycle measures the failure of a section to be a homomorphism, and it is defined only up to a change of section.
\end{remark}

The cocycle identity is the compatibility needed to write the action of the group on functions of the momentum. The next proposition does so in an arbitrary gauge, and the transition function is the only datum which appears.

\begin{proposition}[The action in a gauge]\label{prop:action-trivialization}
	Let $s_0$ be a section of $\varpi$ with transition function $h$, and let $f\in\mathcal{E}_\pi$. Then, for all $g\in SL(2,\mathbb{C})$ and $p\in\mathcal{O}_m^+$,
	\begin{equation*}
		\psi_{\mathcal{L}(g)f}(p)=\pi\big(h(g,p)^{-1}\big)\,\psi_f\big(\Lambda(g)^{-1}p\big),
	\end{equation*}
	with $\psi_f$ as per Proposition~\ref{prop:equivariant-base}.
\end{proposition}

\begin{proof}
	By Definition~\ref{def:transition-function} and Equation~\eqref{eq:equivariance}, with $q=\Lambda(g)^{-1}p$,
	\begin{equation*}
		\psi_{\mathcal{L}(g)f}(p)=f\big(g^{-1}s_0(p)\big)=f\big(s_0(q)h(g,p)\big)=\pi\big(h(g,p)^{-1}\big)f\big(s_0(q)\big).\qedhere
	\end{equation*}
\end{proof}

On functions of the momentum the group therefore acts by relocation of the argument composed with a rotation of the value, and the rotation is the only trace of the frame. For the standard boosts the transition function is the inverse of an explicit matrix, which is the Wigner rotation.

\begin{proposition}[Induced action and Wigner rotation]\label{prop:induced-action}
	For $g\in SL(2,\mathbb{C})$ and $p\in\mathcal{O}_m^+$ set
	\begin{equation}\label{eq:wigner-rotation}
		R_W(g,p):=A_p^{-1}\,g\,A_{\Lambda(g)^{-1}p},
	\end{equation}
	and call it the \emph{Wigner rotation}.
	\begin{enumerate}
		\item $R_W(g,p)=h(g,p)^{-1}$ for the transition function $h$ of the section $p\mapsto A_p$. In particular $R_W(g,p)\in SU(2)$, $R_W(\mathbb{1},p)=\mathbb{1}$ and $R_W(gg',p)=R_W(g,p)\,R_W(g',\Lambda(g)^{-1}p)$.
		\item For every $f\in\mathcal{E}_\pi$, with $\psi_f(p)=f(A_p)$,
		\begin{equation}\label{eq:induced-action}
			\psi_{\mathcal{L}(g)f}(p)=\pi\big(R_W(g,p)\big)\,\psi_f\big(\Lambda(g)^{-1}p\big).
		\end{equation}
		\item $R_W(u,p)=u$ for every $u\in SU(2)$.
	\end{enumerate}
\end{proposition}

\begin{proof}
	Item 1. With $q=\Lambda(g)^{-1}p$, Definition~\ref{def:transition-function} for $s_0(p)=A_p$ reads $g^{-1}A_p=A_q\,h(g,p)$, hence $h(g,p)=A_q^{-1}g^{-1}A_p=(A_p^{-1}gA_q)^{-1}$. The remaining properties are item 1 of Lemma~\ref{lem:transition-cocycle} written for $R_W=h^{-1}$.
	
	Item 2 is Proposition~\ref{prop:action-trivialization} combined with item 1.
	
	Item 3. For $u\in SU(2)$ the matrix $u^{-1}A_pu$ is positive definite with square $u^{-1}\hat{p}u/m=\widehat{\Lambda(u)^{-1}p}/m$, hence $A_{\Lambda(u)^{-1}p}=u^{-1}A_pu$ by Equation~\eqref{eq:standard-boost}. It descends that $R_W(u,p)=A_p^{-1}uu^{-1}A_pu=u$.
\end{proof}

By item 3, rotations act on the spin by the rotation itself, with the same matrix at every momentum. The remaining elements of $SL(2,\mathbb{C})$ are boosts, and for them the Wigner rotation depends on the momentum. The next proposition shows it for the standard boost $A_q$, where the comparison between transported and standard frames reduces to a product of two standard boosts.

\begin{proposition}[Momentum dependence of the Wigner rotation]\label{prop:wigner-rotation}
	Let $q\in\mathcal{O}_m^+$ with $\mathbf{q}\neq\mathbf{0}$, let $p'\in\mathcal{O}_m^+$ and $p:=\Lambda(A_q)p'$. Then $R_W(A_q,p)=\mathbb{1}$ if and only if $\mathbf{q}\times\mathbf{p}'=\mathbf{0}$. In particular $p\mapsto R_W(A_q,p)$ is not constant.
\end{proposition}

\begin{proof}
	Since $\Lambda(A_q)^{-1}p=p'$, $R_W(A_q,p)=A_p^{-1}A_qA_{p'}$, which equals $\mathbb{1}$ if and only if $A_p=A_qA_{p'}$. If this holds, $A_qA_{p'}$ is Hermitian, hence $A_qA_{p'}=(A_qA_{p'})^\dagger=A_{p'}A_q$, so $\hat{q}=mA_q^2$ and $\hat{p}'=mA_{p'}^2$ commute. From $[\mathbf{a}\cdot\boldsymbol{\sigma},\mathbf{b}\cdot\boldsymbol{\sigma}]=2i(\mathbf{a}\times\mathbf{b})\cdot\boldsymbol{\sigma}$ it descends that $\mathbf{q}\times\mathbf{p}'=\mathbf{0}$. Conversely, if $\mathbf{q}\times\mathbf{p}'=\mathbf{0}$, then $\hat{q}$ and $\hat{p}'$ commute, and so do $A_q$ and $A_{p'}$, which are affine in them by Equation~\eqref{eq:standard-boost}. Then $A_qA_{p'}$ is Hermitian and positive definite of determinant one, and $\Lambda(A_qA_{p'})\mathring{p}=\Lambda(A_q)p'=p$, so $A_qA_{p'}=A_p$ by uniqueness in item 3 of Proposition~\ref{prop:orbit}. For the last claim, $p'=\mathring{p}$ yields $p=q$ and $R_W(A_q,q)=\mathbb{1}$, whereas every $p'$ with $\mathbf{p}'$ not parallel to $\mathbf{q}$ yields $R_W(A_q,p)\neq\mathbb{1}$.
\end{proof}

The proposition shows that the Wigner rotation of a boost is trivial only for collinear momenta. The following example computes it for perpendicular momenta, where the rotation is a single explicit matrix.

\begin{example}[Perpendicular boosts and the Thomas--Wigner rotation]\label{ex:thomas-wigner}
	Let $w_1,w_2>0$, and let $q,p'\in\mathcal{O}_m^+$ have rapidity vectors $w_1\mathbf{e}_1$ and $w_2\mathbf{e}_2$, so that $A_q=\exp(\tfrac12w_1\sigma_1)$ and $A_{p'}=\exp(\tfrac12w_2\sigma_2)$. Let $p:=\Lambda(A_q)p'$, $c_i:=\cosh(w_i/2)$ and $s_i:=\sinh(w_i/2)$. Then
	\begin{equation*}
		g:=A_qA_{p'}=c_1c_2\,\mathbb{1}+s_1c_2\,\sigma_1+c_1s_2\,\sigma_2+i\,s_1s_2\,\sigma_3.
	\end{equation*}
	By item 1 of Proposition~\ref{prop:induced-action}, $g=A_p\,R_W(A_q,p)$ with $A_p$ Hermitian positive definite and $R_W(A_q,p)\in SU(2)$, which is the polar decomposition of Lemma~\ref{lem:polar}. Define $\vartheta\in(0,\pi/2)$ by $\tan\tfrac\vartheta2=\tanh\tfrac{w_1}{2}\tanh\tfrac{w_2}{2}$ and $U:=\exp(\tfrac i2\vartheta\sigma_3)$. The coefficient of $\sigma_3$ in $gU^{-1}$ is $i(s_1s_2\cos\tfrac\vartheta2-c_1c_2\sin\tfrac\vartheta2)=0$, the coefficients of $\sigma_1$ and $\sigma_2$ are real, and the scalar part $c_1c_2\cos\tfrac\vartheta2+s_1s_2\sin\tfrac\vartheta2$ is positive. Hence $gU^{-1}$ is Hermitian, positive definite and of determinant one, and uniqueness in Lemma~\ref{lem:polar} yields $R_W(A_q,p)=U$. This is the rotation $R_U$ of angle $-\vartheta$ about $\mathbf{e}_3$, orthogonal to the plane of $\mathbf{q}$ and $\mathbf{p}'$, with $\vartheta\simeq w_1w_2/2$ for small rapidities. It is the Thomas rotation, and it is the elementary reason why the spin of a relativistic particle cannot be given independently of its momentum.
\end{example}

\begin{remark}[The Wigner angle as an area]
	Equation~\eqref{eq:wigner-rotation} reads $A_qA_{p'}=A_p\,R_W(A_q,p)$, which has the form of the defect of a section in Chapter~\ref{ch:analytical-hamiltonians}, now for the section $p\mapsto A_p$ of $\varpi$ and with values in the non-central group $SU(2)$. The analogy extends to the geometry. Rescale $\mathcal{O}_m^+$ to unit radius, so that the geodesic distance between $\mathring{p}$ and $p$ is $|\mathbf{w}_p|$. In Example~\ref{ex:thomas-wigner} the points $\mathring{p}$, $q$ and $p$ are the vertices of a geodesic triangle with legs of length $w_1$ and $w_2$ and a right angle at $q$. Indeed $\Lambda(A_q)$ is an isometry carrying $(\mathring{p},p')$ to $(q,p)$, and it carries the directions $-\mathbf{e}_1$ and $\mathbf{e}_2$ at $\mathring{p}$ to the directions at $q$ towards $\mathring{p}$ and towards $p$. The area $\Delta$ of a right hyperbolic triangle with legs $a$ and $b$ satisfies $\tan\tfrac\Delta2=\tanh\tfrac a2\tanh\tfrac b2$ \cite[TODO-location]{TODO-Ratcliffe}, so $\vartheta$ is the area of this triangle. For a general triangle, the Wigner rotation is the holonomy of the Levi-Civita connection of $\mathcal{O}_m^+$ along its boundary, and its angle has modulus equal to the area \cite[TODO-location]{TODO-Aravind1997}. In Chapter~\ref{ch:analytical-hamiltonians} the defect of the Heisenberg group is the Euclidean area of a triangle, as per Example~\ref{ex:area-defect}, and here it is the hyperbolic one. The flat momentum space of the Galilean case has vanishing curvature, in agreement with the trivial action of the boosts on the spin in Corollary~\ref{cor:schrodinger-realization}.
\end{remark}

It remains to decide whether the dependence on $p$ is an artefact of the standard gauge. Since a change of gauge replaces $h$ by a cohomologous cocycle by Lemma~\ref{lem:transition-cocycle}, one could hope to find a section for which the matrices become independent of the momentum. The next proposition excludes it, and the obstruction is the absence of non-trivial unitary finite-dimensional representations of $SL(2,\mathbb{C})$.

\begin{proposition}[No gauge removes the Wigner rotation]\label{prop:wigner-no-gauge}
	Let $(\pi,V)$ be a non-trivial finite-dimensional unitary representation of $SU(2)$, let $s_0$ be a section of $\varpi$ and $h$ its transition function. Then there exist $g\in SL(2,\mathbb{C})$ and $p,p'\in\mathcal{O}_m^+$ with $\pi(h(g,p))\neq\pi(h(g,p'))$. In particular, for $\pi$ the defining representation, $h(g,p)\neq h(g,p')$.
\end{proposition}

\begin{proof}
	Suppose that $\pi(h(g,p))$ does not depend on $p$ for every $g$, and write $T(g):=\pi(h(g,\mathring{p}))\in U(V)$. Item 1 of Lemma~\ref{lem:transition-cocycle} yields $T(gg')=T(g')T(g)$, hence $g\mapsto T(g)^{-1}$ is a homomorphism $SL(2,\mathbb{C})\to U(V)$. It is continuous, since $h(g,\mathring{p})=s_0(\Lambda(g)^{-1}\mathring{p})^{-1}g^{-1}s_0(\mathring{p})$ is smooth in $g$. By Proposition~\ref{prop:no-unitary}, it is trivial, so $T(g)=\mathbb{1}$ for every $g$. For $g=u\in SU(2)$ one has $\Lambda(u)^{-1}\mathring{p}=\mathring{p}$, and Definition~\ref{def:transition-function} yields
	\begin{equation*}
		h(u,\mathring{p})=s_0(\mathring{p})^{-1}\,u^{-1}\,s_0(\mathring{p}),
	\end{equation*}
	where $s_0(\mathring{p})\in SU(2)$ because $\varpi(s_0(\mathring{p}))=\mathring{p}$. Since $u\mapsto s_0(\mathring{p})^{-1}u^{-1}s_0(\mathring{p})$ is a bijection of $SU(2)$, it descends that $\pi(u)=\mathbb{1}$ for every $u\in SU(2)$, which contradicts the non-triviality of $\pi$.
\end{proof}

\begin{remark}
	This is the geometric form of the statement that spin and momentum cannot be separated. In whichever frame the vectors of $E_p$ are read, the Lorentz group acts on the spin of a particle of momentum $q$ by a rotation depending on $q$. Spin zero is the exception, since for $\pi$ trivial $E=\mathcal{O}_m^+\times\mathbb{C}$ and scalar functions transform by relocation of the argument alone. The proposition concerns the form of the transformation law and not the existence of a representation. By Lemma~\ref{lem:invariant-measure}, the law of Equation~\eqref{eq:induced-action} is unitary on square integrable sections, as will be verified in Definition~\ref{def:induced-rep}. What is lost is the possibility of describing the spin by a fixed matrix at each momentum.
\end{remark}

The obstruction suggests where a local description must look instead. A transformation law with fixed matrices requires a representation $D$ of $SL(2,\mathbb{C})$ on a space $W$, acting on functions $\Psi:\mathcal{O}_m^+\to W$ by $(\mathcal{U}(g)\Psi)(p)=D(g)\Psi(\Lambda(g)^{-1}p)$.

\begin{remark}[Towards fields]
	By Proposition~\ref{prop:no-unitary} such a $D$ cannot be unitary unless it is trivial, and by Proposition~\ref{prop:wigner-no-gauge} the matrix $\pi(R_W(g,p))$ acting on the fibre $V$ cannot be made independent of $p$ in any gauge. The construction of Definition~\ref{def:local-realization} circumvents both facts by enlarging the fibre. The bundle $E$ is embedded as a subbundle of the trivial bundle $\mathcal{O}_m^+\times W$, the frame $A_p$ carrying a subspace $\iota(V)\subseteq W$ to $F(p):=D(A_p)\iota(V)$. The dependence on $p$ which cannot be placed in the matrix is placed in the subspace. The fibres $F(p)$ move with $p$, with $D(g)F(p)=F(\Lambda(g)p)$ by Remark~\ref{rem:local-subbundle}, and the Wigner rotation is encoded in the way they move. A relativistic wave equation is then the constraint $\Psi(p)\in F(p)$ which selects the subbundle, and for spin $\tfrac12$ it is the Dirac equation of Theorem~\ref{thm:dirac}. The price is that the Euclidean norm of $W$ is not invariant, and the invariant norm is the one of Proposition~\ref{prop:norms}.
\end{remark}

\subsection{Finite-dimensional representations and space inversion}
\label{subsec:finite-reps-inversion}

The states of an elementary system of spin $\frac12$ are represented, up to this point, by square integrable functions on the mass shell with values in $\mathbb{C}^2$. A Lorentz transformation acts on the value at the momentum $p$ through a Wigner rotation depending on $p$. Section 6.5 asks whether the same states admit a description in which the value of a field is transformed by a fixed matrix $D(g)$, independent of $p$, and which matrices are available for that purpose. The question belongs to group theory and involves neither the dynamics nor the Hilbert space. What is given is the group $SL(2,\mathbb{C})$ with its chiral decomposition, Lemma~\ref{lem:chiral}, and the semisimplicity of its Lie algebra, Corollary~\ref{cor:poincare-structure}. What is obtained is threefold. Space inversion is an automorphism of $SL(2,\mathbb{C})$ which fixes the rotations and exchanges each standard boost with its inverse. The finite-dimensional representations are labelled by pairs of spins, and exactly two of them restrict to the doublet of $SU(2)$, namely the two Weyl representations. No nontrivial finite-dimensional representation is unitary. These facts are the input of the obstruction to locality and of the doubling of the spinor in Section 6.5, and nothing else of this subsection is used elsewhere.

%--------------------------------------------------------------------
\begin{lemma}[Space inversion]\label{lem:inversion}
	Let $\mathsf{P}(x^0,\mathbf{x}):=(x^0,-\mathbf{x})$ for $x\in M$ and $\theta(g):=(g^\dagger)^{-1}$ for $g\in SL(2,\mathbb{C})$.
	\begin{enumerate}
		\item $\widehat{\mathsf{P}x}=\tilde{x}$. The map $\theta$ is an automorphism of $SL(2,\mathbb{C})$ with $\theta^2=\mathrm{id}$ and $\theta|_{SU(2)}=\mathrm{id}$, its differential is $d\theta(X)=-X^\dagger$, so that $d\theta(r_k)=r_k$ and $d\theta(b_k)=-b_k$, and
		\begin{equation}\label{eq:theta-lambda}
			\Lambda(\theta(g))=\mathsf{P}\,\Lambda(g)\,\mathsf{P}.
		\end{equation}
		\item $\mathsf{P}\mathcal{O}_m^+=\mathcal{O}_m^+$, and for every $p\in\mathcal{O}_m^+$
		\begin{equation}\label{eq:theta-boost}
			\theta(A_p)=A_{\mathsf{P}p}=A_p^{-1}.
		\end{equation}
		\item The measure $\mu$ of Lemma~\ref{lem:invariant-measure} is invariant under $\mathsf{P}$.
	\end{enumerate}
\end{lemma}

\begin{proof}
	Item 1. The first claim is Definition~\ref{def:hermitian-model}. Since $(gg')^\dagger=g'^\dagger g^\dagger$, the map $\theta$ is a homomorphism. It takes values in $SL(2,\mathbb{C})$ because $\det\theta(g)=\overline{\det g}^{\,-1}=1$. Moreover $\theta^2(g)=g$, and $\theta(u)=u$ for $u\in SU(2)$ since $u^\dagger=u^{-1}$. The differential follows from $\frac{d}{dt}(e^{tX\dagger})^{-1}|_{t=0}=-X^\dagger$, and $r_k^\dagger=-r_k$, $b_k^\dagger=b_k$ yield the values on $r_k$ and $b_k$. By item 3 of Lemma~\ref{lem:sl2c-covering},
	\begin{equation*}
		\widehat{\mathsf{P}\Lambda(g)x}=\widetilde{\Lambda(g)x}=(g^\dagger)^{-1}\tilde{x}\,g^{-1}=\theta(g)\,\widehat{\mathsf{P}x}\,\theta(g)^\dagger=\widehat{\Lambda(\theta(g))\mathsf{P}x},
	\end{equation*}
	since $\theta(g)^\dagger=g^{-1}$. It descends that $\mathsf{P}\Lambda(g)=\Lambda(\theta(g))\mathsf{P}$, which is Equation~\eqref{eq:theta-lambda}.
	
	Item 2. The map $\mathsf{P}$ preserves $p^2$ and $p^0$, hence $\mathsf{P}\mathcal{O}_m^+=\mathcal{O}_m^+$. By item 2 of Proposition~\ref{prop:orbit}, the rapidity vector of $p$ is $\mathbf{w}_p=\operatorname{arsinh}(|\mathbf{p}|/m)\,\mathbf{p}/|\mathbf{p}|$ and $\mathbf{w}_{\mathring{p}}=\mathbf{0}$, so that $\mathbf{w}_{\mathsf{P}p}=-\mathbf{w}_p$ and $A_{\mathsf{P}p}=\exp(-\tfrac12\mathbf{w}_p\cdot\boldsymbol{\sigma})=A_p^{-1}$. Since $A_p$ is Hermitian, $\theta(A_p)=(A_p^\dagger)^{-1}=A_p^{-1}$.
	
	Item 3. In the coordinate $\mathbf{p}$ the map $\mathsf{P}$ is $\mathbf{p}\mapsto-\mathbf{p}$, which preserves $d^3\mathbf{p}$ and $\omega(\mathbf{p})$.
\end{proof}

\begin{remark}[Space inversion is an outer automorphism]\label{rem:inversion-outer}
	Rotations are even under space inversion and boosts are odd, as $d\theta(r_k)=r_k$ and $d\theta(b_k)=-b_k$ express. The map $\theta$ is not inner, \textit{i.e.}, there is no $h\in SL(2,\mathbb{C})$ with $\theta(g)=hgh^{-1}$ for all $g$. Such an $h$ would commute with $SU(2)$, hence with the span $M_2(\mathbb{C})$ of $SU(2)$, so that $h=\pm\mathbb{1}$ and $\theta=\mathrm{id}$, which fails for $g=\operatorname{diag}(2,\tfrac12)$. Space inversion does not belong to $\mathcal{P}$, since $\mathsf{P}\notin SO^\uparrow(1,3)$, but it acts on it. By Equation~\eqref{eq:theta-lambda} the map $\vartheta(a,g):=(\mathsf{P}a,\theta(g))$ satisfies
	\begin{equation*}
		\vartheta(a,g)\,\vartheta(a',g')=\big(\mathsf{P}a+\mathsf{P}\Lambda(g)a',\,\theta(gg')\big)=\vartheta\big((a,g)(a',g')\big),
	\end{equation*}
	with $\Lambda(\theta(g))\mathsf{P}a'=\mathsf{P}\Lambda(g)a'$. Hence $\vartheta$ is an automorphism of $\mathcal{P}$ as per Definition~\ref{def:poincare-group}, and it is the one implemented by the parity operators of Section 6.5.
\end{remark}

The automorphism $\theta$ is the only new ingredient needed to name the simplest representations of $SL(2,\mathbb{C})$ beyond the defining one. Composing the defining representation with $\theta$ yields a second two-dimensional representation, which coincides with the first on $SU(2)$ and differs from it on the boosts.

\begin{definition}[Weyl representations]\label{def:weyl-spinors}
	On $\mathbb{C}^2$ the \emph{left-handed} and the \emph{right-handed Weyl representation} of $SL(2,\mathbb{C})$ are
	\begin{equation*}
		D^{(1/2,0)}(g):=g, \qquad D^{(0,1/2)}(g):=\theta(g)=(g^\dagger)^{-1}.
	\end{equation*}
\end{definition}

\begin{remark}[Differentials of the Weyl representations]\label{rem:weyl-differentials}
	The differentials are $dD^{(1/2,0)}(X)=X$ and $dD^{(0,1/2)}(X)=-X^\dagger$. By Lemma~\ref{lem:chiral} and Lemma~\ref{lem:inversion}, $dD^{(1/2,0)}$ sends $A_k\mapsto-\tfrac{i}{2}\sigma_k$ and $B_k\mapsto0$, whereas $dD^{(0,1/2)}$ sends $A_k\mapsto0$ and $B_k\mapsto-\tfrac{i}{2}\sigma_k$. In the first only the factor generated by the $A_k$ acts, with spin $\frac12$, and in the second only the one generated by the $B_k$. Both representations restrict to the defining representation of $SU(2)$, because $\theta|_{SU(2)}=\mathrm{id}$, and $D^{(1/2,0)}\circ\theta=D^{(0,1/2)}$, $D^{(0,1/2)}\circ\theta=D^{(1/2,0)}$, because $\theta^2=\mathrm{id}$.
\end{remark}

The two Weyl representations are the first members of a family labelled by pairs of spins, one for each of the commuting ideals of Lemma~\ref{lem:chiral}. The theorem below states that the family exhausts the finite-dimensional representations and records how restriction to the rotations and composition with $\theta$ act on the labels. Here $V_j:=\mathbb{C}^{2j+1}$, $d_j$ is the irreducible unitary representation of $SU(2)$ as per Definition~\ref{def:spin-data}, $T_k:=-\tfrac{i}{2}\sigma_k$, and differentials are extended complex-linearly to $\mathfrak{sl}(2,\mathbb{C})_{\mathbb{C}}$.

\begin{theorem}[Finite-dimensional representations of $SL(2,\mathbb{C})$]\label{thm:finite-reps}
	\begin{enumerate}
		\item Every finite-dimensional continuous complex representation of $SL(2,\mathbb{C})$ is completely reducible.
		\item For $j_1,j_2\in\{0,\frac12,1,\dots\}$ there is a unique continuous representation $D^{(j_1,j_2)}$ of $SL(2,\mathbb{C})$ on $V_{j_1}\otimes V_{j_2}$ with
		\begin{equation*}
			dD^{(j_1,j_2)}(A_k)=dd_{j_1}(T_k)\otimes\mathbb{1}, \qquad dD^{(j_1,j_2)}(B_k)=\mathbb{1}\otimes dd_{j_2}(T_k).
		\end{equation*}
		It is irreducible, and every finite-dimensional continuous irreducible representation is equivalent to exactly one $D^{(j_1,j_2)}$. The representations of Definition~\ref{def:weyl-spinors} are $D^{(1/2,0)}$ and $D^{(0,1/2)}$.
		\item $D^{(j_1,j_2)}|_{SU(2)}\cong d_{j_1}\otimes d_{j_2}\cong\bigoplus_{J=|j_1-j_2|}^{j_1+j_2}d_J$.
		\item $D^{(j_1,j_2)}\circ\theta\cong D^{(j_2,j_1)}$.
	\end{enumerate}
\end{theorem}

\noindent See \cite[TODO-location]{TODO-Hal03} and \cite[TODO-location]{TODO-Hum10}. Items 3 and 4 are proved below from items 1 and 2.

\begin{proof}
	A continuous representation $D$ of the Lie group $SL(2,\mathbb{C})$ is smooth \cite[TODO-location]{TODO-Lee13}. Since $SL(2,\mathbb{C})$ is connected and simply connected by Lemma~\ref{lem:polar}, the assignment $D\mapsto dD$ is a bijection from the finite-dimensional continuous representations onto the representations of $\mathfrak{sl}(2,\mathbb{C})_{\mathbb{C}}$, and invariant subspaces for $D$ and for $dD$ coincide \cite[TODO-location]{TODO-Hal03}. By Lemma~\ref{lem:chiral}, $\mathfrak{sl}(2,\mathbb{C})_{\mathbb{C}}=\mathfrak{A}\oplus\mathfrak{B}$ with commuting ideals $\mathfrak{A}=\operatorname{span}\{A_k\}$ and $\mathfrak{B}=\operatorname{span}\{B_k\}$, each identified with $\mathfrak{su}(2)_{\mathbb{C}}$ through $A_k,B_k\mapsto T_k$.
	
	Item 1. The algebra $\mathfrak{sl}(2,\mathbb{C})_{\mathbb{C}}$ is semisimple by the proof of item 1 of Corollary~\ref{cor:poincare-structure}, and Weyl's theorem on complete reducibility applies \cite[TODO-location]{TODO-Hum10}.
	
	Item 2. The irreducible finite-dimensional representations of $\mathfrak{A}\oplus\mathfrak{B}$ are the tensor products of an irreducible representation of $\mathfrak{A}$ and one of $\mathfrak{B}$, determined up to equivalence by the pair \cite[TODO-location]{TODO-Hum10}. The irreducible representations of $\mathfrak{su}(2)_{\mathbb{C}}$ are the $dd_j$ with $j\in\{0,\frac12,1,\dots\}$ \cite[TODO-location]{TODO-Hal03}. This yields the representation of the algebra with the stated differential, its irreducibility and the classification, and the first bijection above integrates it uniquely to $SL(2,\mathbb{C})$. The identification of the Weyl representations follows from Remark~\ref{rem:weyl-differentials}.
	
	Item 3. Since $r_k=A_k+B_k$, the differential of the restriction on the real algebra $\mathfrak{su}(2)=\operatorname{span}\{r_k\}$ is
	\begin{equation*}
		dD^{(j_1,j_2)}(r_k)=dd_{j_1}(T_k)\otimes\mathbb{1}+\mathbb{1}\otimes dd_{j_2}(T_k)=d(d_{j_1}\otimes d_{j_2})(T_k),
	\end{equation*}
	with $T_k=r_k$. Two representations of the connected group $SU(2)$ with equal differentials coincide, which yields the first equivalence. The second is the Clebsch--Gordan decomposition \cite[TODO-location]{TODO-Hal03}.
	
	Item 4. By Lemma~\ref{lem:inversion}, $d\theta(A_k)=\tfrac12(r_k+ib_k)=B_k$ and $d\theta(B_k)=A_k$. Hence $D^{(j_1,j_2)}\circ\theta$ has differential $A_k\mapsto\mathbb{1}\otimes dd_{j_2}(T_k)$ and $B_k\mapsto dd_{j_1}(T_k)\otimes\mathbb{1}$ on $V_{j_1}\otimes V_{j_2}$. The flip $V_{j_1}\otimes V_{j_2}\to V_{j_2}\otimes V_{j_1}$ intertwines it with $D^{(j_2,j_1)}$, and the first bijection above concludes.
\end{proof}

The labels $(\frac12,0)$ and $(0,\frac12)$ have a concrete meaning for the spinor at rest, which is the datum on which Section 6.5 builds the constraint.

\begin{remark}[The two chiralities acting on a spinor at rest]\label{rem:chirality-at-rest}
	By item 3 of Theorem~\ref{thm:finite-reps}, $D^{(j_1,j_2)}|_{SU(2)}\cong d_{1/2}$ if and only if $(2j_1+1)(2j_2+1)=2$, \textit{i.e.}, if and only if $\{j_1,j_2\}=\{0,\tfrac12\}$. Hence among the irreducible representations only $D^{(1/2,0)}$ and $D^{(0,1/2)}$ can carry a spinor at rest, and on rotations they act identically, $D^{(1/2,0)}(u)\xi=D^{(0,1/2)}(u)\xi=u\xi$ for $\xi\in\mathbb{C}^2$. They are inequivalent by the uniqueness of the label in item 2, and they differ only on boosts. The generator $b_k$ acts as $+\tfrac12\sigma_k$ in the first and as $-\tfrac12\sigma_k$ in the second. Transporting $\xi=(\xi_+,\xi_-)$ to the momentum $p=m(\cosh w,0,0,\sinh w)$ with the standard boost $A_p=\exp(\tfrac12w\,\sigma_3)$ of Proposition~\ref{prop:orbit} yields
	\begin{equation*}
		D^{(1/2,0)}(A_p)\,\xi=\big(e^{w/2}\xi_+,\,e^{-w/2}\xi_-\big),\qquad D^{(0,1/2)}(A_p)\,\xi=\big(e^{-w/2}\xi_+,\,e^{w/2}\xi_-\big).
	\end{equation*}
	The component of the spin along $\mathbf{p}$ is stretched in the first representation and shrunk in the second, and the roles are exchanged for the component opposite to $\mathbf{p}$. This is the sense of the labels. A representation of chirality $(\frac12,0)$ correlates the amplification under a boost with the spin direction in one way, and $(0,\frac12)$ in the opposite way. Space inversion reverses the momentum and leaves the spin, an axial vector, unchanged, which by Equation~\eqref{eq:theta-boost} replaces $A_p$ with $A_p^{-1}$. It therefore exchanges the two representations, and this is the mechanism behind the parity obstruction of Section 6.5.
\end{remark}

Remark~\ref{rem:chirality-at-rest} shows that boosts act on the Weyl representations through positive Hermitian matrices other than $\mathbb{1}$, so that the Euclidean norm of $\mathbb{C}^2$ is not preserved. The following proposition shows that this is not a defect of the two representations considered, but a property of every finite-dimensional representation of $SL(2,\mathbb{C})$.

\begin{proposition}[No finite-dimensional unitary representations]\label{prop:no-unitary}
	Every continuous homomorphism $\varrho\colon SL(2,\mathbb{C})\to U(n)$ is trivial.
\end{proposition}

\begin{proof}
	The homomorphism $\varrho$ is smooth \cite[TODO-location]{TODO-Lee13}, and its differential takes values in the skew-Hermitian matrices. Set $X_k:=d\varrho(A_k)$ and $Y_k:=d\varrho(B_k)$, with $d\varrho$ extended complex-linearly. Then $d\varrho(r_k)=X_k+Y_k$ is skew-Hermitian, and $d\varrho(b_k)=i(X_k-Y_k)$ is skew-Hermitian, \textit{i.e.}, $X_k-Y_k$ is Hermitian. Adding the two conditions yields $X_k^\dagger=-Y_k$. By Lemma~\ref{lem:chiral}, $[X_j,Y_k]=0$, that is $[X_j,X_k^\dagger]=0$ for all $j,k$, and $X_3=[X_1,X_2]$. Using $X_2X_3^\dagger=X_3^\dagger X_2$ and the cyclicity of the trace,
	\begin{equation*}
		\operatorname{tr}(X_3X_3^\dagger)=\operatorname{tr}(X_1X_2X_3^\dagger)-\operatorname{tr}(X_2X_1X_3^\dagger)=\operatorname{tr}(X_1X_3^\dagger X_2)-\operatorname{tr}(X_2X_1X_3^\dagger)=0.
	\end{equation*}
	Hence $X_3=0$, and by cyclic permutation of the indices $X_1=X_2=0$. It descends that $Y_k=-X_k^\dagger=0$ and $d\varrho=0$. Since $SL(2,\mathbb{C})$ is connected by Lemma~\ref{lem:polar}, it is generated by $\exp(\mathfrak{sl}(2,\mathbb{C}))$, and $\varrho(\exp\xi)=\exp(d\varrho(\xi))=\mathbb{1}$ yields $\varrho=\mathbb{1}$.
\end{proof}

\begin{remark}[Independence of the inner product]\label{rem:non-unitary-boosts}
	The proposition applies to every inner product on the representation space, since the choice of an orthonormal basis turns $D$ into a homomorphism into $U(n)$. Hence a nontrivial finite-dimensional representation $D$ of $SL(2,\mathbb{C})$ preserves no positive definite Hermitian form. The obstruction lies in the boosts, because $SU(2)$ is compact and each $d_j$ is unitary. For $D^{(1/2,0)}$ it is visible directly, as $D^{(1/2,0)}(A_p)^\dagger D^{(1/2,0)}(A_p)=A_p^2=\hat{p}/m\neq\mathbb{1}$ for $p\neq\mathring{p}$ by Proposition~\ref{prop:orbit}. In Section 6.5 this fact fixes the logic. The fixed matrices $D(g)$ of a local realization cannot be unitary, so that the Euclidean norm of the space $W$ cannot be the norm of the representation, and the invariant form has to be found elsewhere.
\end{remark}

\section{Poincar\'e Invariance of the Projective State Space}

The preceding section fixed the group $\mathcal{P}$ and its Lie algebra. This section brings the group to bear on the states. The pure states are the points of the projective Hilbert space $\mathbb{P}\mathcal{H}$ of Chapter~3, carrying the Kähler structure of Equation~(3.24), and the question is how much of the unitary formalism is forced by requiring that $\mathcal{P}$ acts on these states by symmetries. The data are a separable Hilbert space $\mathcal{H}$, which enters only through the Kähler manifold $\mathbb{P}\mathcal{H}$, and a group homomorphism from $\mathcal{P}$ to the Kähler isomorphisms of $\mathbb{P}\mathcal{H}$. No momentum, no position, no spin and no Hamiltonian is assumed. Three results are obtained. First, the action lifts in exactly one way to a strongly continuous unitary representation of $\mathcal{P}$, because $H^2(\mathfrak{p},\mathbb{R})=0$ and $\mathfrak{p}$ is perfect. Second, the generator of the time translations is a selfadjoint operator $H$ with no additive ambiguity, and irreducibility together with the positivity of $H$ singles out the elementary relativistic systems. Third, a realization of the pair, \textit{i.e.}, a Kähler isomorphism onto another projective space intertwining the two actions, is the same thing as a unitary equivalence of the lifts with no phase left over. This last statement is what allows Sections 6.4 to 6.7 to replace $\mathbb{P}\mathcal{H}$ by explicit models without leaving the space of states. The next step of the chain, Section~6.3, extracts from the lift the spectral measure of the translations, which is the four-momentum.

\subsection{The Poincar\'e Action and its Unitary Lift}

The first step is to state what a symmetry acting on states is and to compare it with the unitary operators of the usual formalism. The data are $\mathbb{P}\mathcal{H}$ and a homomorphism from $\mathcal{P}$ to its Kähler isomorphisms. What is obtained is a unitary representation of $\mathcal{P}$ on $\mathcal{H}$ whose projectivization is the given action, and this representation is unique. In the Galilean case of Chapter~\ref{ch:analytical-hamiltonians} the lift required an enlargement of the group and was unique only up to a character. The vanishing of $H^2(\mathfrak{p},\mathbb{R})$ in Theorem~\ref{thm:poincare-cohomology} removes both difficulties, and the uniqueness is used in the last subsection to turn Kähler isomorphisms into unitary equivalences.

\begin{definition}[Poincar\'e action]\label{def:poincare-action}
	Let $\mathcal{H}$ be a separable complex Hilbert space as per Definition~2.2.6, let $\mathbb{P}\mathcal{H}$ be its projective space as per Definition~3.2.2 with the Kähler structure $(J,g,\omega)$ of Equation~(3.24), and let $\mathcal{P}$ be as per Definition~\ref{def:poincare-group}. A \emph{Poincar\'e action} on $\mathbb{P}\mathcal{H}$ is a group homomorphism $\Phi$ from $\mathcal{P}$ to the group of Kähler isomorphisms of $\mathbb{P}\mathcal{H}$, written $\kappa\mapsto\Phi_\kappa$, such that for every $\phi\in\mathbb{P}\mathcal{H}$ the orbit map $\kappa\mapsto\Phi_\kappa(\phi)$ is continuous for the manifold topology of Proposition~3.2.7.
\end{definition}

\begin{remark}[Kähler isomorphisms and continuity]
	A Kähler isomorphism between Kähler manifolds $(M,J,g)$ and $(M',J',g')$, as per Definition~3.1.10, is a diffeomorphism $F$ with $F^*g'=g$ and $dF\circ J=J'\circ dF$. It preserves the fundamental form $\omega=g(J\cdot,\cdot)$ of Definition~3.1.8, hence the whole triple. By Proposition~3.2.9 the manifold topology of $\mathbb{P}\mathcal{H}$ is the weak$^*$ topology, so that continuity of the orbit maps means that expectation values of all observables vary continuously along the orbits, as in the remark following Definition~\ref{def:galilean-action}. Since $\mathcal{P}$ is connected and simply connected by Corollary~\ref{cor:poincare-topology}, the hypotheses of Bargmann's correspondence of Proposition~\ref{prop:bargmann-correspondence} are met.
\end{remark}

A system that cannot be split into two independent sectors by the symmetry group is the object of study of the chapter. This property is stated on the projective space and uses only the Kähler action.

\begin{definition}[Irreducible Poincar\'e action]\label{def:poincare-irreducible}
	A Poincar\'e action $\Phi$ as per Definition~\ref{def:poincare-action} is \emph{irreducible} if the only closed complex subspaces $\mathcal{L}\subseteq\mathcal{H}$ with $\Phi_\kappa(\mathbb{P}\mathcal{L})\subseteq\mathbb{P}\mathcal{L}$ for all $\kappa\in\mathcal{P}$ are $\{0\}$ and $\mathcal{H}$.
\end{definition}

The action $\Phi$ is defined on rays, whereas energy and momentum are operators. Passing from the first to the second requires choosing a unitary operator for each ray transformation, and the phases of this choice may obstruct a genuine representation, as in Chapter~\ref{ch:analytical-hamiltonians}. For $\mathcal{P}$ the obstruction vanishes, and the statement below is the form of Bargmann's theorem needed in the chapter.

\begin{corollary}[Lift of continuous ray representations of $\mathcal{P}$]\label{cor:poincare-lift}
	Let $\mathcal{H}$ be a separable Hilbert space and let $\kappa\mapsto V_\kappa$ be a continuous homomorphism from $\mathcal{P}$ to the group $U(\mathcal{H})/U(1)$ of unitary ray transformations, endowed with the topology in which $V_n\to V$ if and only if $|(x,V_ny)|\to|(x,Vy)|$ for all $x,y\in\mathcal{H}$. Then there exists a strongly continuous unitary representation $\Pi:\mathcal{P}\to U(\mathcal{H})$ such that $\Pi(\kappa)$ belongs to the ray $V_\kappa$ for all $\kappa\in\mathcal{P}$.
\end{corollary}

\begin{proof}
	The group $\mathcal{P}$ is connected and simply connected by Corollary~\ref{cor:poincare-topology}. By items 1 and 2 of Proposition~\ref{prop:bargmann-correspondence}, the class of the multiplier of $V$ corresponds to an element of $H^2(\mathfrak{p},\mathbb{R})$, which vanishes by Theorem~\ref{thm:poincare-cohomology}. The multiplier is therefore equivalent to the trivial one. Item 3 of the same proposition, applied to the continuous real cocycle $\zeta=0$, shows that the phases of $V$ can be chosen so that $V$ is strongly continuous with multiplier exactly $e^{i0}=1$. It descends that $\Pi(\kappa):=V_\kappa$ is a strongly continuous unitary representation of $\mathcal{P}$ with $\Pi(\kappa)$ in the original ray.
\end{proof}

\noindent The Galilean counterpart, Theorem~\ref{thm:bargmann}, yields in addition a mass parameter and a representation of the extended group, while here the lift is a representation of $\mathcal{P}$ itself.

The corollary concerns ray representations, whereas hypothesis H1 provides a Kähler action. The next statement connects the two and records that the lift is unique, a property without Galilean counterpart.

\begin{proposition}[Unitary lift of a Poincar\'e action]\label{prop:poincare-lift}
	Let $\Phi$ be a Poincar\'e action on $\mathbb{P}\mathcal{H}$. There exists a unique strongly continuous unitary representation $\Pi:\mathcal{P}\to U(\mathcal{H})$ such that
	\begin{equation}\label{eq:lift-def}
		\Phi_\kappa([x])=[\Pi(\kappa)x],\qquad \kappa\in\mathcal{P},\ [x]\in\mathbb{P}\mathcal{H}.
	\end{equation}
\end{proposition}

\begin{proof}
	\emph{Existence.} By item 1 of Proposition~3.3.6, each $\Phi_\kappa$ is induced by a unitary operator $V_\kappa$, unique up to a phase. Since $V_\kappa V_{\kappa'}$ and $V_{\kappa\kappa'}$ both induce $\Phi_{\kappa\kappa'}$, the map $\kappa\mapsto[V_\kappa]$ is a homomorphism into $U(\mathcal{H})/U(1)$. For $x,y\in\mathcal{H}\setminus\{0\}$ the quantity $|(x,V_\kappa y)|^2/(\|x\|^2\|y\|^2)$ is the transition probability between $[x]$ and $\Phi_\kappa([y])$. By Theorem~3.8 this probability is a continuous function of the Riemannian distance, and the distance induces the manifold topology by Remark~3.4.11. Together with the continuity of the orbit maps, it follows that $\kappa\mapsto|(x,V_\kappa y)|$ is continuous. Corollary~\ref{cor:poincare-lift} yields a strongly continuous unitary representation $\Pi$ with $\Pi(\kappa)$ in the ray $V_\kappa$, which is Equation~\eqref{eq:lift-def}.
	
	\emph{Uniqueness.} Let $\Pi'$ be another lift. For every $\kappa$ both $\Pi(\kappa)$ and $\Pi'(\kappa)$ induce $\Phi_\kappa$, hence $\Pi'(\kappa)=\chi(\kappa)\Pi(\kappa)$ with $\chi(\kappa)\in U(1)$. For a unit vector $x$ one has $\chi(\kappa)=(\Pi(\kappa)x,\Pi'(\kappa)x)$, which is continuous by strong continuity. Expanding $\Pi'(\kappa\kappa')$ in two ways shows that $\chi$ is a homomorphism. A continuous homomorphism of Lie groups is smooth \cite[TODO-location]{TODO-Lee13}, and its differential $d\chi:\mathfrak{p}\to i\mathbb{R}$ is a Lie algebra homomorphism into an abelian algebra. It vanishes on $[\mathfrak{p},\mathfrak{p}]=\mathfrak{p}$ by item 2 of Corollary~\ref{cor:poincare-structure}, so that $\chi(\exp\xi)=1$ for every $\xi\in\mathfrak{p}$. Since $\mathcal{P}$ is connected it is generated by $\exp(\mathfrak{p})$, and $\chi=1$.
\end{proof}

\begin{remark}[Irreducibility of $\Pi$ and comparison with Chapter~\ref{ch:analytical-hamiltonians}]
	For a closed subspace $\mathcal{L}\subseteq\mathcal{H}$ one has $\Phi_\kappa(\mathbb{P}\mathcal{L})\subseteq\mathbb{P}\mathcal{L}$ if and only if $\Pi(\kappa)\mathcal{L}\subseteq\mathcal{L}$, because $[\Pi(\kappa)x]\in\mathbb{P}\mathcal{L}$ means $\Pi(\kappa)x\in\mathcal{L}$. Hence $\Phi$ is irreducible as per Definition~\ref{def:poincare-irreducible} if and only if the only closed subspaces invariant under $\Pi(\mathcal{P})$ are $\{0\}$ and $\mathcal{H}$. Proposition~\ref{prop:lift} leaves the Galilean lift undetermined up to the character $e^{i\epsilon s}$, and this freedom is the additive constant of the Hamiltonian. Here $\mathcal{P}$ has no nonzero character, and the lift is determined by $\Phi$ alone.
\end{remark}

\subsection{Energy and Elementary Relativistic Systems}

A state of a relativistic system has an energy, and the energy operator must be a consequence of the action rather than an input. The data are the lift $\Pi$ of Proposition~\ref{prop:poincare-lift}. What is obtained is the energy operator, defined as a Stone generator, and three conditions on $\Phi$ that drive the chapter, namely continuity of the orbits, irreducibility and positivity of the energy. Each of them is used at a definite place in the following sections, and the table at the end of the subsection records where. The expression of $H$ through the spectral measure of the translations is not available at this stage and is the subject of Section~6.3.

For $a\in\mathbb{R}^{1,3}$ we write $\tau_a:=(a,\mathbb{1})\in\mathcal{P}$. By item 1 of Lemma~\ref{lem:poincare-relations} the map $s\mapsto\tau_{(s,\mathbf{0})}$ is a one-parameter subgroup of $\mathcal{P}$. Composing it with $\Pi$ yields a strongly continuous one-parameter unitary group, and Stone's theorem, Theorem~2.10, associates with it a unique selfadjoint operator.

\begin{definition}[Energy operator]\label{def:energy-operator}
	Let $\Phi$ be a Poincar\'e action on $\mathbb{P}\mathcal{H}$ with lift $\Pi$ as per Proposition~\ref{prop:poincare-lift}. The \emph{energy operator} of $\Phi$ is the selfadjoint operator $H$ on $\mathcal{H}$ such that
	\begin{equation}\label{eq:energy-stone}
		\Pi\big(\tau_{(s,\mathbf{0})}\big)=e^{isH/\hbar},\qquad s\in\mathbb{R}.
	\end{equation}
\end{definition}

\begin{remark}[$H$ is not assigned]
	The Kähler structure of $\mathbb{P}\mathcal{H}$ does not distinguish any operator, and by Proposition~3.3.6 a one-parameter group of Kähler isomorphisms determines its generator only up to an additive constant. In Definition~\ref{def:energy-operator} the operator is not chosen but read off from the group, and the additive constant is fixed because $\Pi$ is unique by Proposition~\ref{prop:poincare-lift}. In Chapter~\ref{ch:analytical-hamiltonians} the same generator was determined only up to $\hbar\epsilon$. Equation~\eqref{eq:energy-stone} follows the convention of Definition~\ref{def:generators}, so that in the active convention of Remark~\ref{rem:active-convention} the Schr\"odinger flow is $t\mapsto\Phi_{\tau_{(-t,\mathbf{0})}}$, \textit{i.e.}, $[x]\mapsto[e^{-itH/\hbar}x]$. No relation between $H$ and the momentum is assumed. That $H$ is a function of the momentum is derived in Theorem~\ref{thm:relativistic-hamiltonian}, from the support of the spectral measure of Proposition~\ref{prop:snag-translations}.
\end{remark}

The existence of a ground level is a property of $H$ and not of the group alone, so that it is stated as a separate condition. It is the relativistic counterpart of the choice of the sign of $m$ in Theorem~\ref{thm:free-hamiltonian}.

\begin{definition}[Positive energy]\label{def:positive-energy}
	A Poincar\'e action $\Phi$ has \emph{positive energy} if its energy operator $H$, as per Definition~\ref{def:energy-operator}, is bounded below, \textit{i.e.}, if there exists $c\in\mathbb{R}$ with $(x,Hx)\geq c\|x\|^2$ for every $x\in D(H)$.
\end{definition}

\begin{remark}[The lower bound is an output]
	Boundedness from below is invariant under additive constants, so that in the Galilean setting it would not depend on the lift. Here it refers to a canonical operator, and the optimal constant $c$ is not assumed. It will turn out to be the mass, since $\sigma(H)=[m,\infty)$ by Theorem~\ref{thm:relativistic-hamiltonian}.
\end{remark}

Irreducibility and positivity of the energy are the only conditions imposed on a Poincar\'e action beyond H1. Together they define the object studied in the rest of the chapter.

\begin{definition}[Elementary relativistic system]\label{def:elementary-system}
	An \emph{elementary relativistic system} is an irreducible Poincar\'e action $\Phi$, as per Definition~\ref{def:poincare-irreducible}, with positive energy as per Definition~\ref{def:positive-energy}.
\end{definition}

The mass does not enter the definition. It is the value of a Casimir operator, which exists for every elementary relativistic system by Proposition~\ref{prop:mass-shell}, and the distinction between the massive and the massless case is made on that value.

\begin{definition}[Massive elementary relativistic system]\label{def:massive-system}
	Let $\Phi$ be an elementary relativistic system and let $m^2\in\mathbb{R}$ be the value of $P^2$ as per Proposition~\ref{prop:mass-shell}, with $m\geq0$. The system $\Phi$ is \emph{massive} if $m>0$, and $m$ is then called its \emph{mass}.
\end{definition}

The table below records the role of each hypothesis. It separates what is assumed from what is derived and shows that the mass is in the second group.

\begin{table}[ht]
	\centering
	\begin{tabular}{p{3.1cm}p{4.6cm}p{6.1cm}}
		\hline
		\textbf{Hypothesis} & \textbf{Where it enters} & \textbf{What it yields} \\
		\hline
		H1, Poincar\'e action by Kähler isomorphisms with continuous orbits
		& Proposition~\ref{prop:poincare-lift}, with Theorem~\ref{thm:poincare-cohomology}
		& A unique strongly continuous unitary lift $\Pi$, hence an energy operator $H$ without additive constant. \\
		H2, irreducibility
		& Proposition~\ref{prop:mass-shell} through Lemma~\ref{lem:schur-infinite}, and Theorem~\ref{thm:wigner-mackey}
		& The Casimir operator $P^2$ is a number $m^2\mathbb{1}$. The multiplicity space of the spectral measure carries one irreducible representation of $SU(2)$, whose label is the spin $j$. \\
		H3, positive energy
		& Proposition~\ref{prop:mass-shell} and Theorem~\ref{thm:relativistic-hamiltonian}
		& The cases $p^2<0$ and $\mathcal{O}_m^-$ are excluded, $\operatorname{supp}E=\mathcal{O}_m^+$ and $H=\omega(\mathbf{P})$ with $\sigma(H)=[m,\infty)$. \\
		$m>0$, not a hypothesis
		& Definition~\ref{def:massive-system}, Proposition~\ref{prop:orbit} and Theorem~\ref{thm:main}
		& The orbit is $\mathcal{O}_m^+$ with stabilizer $SU(2)$. The case $m=0$ is not treated. \\
		\hline
	\end{tabular}
	\caption{Role of the hypotheses H1 to H3 of Definition~\ref{def:hypotheses}.}
	\label{tab:hypotheses-roles}
\end{table}

The hypotheses H4 to H6 do not appear in the table because they concern the description of the states by fields and not the states themselves. They enter in Section~6.5, and the parity of the state space is never assumed.

\subsection{Realizations of the State Space}

The lift $\Pi$ acts on a Hilbert space that has not been specified, and nothing so far says in which sense two Hilbert spaces carrying the same action describe the same states. The data are two projective spaces with Poincar\'e actions. What is obtained is the notion of realization of $(\mathbb{P}\mathcal{H},\Phi)$ and its characterization by unitary operators intertwining the lifts exactly. The following sections replace $\mathbb{P}\mathcal{H}$ by explicit models, and the characterization shows that this replacement is always a Kähler isomorphism commuting with the action, so that the space of states is not changed. The remainder of the chapter uses only the last remark, which lists the three models.

\begin{definition}[Realization]\label{def:realization}
	Let $\Phi$ be a Poincar\'e action on $\mathbb{P}\mathcal{H}$. A \emph{realization} of $(\mathbb{P}\mathcal{H},\Phi)$ is a triple $(\mathcal{K},F,\Phi')$, where $\mathcal{K}$ is a separable complex Hilbert space, $F:\mathbb{P}\mathcal{H}\to\mathbb{P}\mathcal{K}$ is a Kähler isomorphism and $\Phi'$ is a Poincar\'e action on $\mathbb{P}\mathcal{K}$, such that
	\begin{equation}\label{eq:realization-intertwining}
		F\circ\Phi_\kappa=\Phi'_\kappa\circ F,\qquad \kappa\in\mathcal{P}.
	\end{equation}
\end{definition}

\begin{remark}[Realizations are determined by $F$]
	Equation~\eqref{eq:realization-intertwining} forces $\Phi'_\kappa=F\circ\Phi_\kappa\circ F^{-1}$, which is a homomorphism into the Kähler isomorphisms of $\mathbb{P}\mathcal{K}$ with continuous orbit maps, because $F$ is a diffeomorphism. A realization is therefore the same thing as a Kähler isomorphism $F$ onto some $\mathbb{P}\mathcal{K}$. Two realizations $F$ and $F'$ of the same pair differ by the Kähler isomorphism $F'\circ F^{-1}$, which intertwines the two induced actions. Distinctions between realizations lie in additional structure placed on $\mathbb{P}\mathcal{K}$, such as the fields with fixed matrices of Section~6.5, and not in the state space.
\end{remark}

Kähler isomorphisms between two projective spaces are needed in a form more general than Proposition~3.3.6, which concerns a single space. The statement below identifies them with unitary operators, and its proof makes explicit the exclusion of antiunitary operators.

\begin{lemma}[Kähler isomorphisms between two projective spaces]\label{lem:kahler-iso-two-spaces}
	Let $\mathcal{H}$ and $\mathcal{K}$ be separable complex Hilbert spaces.
	\begin{enumerate}
		\item If $W:\mathcal{H}\to\mathcal{K}$ is unitary, then $[W]:\mathbb{P}\mathcal{H}\to\mathbb{P}\mathcal{K}$, $[x]\mapsto[Wx]$, is a Kähler isomorphism.
		\item Every Kähler isomorphism $F:\mathbb{P}\mathcal{H}\to\mathbb{P}\mathcal{K}$ equals $[W]$ for a unitary $W:\mathcal{H}\to\mathcal{K}$, and $W$ is unique up to a factor in $U(1)$.
	\end{enumerate}
\end{lemma}

\begin{proof}
	\emph{Item 1.} Let $x\in S\mathcal{H}$ and $\psi=[x]$. For $\Phi\in\mathcal{H}$ with $(x,\Phi)\neq0$, the chart of Definition~3.2.5 satisfies
	\begin{equation*}
		F_{[Wx]}\big([W\Phi]\big)=\frac{W\Phi}{(Wx,W\Phi)}-Wx=W F_\psi([\Phi]),
	\end{equation*}
	because $(Wx,W\Phi)=(x,\Phi)$. In the charts of Proposition~3.2.7 the map $[W]$ is therefore the restriction $W|_{x^\perp}:x^\perp\to(Wx)^\perp$, a $\mathbb{C}$-linear isometric bijection. It follows that $[W]$ is a bijection with inverse $[W^{-1}]$, smooth and holomorphic, with differential $W|_{x^\perp}$ at $\psi$. Since $J$ is multiplication by $i$ and $g=2\hbar\,\mathrm{Re}(\cdot,\cdot)$ on $x^\perp$ by Equation~(3.26), $[W]$ preserves $J$ and $g$.
	
	\emph{Item 2.} If $\dim\mathcal{H}=1$, then $\mathbb{P}\mathcal{H}$ is a point, $\dim\mathcal{K}=1$ and the claim is trivial. Let $\dim\mathcal{H}\geq2$. A Kähler isomorphism preserves $g$, hence the length of piecewise $C^1$ curves of Definition~3.4.9, hence the Riemannian distance $d_g$ of Definition~3.4.10, since $F$ is a bijection and $F^{-1}$ preserves $g$ as well. By Theorem~3.8 the transition probability $P([x],[y])=\cos^2\!\big(d_g([x],[y])/\sqrt{2\hbar}\big)$ is preserved by $F$. Wigner's theorem \cite[TODO-location]{TODO-Mor18} yields a unitary or antiunitary operator $W:\mathcal{H}\to\mathcal{K}$ with $F=[W]$, unique up to a phase. If $W$ were antiunitary, the same computation as in item 1 would give $F_{[Wx]}\circ F\circ F_\psi^{-1}=W|_{x^\perp}$, since $W(c\Phi)=\bar{c}\,W\Phi$ and $\overline{(x,\Phi)}=(Wx,W\Phi)$. This map is antilinear, so that $dF\circ J=-J\circ dF$ on $x^\perp\neq\{0\}$, which contradicts the requirement that $F$ preserve $J$. Hence $W$ is unitary. If $W'$ is another unitary with $[W']=[W]$, then $U:=W^{-1}W'$ satisfies $Ux\in\mathbb{C}x$ for every $x$. For linearly independent $x,y$ one finds $U(x+y)=c(x+y)$ together with $Ux=ax$ and $Uy=by$, hence $a=b=c$. Therefore $U$ is a scalar of modulus one.
\end{proof}

\noindent Item 2 is Proposition~3.3.6 with $\mathcal{K}=\mathcal{H}$, with the exclusion of antiunitary operators obtained from the preservation of $J$ and not from continuity of a path.

The lemma identifies Kähler isomorphisms with unitary operators up to a phase. Combined with the uniqueness of the lift in Proposition~\ref{prop:poincare-lift}, it turns the intertwining of Definition~\ref{def:realization} into an exact unitary equivalence.

\begin{proposition}[Realizations are unitary equivalences]\label{prop:realization-equivalence}
	Let $\Phi$ be a Poincar\'e action on $\mathbb{P}\mathcal{H}$ with lift $\Pi$, and let $(\mathcal{K},F,\Phi')$ be a realization of $(\mathbb{P}\mathcal{H},\Phi)$ with lift $\Pi'$ of $\Phi'$. Then there exists a unitary operator $W:\mathcal{H}\to\mathcal{K}$, unique up to a phase, with $F=[W]$, and for every such $W$
	\begin{equation}\label{eq:exact-equivalence}
		W\,\Pi(\kappa)\,W^{-1}=\Pi'(\kappa),\qquad \kappa\in\mathcal{P}.
	\end{equation}
	Moreover $W H W^{-1}=H'$ with $D(H')=W D(H)$, where $H'$ is the energy operator of $\Phi'$, the action $\Phi$ is irreducible if and only if $\Phi'$ is, and $\Phi$ has positive energy if and only if $\Phi'$ has. Conversely, if $W:\mathcal{H}\to\mathcal{K}$ is unitary and $\Pi'$ is a strongly continuous unitary representation of $\mathcal{P}$ on $\mathcal{K}$ satisfying Equation~\eqref{eq:exact-equivalence}, then $\big(\mathcal{K},[W],[\Pi'(\cdot)\,\cdot\,]\big)$ is a realization of $(\mathbb{P}\mathcal{H},\Phi)$.
\end{proposition}

\begin{proof}
	The operator $W$ exists and is unique up to a phase by item 2 of Lemma~\ref{lem:kahler-iso-two-spaces}. The operators $\Pi''(\kappa):=W\Pi(\kappa)W^{-1}$ form a strongly continuous unitary representation of $\mathcal{P}$ on $\mathcal{K}$. By Equations~\eqref{eq:lift-def} and~\eqref{eq:realization-intertwining},
	\begin{equation*}
		[\Pi''(\kappa)Wx]=[W\Pi(\kappa)x]=F\big(\Phi_\kappa([x])\big)=\Phi'_\kappa\big([Wx]\big),
	\end{equation*}
	so that $\Pi''$ is a lift of $\Phi'$. Uniqueness in Proposition~\ref{prop:poincare-lift} yields $\Pi''=\Pi'$, which is Equation~\eqref{eq:exact-equivalence}. Evaluating it on $\tau_{(s,\mathbf{0})}$ gives $e^{isH'/\hbar}=e^{isWHW^{-1}/\hbar}$, and uniqueness in Stone's theorem yields $H'=WHW^{-1}$. A closed subspace $\mathcal{L}\subseteq\mathcal{K}$ is invariant under $\Pi'(\mathcal{P})$ if and only if $W^{-1}\mathcal{L}$ is invariant under $\Pi(\mathcal{P})$, which gives the statement on irreducibility by the remark following Proposition~\ref{prop:poincare-lift}. The operators $H$ and $H'$ are unitarily equivalent, so they have the same lower bound. For the converse, $[W]$ is a Kähler isomorphism by item 1 of Lemma~\ref{lem:kahler-iso-two-spaces}. The map $\Phi'_\kappa:=[\Pi'(\kappa)\,\cdot\,]$ is a homomorphism into the Kähler isomorphisms by the same item, and its orbit maps are continuous because $\Pi'$ is strongly continuous and $x\mapsto[x]$ is continuous for the manifold topology. Finally $[W][\Pi(\kappa)x]=[\Pi'(\kappa)Wx]$ by Equation~\eqref{eq:exact-equivalence}.
\end{proof}

\begin{remark}[Exactness of the equivalence]
	Equation~\eqref{eq:exact-equivalence} holds with no phase. In the Galilean case Proposition~\ref{prop:lift} would allow an intertwiner only up to a character $e^{i\epsilon s}$, which shifts the Hamiltonian by a constant. For $\mathcal{P}$ this freedom is absent because $H^2(\mathfrak{p},\mathbb{R})=0$ and $\mathfrak{p}$ is perfect. It descends that Theorem~\ref{thm:main} is a statement on unitary equivalence of the lifts, and that the energy operator of a realization is the transported one, with no additive constant.
\end{remark}

The notion of realization shows what the state space does and does not determine. Two separable infinite-dimensional spaces cannot be distinguished by their Kähler geometry, and the spin must therefore be sought in the action.

\begin{remark}[The projective space does not carry the spin]\label{rem:ph-no-spin}
	Let $\mathcal{H}_1$ and $\mathcal{H}_2$ be separable infinite-dimensional Hilbert spaces. They are unitarily equivalent, since an orthonormal basis of the first is carried onto an orthonormal basis of the second \cite[TODO-location]{TODO-Mor18}, and by item 1 of Lemma~\ref{lem:kahler-iso-two-spaces} a unitary operator $W:\mathcal{H}_1\to\mathcal{H}_2$ induces a Kähler isomorphism $[W]$. It descends that $\mathbb{P}\mathcal{H}_1$ and $\mathbb{P}\mathcal{H}_2$ are isomorphic Kähler manifolds. The space $L^2(\mathcal{O}_m^+,\mu;\mathbb{C}^{2j+1})$ of Definition~\ref{def:induced-rep} is separable and infinite-dimensional for every $j$, hence the same Kähler manifold hosts the states of particles of any spin. The value of $j$ cannot be read from $\mathbb{P}\mathcal{H}$. It is a property of the action $\Phi$, and Theorem~\ref{thm:wigner-mackey} shows that it is the label of an irreducible representation of the stabilizer of the rest momentum.
\end{remark}

The chapter works with three descriptions of the same pair $(\mathbb{P}\mathcal{H},\Phi)$. They are linked by realizations, and keeping them apart avoids a frequent confusion between states and momenta.

\begin{remark}[The three models]\label{rem:three-models}
	In each of the following models a state is a ray.
	\begin{itemize}
		\item[M1.] The abstract space $\mathbb{P}\mathcal{H}$ with $\Phi_\kappa=[\Pi(\kappa)\,\cdot\,]$ as per Proposition~\ref{prop:poincare-lift}.
		\item[M2.] The space $\mathbb{P}L^2(\mathcal{O}_m^+,\mu;\mathbb{C}^{2j+1})$ with $[U_{m,j}(\kappa)\,\cdot\,]$. It is reached from M1 by $[W]$, with $W$ the unitary operator of Theorem~\ref{thm:wigner-mackey}, and this is a realization by Proposition~\ref{prop:realization-equivalence}.
		\item[M3.] The space $\mathbb{P}\mathcal{F}_D^2$ of Dirac solutions with the norm $\langle\cdot,\cdot\rangle_D$ and the action $[\mathcal{U}_D(\kappa)\,\cdot\,]$. It is reached from M2 by $[J]$, with $J$ the unitary operator of Theorem~\ref{thm:covariant-norm}.
	\end{itemize}
	The set $\mathcal{O}_m^+$ is a set of states in none of them. It is the support of the spectral measure of the translations, \textit{i.e.}, the set of possible values of the four-momentum, and a state in M2 or M3 is a ray of square integrable functions on $\mathcal{O}_m^+$, not a point of $\mathcal{O}_m^+$. Replacing one model by another changes the description and not the space of states. The composition of two Kähler isomorphisms is again a Kähler isomorphism, so that the passage from M1 to M3 is a single realization of $(\mathbb{P}\mathcal{H},\Phi)$. The Kähler structure, and with it the Hamiltonian flow of $\langle H\rangle$ and the transition probabilities, is transported unchanged, and Section~6.7 verifies this for M3.
\end{remark}

\section{The Four-Momentum and the Mass Shell}
\label{sec:poincare-spectral}

The previous section turned an irreducible Poincar\'e action $\Phi$ on $\mathbb{P}\mathcal{H}$ into a unique strongly continuous unitary representation $\Pi$ of $\mathcal{P}$. The question of this section, expressed in terms of states, is which values of the four-momentum the states of the system can take, and how the Lorentz group relates these values to each other. The data are the lift $\Pi$ of Proposition~\ref{prop:poincare-lift}, the irreducibility of $\Phi$ and the positivity of the energy, \textit{i.e.}, hypotheses H1 to H3 of Definition~\ref{def:hypotheses}. The translations form an abelian group, so that the restriction of $\Pi$ to them decomposes over the characters of $\mathbb{R}^4$. A character is labelled by a vector $p\in M^*$, where $M^*$ is the space of momenta, identified with $M$ through $\eta$. The variable $p$ is not a point of $\mathbb{P}\mathcal{H}$, which is a manifold of rays with no momentum coordinate, and it is not a point of spacetime $M$. It is the label of a character of the translation group, \textit{i.e.}, a possible value of the four-momentum. The spectral measure $E$ of the translations records how a ray is distributed over these values, and its support is the set of all values that occur. The outcome of the section is that this support is a single sheet $\mathcal{O}_m^+$ of a mass hyperboloid, with $m$ the eigenvalue of a Casimir operator. The set $\mathcal{O}_m^+$ is therefore a set of values of the four-momentum and not a set of states. This is the first step of the chain. What it leaves open, the multiplicity of $E$ above each point of $\mathcal{O}_m^+$ and the action of the Lorentz group on it, is what the next step has to determine.

The restriction of $\Pi$ to the translations is the only part of the representation that is accessible without information on the Lorentz group. Since $(a,\mathbb{1})(a',\mathbb{1})=(a+a',\mathbb{1})$ by Definition~\ref{def:poincare-group}, the map $a\mapsto\Pi((a,\mathbb{1}))$ is a strongly continuous unitary representation of the additive group $M$. Its restrictions to the time axis and to each spatial axis are therefore strongly continuous one-parameter unitary groups, to which Stone's theorem applies.

\begin{definition}[Generators of the translations]\label{def:poincare-generators}
	Let $\Pi$ be as per Proposition~\ref{prop:poincare-lift}. For $a\in M$ and $s\in\mathbb{R}$ we set
	\begin{equation*}
		T(a):=\Pi\big((a,\mathbb{1})\big), \qquad S(s):=T\big((s,\mathbf{0})\big).
	\end{equation*}
	The selfadjoint operators $H$ and $P_k$, with $k\in\{1,2,3\}$, are the generators, in the sense of Theorem~2.10, of the strongly continuous one-parameter unitary groups
	\begin{equation}\label{eq:generators-convention}
		S(s)=e^{isH/\hbar}, \qquad T\big((0,t\,\mathbf{e}_k)\big)=e^{-itP_k/\hbar}, \qquad s,t\in\mathbb{R},
	\end{equation}
	where $\mathbf{e}_k$ is the canonical basis of $\mathbb{R}^3$. We write $\mathbf{P}=(P_1,P_2,P_3)$.
\end{definition}

\begin{remark}\label{rem:generators-conventions}
	The operator $H$ is the energy operator referred to in Definitions~\ref{def:energy-operator} and~\ref{def:positive-energy}. The signs in Equation~\eqref{eq:generators-convention} are those of the active convention of Remark~\ref{rem:active-convention} and agree with Definition~\ref{def:generators}. With $a\cdot p=a^0p^0-\mathbf{a}\cdot\mathbf{p}$, the relative sign between $H$ and $\mathbf{P}$ is produced by the signature of $\eta$ and not imposed separately. Stone's theorem treats each one-parameter group in isolation, hence Definition~\ref{def:poincare-generators} does not assert that $H,P_1,P_2,P_3$ strongly commute. This is supplied by the joint spectral measure of the next proposition, in which $H$ and $\mathbf{P}$ turn out to be four components of a single vector-valued operator.
\end{remark}

The joint spectral measure of the translations is the object that replaces the choice of a momentum basis in $\mathcal{H}$. Its existence and uniqueness follow from the Stone--Naimark--Ambrose--Godement theorem, and the identification of $H$ and $\mathbf{P}$ with its first moments follows from the uniqueness in Stone's theorem.

\begin{proposition}[Spectral measure of the translations]\label{prop:snag-translations}
	Let $\Pi$ be the lift of a Poincar\'e action $\Phi$ as per Proposition~\ref{prop:poincare-lift}, \textit{i.e.}, assume H1. There is a unique projection valued measure $E$, as per Definition~\ref{def:pvm}, on the Borel sets of $M^*$ such that
	\begin{equation}\label{eq:translation-spectral}
		T(a)=\int_{M^*}e^{ia\cdot p/\hbar}\,dE(p), \qquad a\in M.
	\end{equation}
	Moreover the following hold.
	\begin{enumerate}
		\item $H=\int p^0\,dE$ and $P_k=\int p^k\,dE$ for $k\in\{1,2,3\}$, these four operators strongly commute as per Definition~\ref{def:strongly-commuting}, and $T((0,\mathbf{a}))=e^{-i\mathbf{a}\cdot\mathbf{P}/\hbar}$ with $\mathbf{a}\cdot\mathbf{P}:=\int\mathbf{a}\cdot\mathbf{p}\,dE(p)$.
		\item The operator $P^2:=\int p\cdot p\,dE$, with domain $D(P^2)=\{x\in\mathcal{H}\mid\int(p\cdot p)^2\,d\langle x,Ex\rangle<\infty\}$, is selfadjoint and strongly commutes with $H$ and $\mathbf{P}$. On $D(H^2)\cap D(|\mathbf{P}|^2)$ it acts as $H^2-|\mathbf{P}|^2$.
	\end{enumerate}
\end{proposition}

\begin{proof}
	Every continuous character of $M$ has the form $a\mapsto e^{i\lambda(a)/\hbar}$ for a unique linear form $\lambda$ on $M$. Since $\eta$ is nondegenerate, $\lambda=\eta(\,\cdot\,,p)$ for a unique $p\in M^*$, where $M^*$ is identified with $M$ through $\eta$. Hence $p\mapsto\chi_p$, $\chi_p(a):=e^{ia\cdot p/\hbar}$, is an isomorphism of topological groups from $M^*$ onto the dual group of $M$. Theorem~\ref{thm:snag} applied to $T$ yields a unique projection valued measure on the dual group with $T(a)=\int\chi(a)\,dE(\chi)$, and transporting it through this isomorphism yields $E$ with Equation~\eqref{eq:translation-spectral} and its uniqueness.
	
	Item 1. For $a=(s,\mathbf{0})$ one has $a\cdot p=s\,p^0$, so that $S(s)=\int e^{isp^0/\hbar}dE(p)=e^{isH'/\hbar}$ with $H':=\int p^0\,dE$ selfadjoint by the functional calculus of Proposition~\ref{prop:pvm-calculus}. The uniqueness in Theorem~2.10 yields $H=H'$. For $a=(0,t\,\mathbf{e}_k)$ one has $a\cdot p=-t\,p^k$, and the same argument yields $P_k=\int p^k\,dE$. The spectral projections of $\int f\,dE$ are $E(f^{-1}(B))$ for the Borel sets $B\subseteq\mathbb{R}$, by Proposition~\ref{prop:pvm-calculus}, and those associated with different real Borel functions commute because all of them are built from the single measure $E$. It descends that $H,P_1,P_2,P_3$ strongly commute. The formula for $T((0,\mathbf{a}))$ is Equation~\eqref{eq:translation-spectral} with $a\cdot p=-\mathbf{a}\cdot\mathbf{p}$.
	
	Item 2. Since $p\mapsto p\cdot p$ is a real Borel function, $P^2$ is selfadjoint and strongly commutes with the operators of item 1 by the same argument. If $x\in D(H^2)\cap D(|\mathbf{P}|^2)$, then $\int(p^0)^4d\langle x,Ex\rangle<\infty$ and $\int|\mathbf{p}|^4d\langle x,Ex\rangle<\infty$, and $(p\cdot p)^2\leq2\big((p^0)^4+|\mathbf{p}|^4\big)$ yields $x\in D(P^2)$. On such $x$ one has $P^2x=\int\big((p^0)^2-|\mathbf{p}|^2\big)dE\,x=H^2x-|\mathbf{P}|^2x$.
\end{proof}

\noindent The point $p$ is a character of $M$ and not a point of $\mathbb{P}\mathcal{H}$ or of $M$, and $E(d)$ is the projection onto the vectors whose four-momentum lies in $d$.

The Lorentz group does not act on the translations by conjugation in a trivial way, since $(0,g)(a,\mathbb{1})(0,g)^{-1}=(\Lambda(g)a,\mathbb{1})$ by Lemma~\ref{lem:poincare-relations}. This relation transfers to $E$ through the uniqueness in Proposition~\ref{prop:snag-translations}, and its consequence is that the support of $E$ is a union of Lorentz orbits in $M^*$.

\begin{lemma}[Covariance of the spectral measure]\label{lem:poincare-covariance}
	Let $\Pi$ and $E$ be as per Proposition~\ref{prop:snag-translations}, \textit{i.e.}, assume H1. For every Borel set $d\subseteq M^*$ and every $\kappa=(a,g)\in\mathcal{P}$,
	\begin{equation}\label{eq:covariance}
		\Pi(\kappa)\,E(d)\,\Pi(\kappa)^{-1}=E\big(\Lambda(g)d\big),
	\end{equation}
	where $\Lambda(g)$ acts on $M^*\cong M$ through $\eta$, \textit{i.e.}, $\Lambda(g)a\cdot\Lambda(g)p=a\cdot p$. Consequently
	\begin{enumerate}
		\item for every Borel function $f:M^*\to\mathbb{R}$, with $f(P):=\int f\,dE$ and as operators with their domains,
		\begin{equation*}
			\Pi((0,g))\,f(P)\,\Pi((0,g))^{-1}=\int f\big(\Lambda(g)^{-1}p\big)\,dE(p),
		\end{equation*}
		\item $\Lambda(g)\operatorname{supp}E=\operatorname{supp}E$ for every $g\in SL(2,\mathbb{C})$, with $\operatorname{supp}E$ as per Definition~\ref{def:pvm-support}.
	\end{enumerate}
\end{lemma}

\begin{proof}
	Set $U:=\Pi((0,g))$. By item 2 of Lemma~\ref{lem:poincare-relations}, $UT(a)U^{-1}=T(\Lambda(g)a)$. By the invariance of $\eta$, the right-hand side equals $\int e^{ia\cdot\Lambda(g)^{-1}p/\hbar}dE(p)=\int e^{ia\cdot q/\hbar}dE'(q)$, where $E'(d):=E(\Lambda(g)d)$ is the image of $E$ under $p\mapsto\Lambda(g)^{-1}p$. The left-hand side equals $\int e^{ia\cdot p/\hbar}dE_U(p)$ with $E_U(d):=UE(d)U^{-1}$. Hence $E_U$ and $E'$ are two projection valued measures with the same transform in the sense of Equation~\eqref{eq:translation-spectral}, and they coincide by the uniqueness in Proposition~\ref{prop:snag-translations}. For a general $\kappa$, write $(a,g)=(a,\mathbb{1})(0,g)$. The operator $T(a)=\int e^{ia\cdot p/\hbar}dE$ commutes with every $E(d)$, so that $\Pi(\kappa)E(d)\Pi(\kappa)^{-1}=T(a)E(\Lambda(g)d)T(a)^{-1}=E(\Lambda(g)d)$.
	
	Item 1 is the change of variables in the functional calculus, as in the proof of Lemma~\ref{lem:boost-covariance}.
	
	Item 2. Let $N:=M^*\setminus\operatorname{supp}E$, the largest open $E$-null set as per Definition~\ref{def:pvm-support}. The set $\Lambda(g)N$ is open, and $E(\Lambda(g)N)=UE(N)U^{-1}=0$ by Equation~\eqref{eq:covariance}. Hence $\Lambda(g)N\subseteq N$, and the same argument applied to $g^{-1}$ yields $\Lambda(g)N=N$.
\end{proof}

\noindent Item 2 states that the possible values of the four-momentum are permuted among themselves by changes of inertial frame.

The orbits of $\Lambda(SL(2,\mathbb{C}))$ in $M^*$ are the sheets of the hyperboloids $p\cdot p=\mathrm{const}$ together with the null cone and the origin. The covariance of $E$ and the positivity of the energy select among them. The scalar $P^2$ is invariant under the Lorentz group and is a function of $E$, hence it commutes with the whole representation. The next proposition uses H1 and H2 in item 1, and H3 in addition in items 2 and 3.

\begin{proposition}[The spectral measure lives on a mass shell]\label{prop:mass-shell}
	Let $\Phi$ be an elementary relativistic system as per Definition~\ref{def:elementary-system}, with lift $\Pi$ and spectral measure $E$. Then
	\begin{enumerate}
		\item there is $m^2\in\mathbb{R}$ with $P^2=m^2\mathbb{1}$, and $E\big(\{p\in M^*\mid p\cdot p\neq m^2\}\big)=0$,
		\item $m^2\geq0$,
		\item if $m^2>0$ and $m$ denotes its positive square root, then $E(M^*\setminus\mathcal{O}_m^+)=0$ and $\operatorname{supp}E=\mathcal{O}_m^+$, with $\mathcal{O}_m^+$ as per Definition~\ref{def:mass-shell}.
	\end{enumerate}
\end{proposition}

\begin{proof}
	Item 1. By item 1 of Lemma~\ref{lem:poincare-covariance} and the invariance of $p\mapsto p\cdot p$ under $\Lambda(g)$,
	\begin{equation*}
		\Pi((0,g))\,e^{itP^2}\,\Pi((0,g))^{-1}=\int e^{it(\Lambda(g)^{-1}p)\cdot(\Lambda(g)^{-1}p)}dE(p)=e^{itP^2}, \qquad t\in\mathbb{R}.
	\end{equation*}
	The operator $e^{itP^2}$ also commutes with every $T(a)$, since both are functions of $E$. As $(a,g)=(a,\mathbb{1})(0,g)$, it commutes with $\Pi(\mathcal{P})$. The set $\Pi(\mathcal{P})$ is irreducible. Indeed, a closed subspace $\mathcal{K}$ invariant under $\Pi(\mathcal{P})$ satisfies $\Phi_\kappa(\mathbb{P}\mathcal{K})\subseteq\mathbb{P}\mathcal{K}$ for all $\kappa\in\mathcal{P}$, and Definition~\ref{def:poincare-irreducible} forces $\mathcal{K}\in\{\{0\},\mathcal{H}\}$, as in Step 5 of the proof of Theorem~\ref{thm:free-hamiltonian}. Lemma~\ref{lem:schur-infinite} applied to $A=P^2$ yields $P^2=m^2\mathbb{1}$ for some $m^2\in\mathbb{R}$. It descends that for every $x\in\mathcal{H}$
	\begin{equation*}
		\int\big(p\cdot p-m^2\big)^2\,d\langle x,Ex\rangle=\big\|(P^2-m^2)x\big\|^2=0,
	\end{equation*}
	so that the open set $\{p\cdot p\neq m^2\}$ is $E$-null.
	
	Items 2 and 3. Let $S:=\operatorname{supp}E$. It is nonempty because $E(M^*)=\mathbb{1}$, it is contained in $\{p\cdot p=m^2\}$ by item 1, and it is invariant under $\Lambda(SL(2,\mathbb{C}))$ by item 2 of Lemma~\ref{lem:poincare-covariance}. Every open set meeting $S$ has nonzero $E$-measure. Moreover $E(\{p^0<-n\})$ is the spectral projection of $H$ on $(-\infty,-n)$ by item 1 of Proposition~\ref{prop:snag-translations}, and it vanishes for large $n$ if and only if $H$ is bounded below.
	
	Suppose that $m^2<0$ and set $\nu:=\sqrt{-m^2}$. By item 2 of Lemma~\ref{lem:hyperboloid-orbits} the hyperboloid $\{p\cdot p=-\nu^2\}$ is a single orbit, hence $S$ equals the whole hyperboloid. On it $p^0$ takes every real value, because $|\mathbf{p}|^2=(p^0)^2+\nu^2$ imposes no restriction on $p^0$. Hence every open set $\{p^0<-n\}$ meets $S$, so that $E(\{p^0<-n\})\neq0$ for all $n\in\mathbb{N}$, and $H$ is not bounded below. This contradicts H3, and item 2 follows.
	
	Suppose that $m^2>0$. By item 1, $S\subseteq\mathcal{O}_m^+\cup\mathcal{O}_m^-$. If $S$ met $\mathcal{O}_m^-$, then $\mathcal{O}_m^-\subseteq S$, since $\mathcal{O}_m^-$ is a single orbit by item 1 of Lemma~\ref{lem:hyperboloid-orbits}, and $p^0\leq-m$ would be unbounded below on $S$. The argument above would again contradict H3. Hence $S\cap\mathcal{O}_m^-=\emptyset$, so that $\mathcal{O}_m^-\subseteq M^*\setminus S$ is $E$-null, and together with item 1 this yields $E(M^*\setminus\mathcal{O}_m^+)=0$. Finally $S\subseteq\mathcal{O}_m^+$ is nonempty and invariant, and $\mathcal{O}_m^+$ is a single orbit, so $S=\mathcal{O}_m^+$.
\end{proof}

\noindent Hypothesis H3 enters only through the exclusion of the sheet $\mathcal{O}_m^-$ and of the hyperboloids of negative squared mass.

\begin{remark}[The mass as the eigenvalue of a Casimir operator]\label{rem:casimir-label}
	The operator $P^2$ commutes with $\Pi(\mathcal{P})$, as shown in the proof of Proposition~\ref{prop:mass-shell}, so that it is a Casimir operator of the representation, and Schur's lemma makes it a scalar on an irreducible one. The mass is therefore the label of an irreducible representation, which replaces the Galilean statement of Chapter~\ref{ch:analytical-hamiltonians}, where the mass was the parameter of a central extension in Theorem~\ref{thm:bargmann} and Proposition~\ref{prop:lift}. The origin of the difference is the vanishing of $H^2(\mathfrak{p},\mathbb{R})$ in Theorem~\ref{thm:poincare-cohomology}, which leaves no phase from which a mass could be read. The same vanishing is behind a second difference. In Proposition~\ref{prop:lift} the generator $H$ was determined only up to an additive constant, whereas here the lift is unique by Proposition~\ref{prop:poincare-lift} and the zero of the energy is fixed by the support of $E$. A second Casimir operator, built from the Pauli--Lubanski vector \cite[TODO-location]{TODO-BarutRaczka}, would distinguish the spins, but it is not needed. For $m>0$ the spin is read from the representation of the stabilizer $SU(2)$ of the rest momentum, as shown in Theorem~\ref{thm:wigner-mackey} and Remark~\ref{rem:spin-emerges}.
\end{remark}

\begin{remark}[The cases $m^2<0$ and $m=0$]\label{rem:massless}
	The case $m^2<0$ is excluded by H3 in Proposition~\ref{prop:mass-shell} and would describe tachyons. The case $m=0$ is excluded from now on. The null cone $\{p\cdot p=0\}$ is the union of the origin, which is an orbit by itself, and of two nonzero sheets, and the stabilizer in $SL(2,\mathbb{C})$ of a nonzero null momentum is not $SU(2)$ but the covering of the Euclidean group of the plane \cite[TODO-location]{TODO-BarutRaczka}. The construction of the next section relies on Proposition~\ref{prop:orbit}, in which the stabilizer of the rest momentum is $SU(2)$, so that it does not apply. The same exclusion was made for the Galilean mass in the case $m=0$ of Remark~\ref{rem:weyl}.
\end{remark}

\noindent What remains undetermined is the multiplicity of $E$ above each point of $\mathcal{O}_m^+$ and the action of the Lorentz group on it, which is settled in the next section by Theorem~\ref{thm:wigner-mackey}.

\section{Classification of the Elementary Systems}
\label{sec:classification}

The question of this section is how the pure states of an elementary relativistic system are described when the hypotheses H1 to H3 of Definition~\ref{def:hypotheses} are the only information available, \textit{i.e.}, which data a ray of $\mathbb{P}\mathcal{H}$ has to specify besides the momentum. The datum is the pair $(\mathbb{P}\mathcal{H},\Phi)$ together with its unitary lift $\Pi$ of Proposition~\ref{prop:poincare-lift}. The translations of $\Pi$ have the spectral measure $E$ of Proposition~\ref{prop:snag-translations}, which is concentrated on the forward mass shell $\mathcal{O}_m^+$ by Proposition~\ref{prop:mass-shell}. The set $\mathcal{O}_m^+$ is not a set of states. It is the support of $E$, \textit{i.e.}, the set of the possible values of the four-momentum. What is obtained is a model of the states on a space of square integrable functions on $\mathcal{O}_m^+$ with values in a finite-dimensional spin space, the model denoted M2 in the opening of the chapter. In it a state is the ray of a function $\psi$, the vector $\psi(p)$ is the amplitude of the momentum $p$ together with the amplitudes of the spin components, and the Poincar\'e group acts through an explicit formula containing the Wigner rotation of Proposition~\ref{prop:induced-action}. Each elementary system is labelled by the mass $m$ and by a number $j\in\{0,\tfrac12,1,\dots\}$, which is not an input but the label of an irreducible representation of the stabilizer $SU(2)$ of the rest momentum. This is the classification step of the chain of the chapter. Its output, the model on $\mathcal{O}_m^+$ with the energy $H=\omega(\mathbf{P})$, is the starting point of the next section, where the dependence of the Wigner rotation on the momentum is shown to obstruct a description of the states by fields transforming with fixed matrices.

The form of the representation is motivated as follows. Since $E$ is concentrated on $\mathcal{O}_m^+$, a realization of $\mathcal{H}$ as a space of square integrable functions on $\mathcal{O}_m^+$ with values in a Hilbert space $\mathcal{K}$, in which $E$ acts by multiplication by characteristic functions, makes the translations act by multiplication by $e^{ia\cdot p/\hbar}$. By Lemma~\ref{lem:poincare-covariance} the unitary $\Pi((0,g))$ transforms $E(d)$ into $E(\Lambda(g)d)$, so it carries functions supported on $d$ to functions supported on $\Lambda(g)d$. Since the measure $\mu$ of Lemma~\ref{lem:invariant-measure} is invariant under $\Lambda(g)$, this suggests the form $\psi\mapsto R(g,p)\,\psi(\Lambda(g)^{-1}p)$ with $R(g,p)$ a unitary operator on $\mathcal{K}$. The homomorphism property of $\Pi$ then makes $R$ a cocycle, \textit{i.e.}, $R(gg',p)=R(g,p)R(g',\Lambda(g)^{-1}p)$. The induction developed among the preliminaries of the chapter provides the canonical cocycle $R(g,p)=\pi(R_W(g,p))$, with $\pi$ a unitary representation of the stabilizer $SU(2)$ and $R_W$ the Wigner rotation of Proposition~\ref{prop:induced-action}, see also Theorem~\ref{thm:equivariant-sections}. This is an argument of motivation and not a derivation. That every elementary system has this form is the content of Theorem~\ref{thm:wigner-mackey}.

\begin{definition}[Spin data]\label{def:spin-data}
	Let $j\in\{0,\tfrac12,1,\dots\}$. The \emph{spin data} of label $j$ is the pair $(\pi,V)$ with $V:=\mathbb{C}^{2j+1}$ and $\pi:=d_j$ the irreducible unitary representation of $SU(2)$ of dimension $2j+1$, as per Definition~1.2.25, denoted $d_j$ in Corollary~\ref{cor:schrodinger-realization}.
\end{definition}

The space $V$ collects the internal amplitudes attached to a single value of the momentum, and $\pi$ prescribes how the stabilizer of the rest momentum rotates them. Letting the amplitudes depend on $p\in\mathcal{O}_m^+$ yields a Hilbert space of functions of the momentum, whose vectors are states with $\psi(p)$ the amplitude of the momentum $p$ and of the spin components, and on which the Poincar\'e group acts through the Wigner rotation.

\begin{definition}[Induced representation]\label{def:induced-rep}
	Let $(\pi,V)$ be the spin data of label $j$ as per Definition~\ref{def:spin-data}, let $\mu$ be the measure of Lemma~\ref{lem:invariant-measure} and set $\mathcal{H}_{m,j}:=L^2(\mathcal{O}_m^+,\mu;V)$. The \emph{induced representation} $U_{m,j}$ takes $(a,g)\in\mathcal{P}$ and yields the operator on $\mathcal{H}_{m,j}$ defined by
	\begin{equation}\label{eq:induced-rep}
		\big(U_{m,j}(a,g)\psi\big)(p):=e^{ia\cdot p/\hbar}\,\pi\big(R_W(g,p)\big)\,\psi\big(\Lambda(g)^{-1}p\big),
	\end{equation}
	where $R_W(g,p)=A_p^{-1}\,g\,A_{\Lambda(g)^{-1}p}\in SU(2)$ is the Wigner rotation of Proposition~\ref{prop:induced-action}, with $A_p$ the standard boost of Proposition~\ref{prop:orbit}.
\end{definition}

By Theorem~\ref{thm:equivariant-sections} and Proposition~\ref{prop:action-trivialization}, $\mathcal{H}_{m,j}$ is the space of square integrable sections of the bundle associated to the spin data, written in the trivialization $p\mapsto A_p$, and Equation~\eqref{eq:induced-rep} is the action of the left translations in that trivialization combined with the multiplication by the phase of the translations. The cocycle identity of $R_W$ is the property that turns the formula into a representation, as the following proposition verifies together with the properties of the spectral measure.

\begin{proposition}[Unitarity of the induced representation]\label{prop:induced-unitary}
	$U_{m,j}$ is a strongly continuous unitary representation of $\mathcal{P}$ on $\mathcal{H}_{m,j}$, as per Definition~1.2.28. The spectral measure of its translations is $E_{m,j}(d)=\chi_d$, \textit{i.e.}, the multiplication by the characteristic function of the Borel set $d\subseteq\mathcal{O}_m^+$. Consequently $P^\mu$ acts by multiplication by $p^\mu$, $P^\mu P_\mu=m^2\,\mathbb{1}$, and $H$ acts by multiplication by $\omega(\mathbf{p})\geq m$.
\end{proposition}

\begin{proof}
	We write $U:=U_{m,j}$.
	
	\emph{Homomorphism.} By Proposition~\ref{prop:induced-action}, $R_W(\mathbb{1},p)=\mathbb{1}$ and $R_W(gg',p)=R_W(g,p)R_W(g',\Lambda(g)^{-1}p)$. Together with the identity $a'\cdot\Lambda(g)^{-1}p=\Lambda(g)a'\cdot p$ and with the product $(a,g)(a',g')=(a+\Lambda(g)a',gg')$ of Definition~\ref{def:poincare-group}, this yields
	\begin{align*}
		\big(U(a,g)U(a',g')\psi\big)(p)&=e^{ia\cdot p/\hbar}\pi\big(R_W(g,p)\big)\,e^{i\Lambda(g)a'\cdot p/\hbar}\pi\big(R_W(g',\Lambda(g)^{-1}p)\big)\psi\big(\Lambda(gg')^{-1}p\big)\\
		&=e^{i(a+\Lambda(g)a')\cdot p/\hbar}\,\pi\big(R_W(gg',p)\big)\,\psi\big(\Lambda(gg')^{-1}p\big)=\big(U((a,g)(a',g'))\psi\big)(p).
	\end{align*}
	Since $U(0,\mathbb{1})=\mathbb{1}$, every $U(a,g)$ is invertible with inverse $U((a,g)^{-1})$.
	
	\emph{Unitarity.} The operator $\pi(R_W(g,p))$ is unitary and $\mu$ is invariant under $\Lambda(g)$ by Lemma~\ref{lem:invariant-measure}, so that
	\begin{equation*}
		\|U(a,g)\psi\|^2=\int\|\psi(\Lambda(g)^{-1}p)\|^2\,d\mu(p)=\|\psi\|^2.
	\end{equation*}
	
	\emph{Strong continuity.} Let $\psi$ be continuous with compact support. Since $p\mapsto A_p$ is smooth, the function $(a,g,p)\mapsto e^{ia\cdot p/\hbar}\pi(R_W(g,p))\psi(\Lambda(g)^{-1}p)$ is continuous. For $(a,g)$ in a compact set the support of $U(a,g)\psi$ lies in a fixed compact set, and its norm is bounded by $\sup\|\psi\|$. Dominated convergence yields the continuity of $(a,g)\mapsto U(a,g)\psi$ in $L^2$. Such vectors are dense and the operators $U(a,g)$ are isometries, so continuity holds for every $\psi\in\mathcal{H}_{m,j}$.
	
	\emph{Translations.} By Equation~\eqref{eq:induced-rep}, $U(a,\mathbb{1})$ is the multiplication by $e^{ia\cdot p/\hbar}$, \textit{i.e.}, $\int e^{ia\cdot p/\hbar}\,d\chi_d(p)$, and the uniqueness in Proposition~\ref{prop:snag-translations} identifies $E_{m,j}$ with $\chi$. The remaining statements follow from $p\cdot p=m^2$ and $p^0=\omega(\mathbf{p})\geq m$ on $\mathcal{O}_m^+$.
\end{proof}

The representation $U_{m,j}$ contains a topological datum, the value taken by the element of $\mathcal{P}$ that covers the identity of the proper orthochronous isometries. It is the same datum that distinguishes rotations of $2\pi$ on integer and half-integer spin in Example~\ref{ex:so3-su2}, and the next proposition shows that it is a sign in every irreducible representation.

\begin{proposition}[The central sign]\label{prop:central-sign}
	Let $z:=(0,-\mathbb{1})\in\mathcal{P}$. Then $z$ is central, $z^2$ is the identity of $\mathcal{P}$, and for every irreducible strongly continuous unitary representation $\Pi$ of $\mathcal{P}$ there is $s_\Pi\in\{1,-1\}$ with $\Pi(z)=s_\Pi\,\mathbb{1}$.
\end{proposition}

\begin{proof}
	By Lemma~\ref{lem:sl2c-covering}, $\Lambda(-\mathbb{1})=\mathrm{id}$, so $z(a,g)=(a,-g)=(a,g)z$ and $z^2=(0,\mathbb{1})$. The operator $\Pi(z)$ is unitary with $\Pi(z)^2=\mathbb{1}$, hence $\Pi(z)^*=\Pi(z)$ and $\Pi(z)$ is selfadjoint. It commutes with $\Pi(\mathcal{P})$ because $z$ is central, and so does $e^{it\Pi(z)}=\cos t\,\mathbb{1}+i\sin t\,\Pi(z)$. Lemma~\ref{lem:schur-infinite} yields $\Pi(z)=s_\Pi\mathbb{1}$ with $s_\Pi\in\mathbb{R}$, and $s_\Pi^2=1$.
\end{proof}

\noindent The lift of an irreducible Poincar\'e action is irreducible in the sense required here, since a closed subspace invariant under $\Pi$ is invariant under $\Phi$ as per Definition~\ref{def:poincare-irreducible}.

The proposition leaves $s_\Pi$ undetermined, and the induced representation allows one to compute it. The computation connects the sign to the label $j$ of the spin data, without any reference to the particular value of the momentum.

\begin{corollary}[The central sign is $(-1)^{2j}$]\label{cor:central-sign-spin}
	In $U_{m,j}$ one has $U_{m,j}(z)=(-1)^{2j}\,\mathbb{1}$. Hence, whenever $U_{m,j}$ is irreducible, as shown in Theorem~\ref{thm:wigner-mackey}, its central sign is $s_\Pi=(-1)^{2j}$. For integer $j$ it factors through the proper orthochronous Poincar\'e group of Corollary~\ref{cor:poincare-topology}, and for half-integer $j$ it does not.
\end{corollary}

\begin{proof}
	Since $\Lambda(-\mathbb{1})=\mathrm{id}$, Equation~\eqref{eq:induced-rep} yields $(U_{m,j}(z)\psi)(p)=\pi(R_W(-\mathbb{1},p))\psi(p)$ and $R_W(-\mathbb{1},p)=A_p^{-1}(-\mathbb{1})A_p=-\mathbb{1}$. By Example~1.2.17, $-\mathbb{1}=\exp(-2\pi i\,\sigma_3/2)$ in $SU(2)$, so $d_j(-\mathbb{1})=e^{-2\pi iJ_3}$. The eigenvalues of $J_3$ in $d_j$ are $j,j-1,\dots,-j$ \cite[TODO-location]{TODO-Hal03}, all congruent to $j$ modulo integers, and therefore $d_j(-\mathbb{1})=e^{-2\pi ij}\,\mathbb{1}=(-1)^{2j}\,\mathbb{1}$.
\end{proof}

The representations $U_{m,j}$ are candidates, and nothing so far guarantees that an elementary system is equivalent to one of them. The classification theorem closes this gap, and it does so by showing that the label $j$ is determined by the action $\Phi$ and not chosen.

\begin{theorem}[Wigner, Mackey]\label{thm:wigner-mackey}
	Let $m>0$.
	\begin{enumerate}
		\item For every $j\in\{0,\tfrac12,1,\dots\}$ the representation $U_{m,j}$ of Definition~\ref{def:induced-rep} is an irreducible strongly continuous unitary representation of $\mathcal{P}$ of positive energy with $P^\mu P_\mu=m^2\mathbb{1}$.
		\item Let $\Phi$ be an elementary relativistic system as per Definition~\ref{def:elementary-system}, \textit{i.e.}, satisfying H1 to H3 of Definition~\ref{def:hypotheses}, with lift $\Pi$ as per Proposition~\ref{prop:poincare-lift} and mass $m>0$ as per Proposition~\ref{prop:mass-shell}. Then there exist a unique $j\in\{0,\tfrac12,1,\dots\}$ and a unitary operator $W:\mathcal{H}\to\mathcal{H}_{m,j}$ such that $W\,\Pi(\kappa)\,W^{-1}=U_{m,j}(\kappa)$ for all $\kappa\in\mathcal{P}$ and $W\,E(d)\,W^{-1}=\chi_d$ for every Borel set $d\subseteq\mathcal{O}_m^+$.
	\end{enumerate}
\end{theorem}

\noindent See \cite[TODO-location]{TODO-Wigner1939} and \cite[TODO-location]{TODO-Mackey1958}.

The proof is not reproduced. The following remark isolates the role of each ingredient, and it is the scheme of the proof.

\begin{remark}[Reading the theorem in four steps]\label{rem:mackey-steps}
	\emph{Step 1, a system of imprimitivity.} By Lemma~\ref{lem:poincare-covariance} and Proposition~\ref{prop:mass-shell}, $E$ is a projection valued measure concentrated on $\mathcal{O}_m^+$ with $\Pi((0,g))E(d)\Pi((0,g))^{-1}=E(\Lambda(g)d)$. Hence $g\mapsto\Pi((0,g))$ and $E$ form a system of imprimitivity for $SL(2,\mathbb{C})$ based on the homogeneous space $\mathcal{O}_m^+\cong SL(2,\mathbb{C})/SU(2)$ of Proposition~\ref{prop:orbit}.
	
	\emph{Step 2, the induced form.} Mackey's imprimitivity theorem identifies $\mathcal{H}$ with $L^2(\mathcal{O}_m^+,\mu;\mathcal{K})$ for a separable Hilbert space $\mathcal{K}$, with $E$ acting by multiplication by characteristic functions, and $SL(2,\mathbb{C})$ acting as in the representation induced from a unitary representation $d$ of $SU(2)$ on $\mathcal{K}$. In the trivialization of Proposition~\ref{prop:action-trivialization} this action is $\psi\mapsto d(R_W(g,p))\psi(\Lambda(g)^{-1}p)$. The measure $\mu$ of Lemma~\ref{lem:invariant-measure} belongs to the unique quasi-invariant class on the homogeneous space, so that no Radon--Nikodym factor appears.
	
	\emph{Step 3, the translations.} By Proposition~\ref{prop:snag-translations}, $T(a)=\int e^{ia\cdot p/\hbar}\,dE(p)$, so the translations act by multiplication by $e^{ia\cdot p/\hbar}$. Since $(a,g)=(a,\mathbb{1})(0,g)$, the representation $\Pi$ is given by Equation~\eqref{eq:induced-rep} with $(d,\mathcal{K})$ in place of $(\pi,V)$.
	
	\emph{Step 4, irreducibility selects the spin.} If $\mathcal{K}_1\subseteq\mathcal{K}$ is a closed subspace invariant under $d$, then $L^2(\mathcal{O}_m^+,\mu;\mathcal{K}_1)$ is invariant under $\Pi$, so irreducibility forces $d$ to be irreducible. Since $SU(2)$ is compact, $d$ is finite-dimensional and equivalent to $d_j$ for a unique $j$ \cite[TODO-location]{TODO-Hal03}, so that $\mathcal{K}\cong\mathbb{C}^{2j+1}$. Equivalence classes of induced representations correspond bijectively to equivalence classes of the inducing representations, which yields the uniqueness of $j$. Conversely, the irreducibility of $d_j$ yields that of $U_{m,j}$.
\end{remark}

The theorem fixes the multiplicity of $E$ and the action of the Lorentz group, and hence the energy as a function of the momentum. Since the energy is the generator of the time translations, its expression follows from the form of $E$ alone, with no further use of the group structure.

\begin{corollary}[Hamiltonian of the relativistic free particle]\label{thm:relativistic-hamiltonian}
	Let $\Phi$ be an elementary relativistic system of mass $m>0$, \textit{i.e.}, satisfying H1 to H3 of Definition~\ref{def:hypotheses}, and let $H$ and $\mathbf{P}$ be as per Definition~\ref{def:poincare-generators}. Then
	\begin{equation}\label{eq:relativistic-hamiltonian}
		H=\omega(\mathbf{P})=\sqrt{|\mathbf{P}|^2+m^2},\qquad D(H)=D(|\mathbf{P}|),
	\end{equation}
	the spectrum is $\sigma(H)=[m,\infty)$, and the spectral measure of $\mathbf{P}$ has support $\mathbb{R}^3$.
\end{corollary}

\begin{proof}
	Let $W$ and $j$ be as in item 2 of Theorem~\ref{thm:wigner-mackey}. Since $W E(d)W^{-1}=\chi_d$, the functional calculus yields that $WHW^{-1}=\int p^0\,d\chi$ and $W\mathbf{P}W^{-1}=\int\mathbf{p}\,d\chi$ are the multiplication operators by $p^0=\omega(\mathbf{p})$ and by $\mathbf{p}$ on $\mathcal{H}_{m,j}$, with maximal domains. Through the homeomorphism $\mathbf{p}\mapsto(\omega(\mathbf{p}),\mathbf{p})$ of Definition~\ref{def:mass-shell}, $\mathcal{O}_m^+\cong\mathbb{R}^3$ and $\mu=d^3\mathbf{p}/2\omega(\mathbf{p})$, which has a positive continuous density. Hence $H=\omega(\mathbf{P})$, and every nonempty open set has positive $\mu$ measure, so that the spectral measure of $\mathbf{P}$ has support $\mathbb{R}^3$. Since $|\mathbf{p}|\leq\omega(\mathbf{p})\leq|\mathbf{p}|+m$, one has $D(H)=D(|\mathbf{P}|)$.
	
	For the spectrum, let $\lambda\geq m$ and $\mathbf{p}_0$ with $\omega(\mathbf{p}_0)=\lambda$, which exists because $\omega$ is continuous with $\omega(\mathbb{R}^3)=[m,\infty)$. For $\varepsilon>0$ let $\psi_\varepsilon$ be the normalized characteristic function of a ball centred at $\mathbf{p}_0$ on which $|\omega-\lambda|<\varepsilon$. Then $\|(H-\lambda)\psi_\varepsilon\|\leq\varepsilon$, so $\lambda\in\sigma(H)$. If $\lambda\notin[m,\infty)$, then $|\omega-\lambda|\geq\operatorname{dist}(\lambda,[m,\infty))>0$, and $(H-\lambda)^{-1}$ is the bounded multiplication by $(\omega-\lambda)^{-1}$. It descends that $\sigma(H)=[m,\infty)$.
\end{proof}

\noindent The statement is independent of $j$, and no relation between boosts and time translations was used, since the transitivity of the Lorentz group on the mass shell is already contained in the support of $E$.

\begin{remark}[Spin is an output]\label{rem:spin-emerges}
	Theorem~\ref{thm:wigner-mackey} answers the question left open by Remark~\ref{rem:ph-no-spin}. The state space $\mathbb{P}\mathcal{H}$ does not contain the spin, which is a property of the action $\Phi$. It enters as the label of the irreducible representation of the stabilizer of the rest momentum that irreducibility selects, and integer and half-integer values are distinguished by the central sign $s_\Pi=(-1)^{2j}$ of Corollary~\ref{cor:central-sign-spin}. In the model $\mathcal{H}_{m,j}$ the vector $\psi(p)$ collects the amplitudes of the spin components at the momentum $p$, and for a unit vector $\psi$ one has $\langle\psi,E(d)\psi\rangle=\int_d\|\psi(p)\|^2\,d\mu(p)$, so that $\|\psi(p)\|^2d\mu(p)$ is the distribution of the four-momentum in the state $[\psi]$. The components are defined through the standard boost $A_p$, hence they are the components of the spin in the rest frame reached by $A_p^{-1}$, and a boost acts on them through the Wigner rotation of Proposition~\ref{prop:wigner-rotation}, which depends on $p$. In the Galilean model of Corollary~\ref{cor:schrodinger-realization} the boosts act trivially on the factor $\mathbb{C}^{2j+1}$, and this difference is what the next section turns into an obstruction. From the next section on we fix $j=\tfrac12$, which is a choice of label and not a hypothesis. The hypotheses H4 to H6 of Definition~\ref{def:hypotheses} have not been used.
\end{remark}

\begin{proof}[Proof of Theorem~\ref{thm:main}, part (A)]
	Only H1 to H3 are used. By Proposition~\ref{prop:poincare-lift}, $\Phi$ has a unique unitary lift $\Pi$, and by item 2 of Theorem~\ref{thm:wigner-mackey} there exist a unique $j$ and a unitary operator $W:\mathcal{H}\to\mathcal{H}_{m,j}$ with $W\Pi(\kappa)W^{-1}=U_{m,j}(\kappa)$. By Lemma~\ref{lem:kahler-iso-two-spaces}, $F:[x]\mapsto[Wx]$ is an isomorphism of Kähler manifolds from $\mathbb{P}\mathcal{H}$ onto $\mathbb{P}\mathcal{H}_{m,j}$, and
	\begin{equation*}
		F\big(\Phi_\kappa[x]\big)=[W\Pi(\kappa)x]=[U_{m,j}(\kappa)Wx]=[U_{m,j}(\kappa)\,\cdot\,]\,F[x],\qquad \kappa\in\mathcal{P},\ [x]\in\mathbb{P}\mathcal{H}.
	\end{equation*}
	This is the existence of the isomorphism, \textit{i.e.}, a realization as per Definition~\ref{def:realization}.
	
	For the uniqueness, let $F':\mathbb{P}\mathcal{H}\to\mathbb{P}\mathcal{H}_{m,j'}$ be a Kähler isomorphism intertwining $\Phi_\kappa$ with $[U_{m,j'}(\kappa)\,\cdot\,]$. The latter is a Poincar\'e action, since the strong continuity of $U_{m,j'}$ makes the expectation values continuous along orbits, and hence the orbit maps by Proposition~3.2.9. By the uniqueness in Proposition~\ref{prop:poincare-lift} its lift is $U_{m,j'}$, and Proposition~\ref{prop:realization-equivalence} yields a unitary operator $\mathcal{H}\to\mathcal{H}_{m,j'}$ intertwining $\Pi$ with $U_{m,j'}$. The uniqueness of $j$ in Theorem~\ref{thm:wigner-mackey} gives $j'=j$.
	
	Under $W$ the energy operator of $\Pi$ corresponds to the multiplication by $\omega(\mathbf{p})$, and Corollary~\ref{thm:relativistic-hamiltonian} yields $H=\omega(\mathbf{P})$ with $D(H)=D(|\mathbf{P}|)$.
\end{proof}

\section{Locality, Parity and the Dirac Space}
\label{sec:locality}

Theorem~\ref{thm:wigner-mackey} replaces the abstract state space by a concrete one. The states of an elementary relativistic system of mass $m>0$ and spin $j$ are the rays of $L^2(\mathcal{O}_m^+,\mu;\mathbb{C}^{2j+1})$, and $\mathcal{P}$ acts on them through the induced representation $U_{m,j}$. In this model the value of a state at the momentum $p$ is a spin vector. A Lorentz transformation moves the momentum and rotates the spin vector by the Wigner rotation $R_W(g,p)$, which depends on $p$ by Proposition~\ref{prop:wigner-rotation}. The question of this section is whether the same states, with the same action of $\mathcal{P}$, admit a model in which the value of a state at $p$ is transformed by a matrix independent of $p$. This is the form in which a field theory describes states, since a field transforms pointwise under the Lorentz group.

The data are the model $\mathbb{P}L^2(\mathcal{O}_m^+,\mu;\mathbb{C}^{2j+1})$ with the action $[U_{m,j}(\kappa)\,\cdot\,]$, the Wigner cocycle of Proposition~\ref{prop:induced-action} and the finite-dimensional representations of $SL(2,\mathbb{C})$ of Theorem~\ref{thm:finite-reps}. The result has three parts. A local model exists, at the price of constraining the values of a state at each $p$ to a subspace $F(p)$ of a larger space $W$, and Proposition~\ref{prop:no-local-unitary} shows that the matrices of such a model cannot be unitary. When space inversion is required to act on the model by a fixed matrix, the representation of $SL(2,\mathbb{C})$ must contain both chiralities. For spin $\frac12$ the smallest such space is the four-dimensional Dirac space.

Locality is not a consequence of H1 to H3. It is a requirement on the model of the states and is recorded, together with the requirements of parity and minimality, as H4 to H6 in Definition~\ref{def:locality-hypotheses}. This section is the step between the induced model and the Dirac equation. It leaves open the subspace $F(\mathring{p})\subseteq W$, which is determined in the next section by space inversion and $SU(2)$ and then transported along the mass shell.

%--------------------------------------------------------------------
\subsection{Local Covariant Realizations}
\label{subsec:local-realizations}

The states of $\mathcal{H}_{m,j}$ are functions of the momentum, and the local description of the same states replaces the matrix $\pi(R_W(g,p))$ by a matrix $D(g)$ independent of $p$. The data are the spin representation $(\pi,V)$ of Definition~\ref{def:spin-data} and the cocycle $R_W$. This subsection defines the models of this type, shows that they describe the same states with the same action of $\mathcal{P}$ as $\mathcal{H}_{m,j}$, and identifies them with subbundles of a trivial bundle. The unitarity of $D$ is treated in the next subsection.

\begin{definition}[Local covariant realization]\label{def:local-realization}
	Let $(\pi,V)$ be as per Definition~\ref{def:spin-data}. A \emph{local covariant realization} of $(\pi,V)$ is a triple $(D,W,\iota)$ where $W$ is a finite-dimensional complex vector space, $D:SL(2,\mathbb{C})\to GL(W)$ is a continuous homomorphism and $\iota:V\to W$ is an injective linear map such that
	\begin{equation}\label{eq:iota-intertwiner}
		\iota\,\pi(u)=D(u)\,\iota,\qquad u\in SU(2).
	\end{equation}
\end{definition}

The map $\iota$ singles out, inside the restriction of $D$ to $SU(2)$, a copy of the spin representation. No unitarity is required of $D$, in agreement with Proposition~\ref{prop:no-unitary}. The triple determines a family of subspaces of $W$ indexed by the mass shell, and a space of fields with values in them.

\begin{definition}[Fibres and fields]\label{def:local-fields}
	Let $(D,W,\iota)$ be as per Definition~\ref{def:local-realization}. The \emph{rest subspace} is $F(\mathring{p}):=\iota(V)$ and the \emph{fibre} at $p\in\mathcal{O}_m^+$ is $F(p):=D(A_p)F(\mathring{p})$, with $A_p$ as per Proposition~\ref{prop:orbit}. The \emph{space of fields} is
	\begin{equation*}
		\mathcal{F}_D:=\{\Psi:\mathcal{O}_m^+\to W\ \text{Borel}\mid\Psi(p)\in F(p)\ \ \forall p\in\mathcal{O}_m^+\},
	\end{equation*}
	and $\mathcal{P}$ acts on it by
	\begin{equation}\label{eq:local-action}
		\big(\mathcal{U}_D(a,g)\Psi\big)(p):=e^{ia\cdot p/\hbar}\,D(g)\,\Psi\big(\Lambda(g)^{-1}p\big).
	\end{equation}
\end{definition}

The matrix $D(g)$ in \eqref{eq:local-action} does not depend on $p$, whereas $\pi(R_W(g,p))$ in Definition~\ref{def:induced-rep} does. The following lemma shows that the price of this simplification is the constraint $\Psi(p)\in F(p)$, and that the two descriptions are equivalent.

\begin{lemma}[Local realization of the induced representation]\label{lem:local-realization}
	Let $(D,W,\iota)$ be a local covariant realization of $(\pi,V)$.
	\begin{enumerate}
		\item If $A'\in SL(2,\mathbb{C})$ satisfies $\Lambda(A')\mathring{p}=p$, then $D(A')F(\mathring{p})=F(p)$. Moreover $\dim F(p)=\dim V$.
		\item $D(g)F(p)=F(\Lambda(g)p)$ for all $g\in SL(2,\mathbb{C})$ and $p\in\mathcal{O}_m^+$.
		\item The map $(J\psi)(p):=D(A_p)\,\iota\,\psi(p)$ is a linear bijection from the Borel functions $\mathcal{O}_m^+\to V$ onto $\mathcal{F}_D$. The space $\mathcal{F}_D$ is invariant under $\mathcal{U}_D$, which is a representation of $\mathcal{P}$, and
		\begin{equation}\label{eq:J-intertwines}
			J\,U(a,g)=\mathcal{U}_D(a,g)\,J,
		\end{equation}
		where $U(a,g)$ acts on Borel functions by the formula of Definition~\ref{def:induced-rep}.
	\end{enumerate}
\end{lemma}

\begin{proof}
	Item 1. By item 3 of Proposition~\ref{prop:orbit}, $A'=A_pu$ with $u\in SU(2)$. By \eqref{eq:iota-intertwiner}, $D(u)F(\mathring{p})=\iota\pi(u)V=F(\mathring{p})$, so that $D(A')F(\mathring{p})=D(A_p)F(\mathring{p})=F(p)$. The dimension statement holds because $\iota$ is injective and $D(A_p)$ is invertible.
	
	Item 2. Let $q:=\Lambda(g)p$. Since $\Lambda(gA_p)\mathring{p}=q$, item 3 of Proposition~\ref{prop:orbit} yields $gA_p=A_qu$ with $u\in SU(2)$. It descends that $D(g)F(p)=D(A_q)D(u)F(\mathring{p})=F(q)$.
	
	Item 3. The map $J$ is linear and Borel because $p\mapsto A_p$ is smooth. It takes values in $\mathcal{F}_D$, and its inverse is $\psi(p)=\iota^{-1}D(A_p)^{-1}\Psi(p)$. Invariance of $\mathcal{F}_D$ follows from item 2, since $D(g)F(\Lambda(g)^{-1}p)=F(p)$. The homomorphism property of $\mathcal{U}_D$ is the identity $a'\cdot\Lambda(g)^{-1}p=\Lambda(g)a'\cdot p$ together with $D(g)D(g')=D(gg')$. For \eqref{eq:J-intertwines}, let $q:=\Lambda(g)^{-1}p$. By Proposition~\ref{prop:induced-action}, $gA_q=A_pR_W(g,p)$ with $R_W(g,p)\in SU(2)$, and \eqref{eq:iota-intertwiner} yields
	\begin{align*}
		\big(\mathcal{U}_D(a,g)J\psi\big)(p)&=e^{ia\cdot p/\hbar}D(g)D(A_q)\,\iota\,\psi(q)=e^{ia\cdot p/\hbar}D(A_p)D\big(R_W(g,p)\big)\iota\,\psi(q)\\
		&=e^{ia\cdot p/\hbar}D(A_p)\,\iota\,\pi\big(R_W(g,p)\big)\psi(q)=\big(J\,U(a,g)\psi\big)(p).\qedhere
	\end{align*}
\end{proof}

\noindent Two local realizations which differ by a choice of basis of $W$ describe the same states, and the next definition makes this precise.

\begin{definition}[Equivalent local realizations]\label{def:equivalent-local-realizations}
	Two local covariant realizations $(D,W,\iota)$ and $(D',W',\iota')$ of $(\pi,V)$ are \emph{equivalent} if there is a linear bijection $S:W\to W'$ with $S\,D(g)=D'(g)\,S$ for all $g\in SL(2,\mathbb{C})$ and $S\,\iota=\iota'$.
\end{definition}

The lemma has a geometric content, which is the form in which a field theory describes states, and which is stated in terms of the associated bundle of Appendix~\ref{app:induced-geometry}.

\begin{remark}[Local realizations as subbundles]\label{rem:local-subbundle}
	Let $E_\pi:=SL(2,\mathbb{C})\times_{SU(2)}V$ be the associated bundle of Definition~\ref{def:associated-bundle}. The assignment
	\begin{equation*}
		\widehat{J}:E_\pi\to\mathcal{O}_m^+\times W,\qquad[g,v]\mapsto\big(\Lambda(g)\mathring{p},\,D(g)\,\iota v\big),
	\end{equation*}
	is well defined, since $D(gu)\iota v=D(g)\iota\pi(u)v$ by \eqref{eq:iota-intertwiner}. It is linear and injective on each fibre, and its image is the subbundle of the trivial bundle with fibre $W$ whose fibre at $p$ is $F(p)$. No section of $SL(2,\mathbb{C})\to\mathcal{O}_m^+$ enters its definition. By Proposition~\ref{prop:action-equivariant}, $SL(2,\mathbb{C})$ acts on the sections of $E_\pi$ through $g\cdot[g',v]=[gg',v]$, and $\widehat{J}$ carries this action to $(p,w)\mapsto(\Lambda(g)p,D(g)w)$, in which the fibre is acted upon by a fixed matrix. The map $J$ of Lemma~\ref{lem:local-realization} is the expression of $\widehat{J}$ in the trivialization $s_0(p)=A_p$ of Definition~\ref{def:section-principal}, \textit{i.e.}, $\widehat{J}[A_p,\psi(p)]=(p,D(A_p)\iota\psi(p))$.
	
	The dependence on the momentum has not disappeared. In the trivialization of $E_\pi$ it sits in the transition function $h(g,p)=R_W(g,p)^{-1}$ of Definition~\ref{def:transition-function}, which enters the action by Proposition~\ref{prop:action-trivialization}, and Proposition~\ref{prop:wigner-no-gauge} states that no change of section removes it. In $\mathcal{O}_m^+\times W$ the transition function is the constant matrix $D(g)$, and the dependence is transferred to the position of the fibre $F(p)$ inside $W$, which moves with $p$ according to item 2 of Lemma~\ref{lem:local-realization}. A state is therefore no longer a function with values in a fixed spin space, but a function with values in $W$ constrained to a moving subspace.
	
	Conversely, the subbundle is determined by its fibre at the rest momentum. Every $SU(2)$-invariant subspace $L\subseteq W$ isomorphic to $V$ defines $F_L(p):=D(A_p)L$, and the proof of item 2 of the lemma uses only the $SU(2)$-invariance of $L$. Hence $F_L$ is an $SL(2,\mathbb{C})$-invariant subbundle with $D(g)F_L(p)=F_L(\Lambda(g)p)$. The fibres of a local realization are thus fixed by the pair formed by a finite-dimensional representation $D$ and an $SU(2)$-invariant subspace of $W$ isomorphic to $V$. The condition $\Psi(p)\in F(p)$ is a linear constraint on the values of the field, covariant by item 2, \textit{i.e.}, a covariant field equation in momentum space. Once $(D,W)$ is fixed, the only datum left is $F(\mathring{p})$, which Proposition~\ref{prop:rest-selection} determines for the Dirac space.
\end{remark}

\begin{remark}[Hilbert structure and realizations]\label{rem:transported-norm}
	Let $\Phi$ be an elementary relativistic system of mass $m>0$ and spin $j$, let $\mathcal{W}:\mathcal{H}\to\mathcal{H}_{m,j}$ be the unitary operator of Theorem~\ref{thm:wigner-mackey} and let $(D,W,\iota)$ be a local covariant realization of $(d_j,\mathbb{C}^{2j+1})$. On $\mathcal{F}_D^2:=J(\mathcal{H}_{m,j})$ we set $\langle J\psi,J\varphi\rangle_D:=\langle\psi,\varphi\rangle_{\mathcal{H}_{m,j}}$. Then $J$ is unitary and $\mathcal{U}_D$ acts unitarily on $\mathcal{F}_D^2$ by \eqref{eq:J-intertwines}. The unitary operator $J\mathcal{W}:\mathcal{H}\to\mathcal{F}_D^2$ induces a Kähler isomorphism $\mathbb{P}\mathcal{H}\to\mathbb{P}\mathcal{F}_D^2$ by Lemma~\ref{lem:kahler-iso-two-spaces}, which intertwines $\Phi_\kappa$ with $[\mathcal{U}_D(\kappa)\,\cdot\,]$. Every local covariant realization therefore yields a realization of $(\mathbb{P}\mathcal{H},\Phi)$ as per Definition~\ref{def:realization}. If $(D',W',\iota')$ is equivalent to $(D,W,\iota)$ through $S$, then $J'=SJ$, so that $S$ is unitary for the transported structures and the two realizations are isomorphic.
	
	The scalar product is defined here through $J$, and the Euclidean norm of $W$ plays no role. On a field it reads $\int\|\iota^{-1}D(A_p)^{-1}\Psi(p)\|^2\,d\mu(p)$ and depends on $p$ through $A_p$. Whether it is expressed by a fixed sesquilinear form on $W$ is a further property, which holds for the Dirac space by Theorem~\ref{thm:covariant-norm}.
\end{remark}

\begin{remark}[Locality in position space]\label{rem:position-space}
	For a field $\Psi\in\mathcal{F}_D$ which is continuous with compact support, set $\Psi(x):=\int e^{-ip\cdot x/\hbar}\Psi(p)\,d\mu(p)$ for $x\in M$. The substitution $p=\Lambda(g)q$ and the invariance of $\mu$ by Lemma~\ref{lem:invariant-measure} yield, with $\Lambda(g)^{-1}$ acting on $a\cdot p$ and $x\cdot p$,
	\begin{equation*}
		\big(\mathcal{U}_D(a,g)\Psi\big)(x)=D(g)\,\Psi\big(\Lambda(g)^{-1}(x-a)\big).
	\end{equation*}
	The value at $x$ is mapped to the value at the transformed point by a fixed matrix, which is the definition of a local transformation law. For the induced representation the same substitution produces the integral
	\begin{equation*}
		\int e^{-ip\cdot x/\hbar}\,\pi\big(R_W(g,p)\big)\,\psi\big(\Lambda(g)^{-1}p\big)\,d\mu(p),
	\end{equation*}
	in which the matrix depends on the integration variable and cannot be taken out. By Proposition~\ref{prop:wigner-rotation} it is not constant in $p$, and by Proposition~\ref{prop:wigner-no-gauge} this does not depend on the choice of the section $p\mapsto A_p$. The field $\Psi(x)$ is a Fourier transform of a state and does not define a position observable.
\end{remark}

%--------------------------------------------------------------------
\subsection{The Obstruction to Unitary Locality}
\label{subsec:no-unitary-locality}

The matrices $D(g)$ of a local realization are not required to be unitary, and one asks whether this freedom is necessary. The data are the induced representations $U_{m,j}$, the central sign of Corollary~\ref{cor:central-sign-spin} and the absence of nontrivial unitary finite-dimensional representations of $SL(2,\mathbb{C})$ of Proposition~\ref{prop:no-unitary}. The result is that a local model with unitary matrices exists only for $j=0$. It explains why the constraint $\Psi(p)\in F(p)$ and a norm different from the Euclidean one of $W$ are unavoidable, and it prepares the introduction of parity in the next subsection.

\begin{proposition}[Obstruction to a unitary local realization]\label{prop:no-local-unitary}
	Let $V'$ be a nonzero finite-dimensional Hilbert space and let $M:SL(2,\mathbb{C})\to GL(V')$ be a continuous homomorphism such that
	\begin{equation*}
		\big(\mathcal{U}_M(a,g)\psi\big)(p):=e^{ia\cdot p/\hbar}\,M(g)\,\psi\big(\Lambda(g)^{-1}p\big)
	\end{equation*}
	is unitary on $L^2(\mathcal{O}_m^+,\mu;V')$ for all $(a,g)\in\mathcal{P}$. Then $M$ is trivial. Consequently, for $j\neq0$ the representation $U_{m,j}$ is not unitarily equivalent to any representation of this form.
\end{proposition}

\begin{proof}
	Let $v\in V'$ and let $B\subseteq\mathcal{O}_m^+$ be a Borel set with $0<\mu(B)<\infty$, and set $\psi=\chi_Bv$. By Lemma~\ref{lem:invariant-measure}, $\|\mathcal{U}_M(0,g)\psi\|^2=\mu(B)\|M(g)v\|^2$, which equals $\|\psi\|^2=\mu(B)\|v\|^2$. Hence $M(g)$ is an isometry of a finite-dimensional space, \textit{i.e.}, a unitary matrix, and Proposition~\ref{prop:no-unitary} yields $M=\mathbb{1}$.
	
	Let $\mathcal{V}$ be a unitary operator with $\mathcal{V}U_{m,j}\mathcal{V}^{-1}=\mathcal{U}_M$. Then $M=\mathbb{1}$ and $\mathcal{U}_M(a,g)\psi(p)=e^{ia\cdot p/\hbar}\psi(\Lambda(g)^{-1}p)$. If $j$ is a half-integer, Corollary~\ref{cor:central-sign-spin} yields $U_{m,j}(0,-\mathbb{1})=-\mathbb{1}$, whereas $\mathcal{U}_M(0,-\mathbb{1})=\mathbb{1}$, which is a contradiction. If $j\geq1$ is an integer, then $\mathcal{U}_M$ is reducible for $\dim V'\geq2$, since $L^2(\mathcal{O}_m^+,\mu;\mathbb{C}v)$ is a proper invariant subspace, whereas $U_{m,j}$ is irreducible by item 1 of Theorem~\ref{thm:wigner-mackey}. For $\dim V'=1$ one has $\mathcal{U}_M=U_{m,0}$, and the uniqueness of the label in item 2 of Theorem~\ref{thm:wigner-mackey}, applied to $U_{m,j}$, forces $j=0$.
\end{proof}

\noindent The proposition is the precise form of the statement that spin and momentum cannot be separated by a unitary local model.

\begin{remark}
	For $j=0$ the obstruction is absent, since $W=\mathbb{C}$, $D=\mathbb{1}$ and $\iota=\mathrm{id}$ is a local covariant realization with $F(p)=W$ and unitary $D$. The obstruction is therefore a property of the spin. For $j\neq0$ it dictates the first move of the construction. The matrices $D(g)$ are not unitary, the Euclidean norm of $W$ is not the norm of the representation, and the norm is carried by the scalar product transported in Remark~\ref{rem:transported-norm}.
\end{remark}

%--------------------------------------------------------------------
\subsection{Space Inversion and the Hypotheses of Locality}
\label{subsec:parity-hypotheses}

Space inversion does not enter H1 to H3, and the induced model carries an operator implementing it for each sign $\epsilon=\pm1$. The question is whether the local model can implement it by a fixed matrix. The data are the automorphism $\theta$ and the properties of $A_p$ of Lemma~\ref{lem:inversion}. This subsection defines the parity matrix, shows that the induced representation realizes inversion without Wigner rotations, and gives the condition under which the space of fields is stable. It then collects the requirements on the model as H4 to H6. The next subsection shows that these requirements select the Dirac space.

\begin{definition}[Parity matrix]\label{def:parity-realization}
	Let $(D,W,\iota)$ be a local covariant realization. A \emph{parity matrix} for $D$ is a matrix $M\in GL(W)$ with $M^2=\mathbb{1}$ and
	\begin{equation}\label{eq:parity-matrix}
		M\,D(g)\,M^{-1}=D(\theta(g)),\qquad g\in SL(2,\mathbb{C}),
	\end{equation}
	with $\theta(g)=(g^\dagger)^{-1}$ as per Lemma~\ref{lem:inversion}.
\end{definition}

\begin{remark}[Space inversion on fields]\label{rem:parity-action}
	Space inversion acts on fields by $(\mathcal{U}_D(\mathsf{P})\Psi)(p):=M\Psi(\mathsf{P}p)$. Since $\mathsf{P}a\cdot p=a\cdot\mathsf{P}p$ and $\mathsf{P}\Lambda(g)^{-1}\mathsf{P}=\Lambda(\theta(g))^{-1}$ by item 1 of Lemma~\ref{lem:inversion}, and $\mathcal{U}_D(\mathsf{P})^{-1}=\mathcal{U}_D(\mathsf{P})$, Equation~\eqref{eq:parity-matrix} yields
	\begin{equation*}
		\mathcal{U}_D(\mathsf{P})\,\mathcal{U}_D(a,g)\,\mathcal{U}_D(\mathsf{P})^{-1}=\mathcal{U}_D(\mathsf{P}a,\theta(g)).
	\end{equation*}
	The condition $M^2=\mathbb{1}$ corresponds to the extension of $\mathcal{P}$ by an involution acting by $(a,g)\mapsto(\mathsf{P}a,\theta(g))$.
\end{remark}

The model has to reproduce a property of the induced representation, which is that inversion acts without Wigner rotations. The following lemma verifies it for every $j$.

\begin{lemma}[Parity on the induced representation]\label{lem:parity-induced}
	Let $\epsilon\in\{+1,-1\}$ and let $(U^\epsilon(\mathsf{P})\psi)(p):=\epsilon\,\psi(\mathsf{P}p)$ on $\mathcal{H}_{m,j}$.
	\begin{enumerate}
		\item $R_W(\theta(g),p)=R_W(g,\mathsf{P}p)$ for all $g\in SL(2,\mathbb{C})$ and $p\in\mathcal{O}_m^+$.
		\item $U^\epsilon(\mathsf{P})$ is unitary, $U^\epsilon(\mathsf{P})^2=\mathbb{1}$ and
		\begin{equation}\label{eq:parity-covariance-induced}
			U^\epsilon(\mathsf{P})\,U(a,g)\,U^\epsilon(\mathsf{P})^{-1}=U(\mathsf{P}a,\theta(g)).
		\end{equation}
	\end{enumerate}
\end{lemma}

\begin{proof}
	Item 1. By item 1 of Lemma~\ref{lem:inversion}, $\Lambda(\theta(g))^{-1}=\mathsf{P}\Lambda(g)^{-1}\mathsf{P}$, and by item 2 of the same lemma $A_p=\theta(A_{\mathsf{P}p})$ and $A_{\mathsf{P}\Lambda(g)^{-1}\mathsf{P}p}=\theta(A_{\Lambda(g)^{-1}\mathsf{P}p})$. Since $\theta$ is an automorphism,
	\begin{equation*}
		R_W(\theta(g),p)=\theta(A_{\mathsf{P}p})^{-1}\,\theta(g)\,\theta\big(A_{\Lambda(g)^{-1}\mathsf{P}p}\big)=\theta\big(R_W(g,\mathsf{P}p)\big)=R_W(g,\mathsf{P}p),
	\end{equation*}
	where the last equality holds because $R_W(g,\mathsf{P}p)\in SU(2)$ and $\theta|_{SU(2)}=\mathrm{id}$.
	
	Item 2. Unitarity follows from item 3 of Lemma~\ref{lem:inversion}, and $U^\epsilon(\mathsf{P})^2=\mathbb{1}$ from $\mathsf{P}^2=\mathrm{id}$ and $\epsilon^2=1$. With $U^\epsilon(\mathsf{P})^{-1}=U^\epsilon(\mathsf{P})$, item 1 and $\mathsf{P}a\cdot p=a\cdot\mathsf{P}p$ yield
	\begin{align*}
		\big(U^\epsilon(\mathsf{P})U(a,g)U^\epsilon(\mathsf{P})^{-1}\psi\big)(p)&=\epsilon^2e^{ia\cdot\mathsf{P}p/\hbar}\,\pi\big(R_W(g,\mathsf{P}p)\big)\,\psi\big(\mathsf{P}\Lambda(g)^{-1}\mathsf{P}p\big)\\
		&=e^{i\mathsf{P}a\cdot p/\hbar}\,\pi\big(R_W(\theta(g),p)\big)\,\psi\big(\Lambda(\theta(g))^{-1}p\big),
	\end{align*}
	which is $(U(\mathsf{P}a,\theta(g))\psi)(p)$.
\end{proof}

\noindent The two signs $\epsilon=\pm1$ are the two intrinsic parities, and the group relations do not distinguish them. The next lemma gives the condition under which the space of fields of a local realization is stable under inversion, and shows that the sign $\epsilon$ is then determined by $M$.

\begin{lemma}[Parity stability]\label{lem:parity-stable}
	Let $(D,W,\iota)$ be a local covariant realization of $(d_j,\mathbb{C}^{2j+1})$ with a parity matrix $M$.
	\begin{enumerate}
		\item $MF(p)=D(A_{\mathsf{P}p})\,MF(\mathring{p})$ for all $p\in\mathcal{O}_m^+$.
		\item $\mathcal{U}_D(\mathsf{P})$ maps $\mathcal{F}_D$ into itself if and only if $MF(\mathring{p})=F(\mathring{p})$. In this case $MF(p)=F(\mathsf{P}p)$ for all $p$.
		\item If $MF(\mathring{p})=F(\mathring{p})$, there is $\epsilon\in\{+1,-1\}$ with $M\iota=\epsilon\,\iota$, and $J\,U^\epsilon(\mathsf{P})=\mathcal{U}_D(\mathsf{P})\,J$.
	\end{enumerate}
\end{lemma}

\begin{proof}
	Item 1. By \eqref{eq:parity-matrix} and item 2 of Lemma~\ref{lem:inversion}, $MD(A_p)=D(\theta(A_p))M=D(A_{\mathsf{P}p})M$.
	
	Item 2. Let $MF(\mathring{p})=F(\mathring{p})$. Item 1 yields $MF(p)=D(A_{\mathsf{P}p})F(\mathring{p})=F(\mathsf{P}p)$. Hence $M\Psi(\mathsf{P}p)\in MF(\mathsf{P}p)=F(p)$ for every $\Psi\in\mathcal{F}_D$. Conversely, let $v\in F(\mathring{p})$ and let $\Psi_v(p):=D(A_p)v$, which belongs to $\mathcal{F}_D$. Since $\mathsf{P}\mathring{p}=\mathring{p}$, one has $(\mathcal{U}_D(\mathsf{P})\Psi_v)(\mathring{p})=Mv\in F(\mathring{p})$. Hence $MF(\mathring{p})\subseteq F(\mathring{p})$, with equality because $M$ is invertible.
	
	Item 3. Let $T:=\iota^{-1}M\iota:V\to V$, which is defined because $MF(\mathring{p})=F(\mathring{p})$. By \eqref{eq:iota-intertwiner} and \eqref{eq:parity-matrix}, $T\pi(u)=\iota^{-1}MD(u)\iota=\iota^{-1}D(\theta(u))M\iota=\pi(u)T$, since $\theta|_{SU(2)}=\mathrm{id}$. By Lemma 1.2.26, $T=\epsilon\mathbb{1}$, and $T^2=\mathbb{1}$ yields $\epsilon=\pm1$. By item 1 applied with $\theta(A_{\mathsf{P}p})=A_p$,
	\begin{equation*}
		(\mathcal{U}_D(\mathsf{P})J\psi)(p)=MD(A_{\mathsf{P}p})\,\iota\,\psi(\mathsf{P}p)=D(A_p)\,M\iota\,\psi(\mathsf{P}p)=D(A_p)\,\iota\,\epsilon\psi(\mathsf{P}p)=(J\,U^\epsilon(\mathsf{P})\psi)(p).\qedhere
	\end{equation*}
\end{proof}

\noindent It is now possible to state the requirements on the model of the states. They are conditions on a triple and not on $\Phi$.

\begin{definition}[Locality hypotheses]\label{def:locality-hypotheses}
	Let $\Phi$ be an elementary relativistic system of mass $m>0$ and spin $j$ as per Theorem~\ref{thm:wigner-mackey}. The hypotheses H4, H5 and H6 of Definition~\ref{def:hypotheses} on a triple $(D,W,\iota)$ read as follows.
	\begin{enumerate}
		\item[\textup{H4}] \emph{Locality.} $(D,W,\iota)$ is a local covariant realization of $(d_j,\mathbb{C}^{2j+1})$ as per Definition~\ref{def:local-realization}.
		\item[\textup{H5}] \emph{Parity.} There is a parity matrix $M$ for $D$ as per Definition~\ref{def:parity-realization} with $MF(\mathring{p})=F(\mathring{p})$.
		\item[\textup{H6}] \emph{Minimality.} $\dim W\leq\dim W'$ for every triple $(D',W',\iota')$ which satisfies H4 and H5.
	\end{enumerate}
\end{definition}

\begin{remark}[Scope of H4 to H6]\label{rem:locality-hypotheses}
	The hypotheses concern the model of the states. By Remark~\ref{rem:transported-norm} every triple satisfying H4 yields a realization of $(\mathbb{P}\mathcal{H},\Phi)$, and all of them describe the same states. Hypothesis H4 alone is not selective. For every $j$ the representation $D^{(j,0)}$ of Theorem~\ref{thm:finite-reps} restricts to $d_j$ on $SU(2)$, so that $(D^{(j,0)},\mathbb{C}^{2j+1},\mathrm{id})$ satisfies H4 with $F(p)=\mathbb{C}^{2j+1}$ for every $p$ and with no constraint on the fields. The selective content is in H5 and H6, and the next subsection shows that for $j=\frac12$ they select a single model. Space inversion is not part of H1 to H3, and H5 does not assert that it is a symmetry of nature, but only that the model is compatible with it.
\end{remark}

%--------------------------------------------------------------------
\subsection{The Dirac Space}
\label{subsec:dirac-space}

Hypothesis H5 requires a matrix exchanging $D$ with $D\circ\theta$, and the Weyl representations of Definition~\ref{def:weyl-spinors} are exchanged by $\theta$. The question is which local realizations of spin $\frac12$ satisfy H5, and which is the smallest. From now on $j=\frac12$, so that $V=\mathbb{C}^2$ and $\pi(u)=u$. The data are Theorem~\ref{thm:finite-reps} and Lemma~\ref{lem:parity-stable}. The results are that neither Weyl representation supports a parity matrix, that on the sum of the two the parity matrix is unique up to equivalence, and that the smallest parity-covariant space is the Dirac space of dimension four. The subspace $F(\mathring{p})$ and the transport along the mass shell are the object of the next section.

\begin{proposition}[The parity problem]\label{prop:parity-problem}
	\begin{enumerate}
		\item There is no $T\in GL(2,\mathbb{C})$ with $T\,g\,T^{-1}=(g^\dagger)^{-1}$ for all $g\in SL(2,\mathbb{C})$. In particular $D^{(1/2,0)}$ and $D^{(0,1/2)}$ are not equivalent.
		\item Neither $D^{(1/2,0)}$ nor $D^{(0,1/2)}$ admits a parity matrix.
	\end{enumerate}
\end{proposition}

\begin{proof}
	Item 1. Suppose that $T$ exists. For $u\in SU(2)$ one has $(u^\dagger)^{-1}=u$, so that $T$ commutes with $SU(2)$ and $T=c\,\mathbb{1}$ by Lemma 1.2.26, since the defining representation of $SU(2)$ is irreducible. Then $g=(g^\dagger)^{-1}$ for all $g$, which fails for $g=\operatorname{diag}(2,\tfrac12)$.
	
	Item 2. A parity matrix for $D^{(1/2,0)}$ would satisfy $MgM^{-1}=\theta(g)$, which contradicts item 1. For $D^{(0,1/2)}$, Equation~\eqref{eq:parity-matrix} reads $M\theta(g)M^{-1}=g$. Replacing $g$ by $\theta(g)$ and using $\theta^2=\mathrm{id}$ yields $Mg'M^{-1}=\theta(g')$ for all $g'$, which is the same contradiction.
\end{proof}

The failure is a failure of equivalence between $D$ and $D\circ\theta$, and a parity-covariant realization has to contain both chiralities. On the sum of the two Weyl representations a parity matrix exists, and the following lemma shows that it is unique up to a rescaling.

\begin{lemma}[Uniqueness of the parity matrix]\label{lem:parity-matrix-unique}
	Let $D:=D^{(1/2,0)}\oplus D^{(0,1/2)}$, \textit{i.e.}, $D(g)=\operatorname{diag}(g,\theta(g))$ on $\mathbb{C}^4=\mathbb{C}^2\oplus\mathbb{C}^2$, and let $\gamma^0:=\begin{pmatrix}0&\mathbb{1}\\\mathbb{1}&0\end{pmatrix}$.
	\begin{enumerate}
		\item The commutant of $D(SL(2,\mathbb{C}))$ consists of the matrices $\operatorname{diag}(s_1\mathbb{1},s_2\mathbb{1})$ with $s_1,s_2\in\mathbb{C}$.
		\item The parity matrices for $D$ are exactly $M_b:=\begin{pmatrix}0&b\mathbb{1}\\b^{-1}\mathbb{1}&0\end{pmatrix}$ with $b\in\mathbb{C}\setminus\{0\}$.
		\item $M_b=S_b\,\gamma^0S_b^{-1}$ with $S_b:=\operatorname{diag}(\mathbb{1},b^{-1}\mathbb{1})$, which belongs to the commutant. Consequently, if $(D,\mathbb{C}^4,\iota)$ is a local covariant realization with parity matrix $M_b$, then $(D,\mathbb{C}^4,S_b^{-1}\iota)$ is equivalent to it through $S_b^{-1}$ as per Definition~\ref{def:equivalent-local-realizations}, and $\gamma^0=S_b^{-1}M_bS_b$ is a parity matrix for it. The relation $M F(\mathring{p})=F(\mathring{p})$ is preserved.
	\end{enumerate}
\end{lemma}

\begin{proof}
	Both Weyl representations are irreducible, because their restrictions to $SU(2)$ are the defining representation, and they are not equivalent by item 1 of Proposition~\ref{prop:parity-problem}. Write a matrix as $X=(X_{kl})_{k,l=1,2}$ with $2\times2$ blocks.
	
	Item 1. The relation $XD(g)=D(g)X$ reads $X_{11}g=gX_{11}$, $X_{22}\theta(g)=\theta(g)X_{22}$, $X_{12}\theta(g)=gX_{12}$ and $X_{21}g=\theta(g)X_{21}$. By Lemma 1.2.26, item 1 for the diagonal blocks, $X_{11}$ and $X_{22}$ are scalar. By item 2 of the same lemma the intertwiners $X_{12}$ and $X_{21}$ between inequivalent irreducible representations vanish.
	
	Item 2. Since $D(\theta(g))=\operatorname{diag}(\theta(g),g)$, Equation~\eqref{eq:parity-matrix} reads $MD(g)=D(\theta(g))M$, which yields $M_{11}g=\theta(g)M_{11}$, $M_{22}\theta(g)=gM_{22}$, $M_{12}\theta(g)=\theta(g)M_{12}$ and $M_{21}g=gM_{21}$. By Lemma 1.2.26, $M_{11}$ and $M_{22}$ vanish as intertwiners between inequivalent irreducible representations, $M_{12}=b\mathbb{1}$ and $M_{21}=c\mathbb{1}$. It descends that $M^2=bc\,\mathbb{1}$, so that $M^2=\mathbb{1}$ yields $c=b^{-1}$. Conversely, a direct computation shows that $M_b$ satisfies $M_b^2=\mathbb{1}$ and \eqref{eq:parity-matrix}.
	
	Item 3. The identity $M_b=S_b\gamma^0S_b^{-1}$ is the multiplication of block-diagonal matrices. The matrix $S_b$ is in the commutant by item 1, so that $S_b^{-1}D(g)S_b=D(g)$ and $S_b^{-1}M_bS_b=\gamma^0$ satisfies \eqref{eq:parity-matrix} for $D$. Moreover $\gamma^0S_b^{-1}F(\mathring{p})=S_b^{-1}M_bF(\mathring{p})=S_b^{-1}F(\mathring{p})$ whenever $M_bF(\mathring{p})=F(\mathring{p})$.
\end{proof}

\noindent The lemma fixes the parity matrix once $D$ is fixed. The next proposition shows that $D$ is forced as well, once the dimension is minimal.

\begin{proposition}[Minimal parity-covariant space]\label{prop:minimal-extension}
	Let $(D,W,\iota)$ be a local covariant realization of $(d_{1/2},\mathbb{C}^2)$ with a parity matrix $M$. Then $\dim W\geq4$, and $\dim W=4$ if and only if $D\cong D^{(1/2,0)}\oplus D^{(0,1/2)}$.
\end{proposition}

\begin{proof}
	By Theorem~\ref{thm:finite-reps}, $D\cong\bigoplus_{(j_1,j_2)}n(j_1,j_2)\,D^{(j_1,j_2)}$ with multiplicities $n(j_1,j_2)\in\mathbb{N}$. Equation~\eqref{eq:parity-matrix} states that $M$ is an equivalence between $D$ and $D\circ\theta\cong\bigoplus n(j_1,j_2)\,D^{(j_2,j_1)}$, and the isotypic decomposition is unique, so that $n(j_1,j_2)=n(j_2,j_1)$. The subspace $\iota(V)$ is an irreducible $SU(2)$-invariant subspace of $W$ isomorphic to $d_{1/2}$. Some projection onto a summand with $n(j_1,j_2)>0$ is nonzero on it, hence injective by Lemma 1.2.26, and $d_{1/2}$ is contained in $D^{(j_1,j_2)}|_{SU(2)}\cong d_{j_1}\otimes d_{j_2}$. By the Clebsch--Gordan decomposition \cite[TODO-location]{TODO-Hal03},
	\begin{equation*}
		d_{j_1}\otimes d_{j_2}\cong\bigoplus_{J=|j_1-j_2|}^{j_1+j_2}d_J.
	\end{equation*}
	This contains $d_{1/2}$ if and only if $|j_1-j_2|\leq\tfrac12\leq j_1+j_2$ and $\tfrac12-|j_1-j_2|\in\mathbb{Z}$, \textit{i.e.}, $|j_1-j_2|=\tfrac12$. In particular $j_1\neq j_2$, so that $(j_2,j_1)$ occurs as well, with the same multiplicity, and
	\begin{equation*}
		\dim W\geq2(2j_1+1)(2j_2+1)\geq4,
	\end{equation*}
	since $|j_1-j_2|=\tfrac12$ implies $(2j_1+1)(2j_2+1)\geq2$. Equality holds if and only if $\{j_1,j_2\}=\{0,\tfrac12\}$ and $W$ contains no other summand.
\end{proof}

\noindent The minimal space is now given a name. It is the space on which the matrix $\gamma^0$ of the next sections acts.

\begin{definition}[Dirac space]\label{def:dirac-space}
	The \emph{Dirac space} is $W_{\mathrm{D}}:=\mathbb{C}^4=\mathbb{C}^2\oplus\mathbb{C}^2$ with the representation and the matrix
	\begin{equation*}
		D(g):=\begin{pmatrix}g&0\\0&(g^\dagger)^{-1}\end{pmatrix},\qquad\gamma^0:=\begin{pmatrix}0&\mathbb{1}\\\mathbb{1}&0\end{pmatrix},
	\end{equation*}
	\textit{i.e.}, $D=D^{(1/2,0)}\oplus D^{(0,1/2)}$. We write $M:=\gamma^0$.
\end{definition}

\begin{lemma}\label{lem:dirac-space}
	\begin{enumerate}
		\item $(\gamma^0)^2=\mathbb{1}$, $\gamma^0$ is Hermitian and $\gamma^0D(g)\gamma^0=D(\theta(g))$, so that $\gamma^0$ is a parity matrix for $D$.
		\item $D(g)^\dagger\gamma^0D(g)=\gamma^0$ for all $g\in SL(2,\mathbb{C})$.
		\item $D(u)=\operatorname{diag}(u,u)$ is unitary for $u\in SU(2)$, and $D(A_p)=\operatorname{diag}(A_p,A_p^{-1})$.
		\item For $\epsilon\in\{+1,-1\}$ the map $\iota_\epsilon:\mathbb{C}^2\to\mathbb{C}^4$, $\iota_\epsilon\eta:=(\eta,\epsilon\eta)$, makes $(D,W_{\mathrm{D}},\iota_\epsilon)$ a local covariant realization of $(d_{1/2},\mathbb{C}^2)$ with $\gamma^0\iota_\epsilon=\epsilon\,\iota_\epsilon$ and $\gamma^0F(\mathring{p})=F(\mathring{p})$.
	\end{enumerate}
\end{lemma}

\begin{proof}
	Item 1. $D$ is a homomorphism because $\theta$ is, and $\gamma^0D(g)\gamma^0=\operatorname{diag}(\theta(g),g)=\operatorname{diag}(\theta(g),\theta(\theta(g)))=D(\theta(g))$ since $\theta^2=\mathrm{id}$. Item 2. Since $\theta(g)^\dagger=g^{-1}$,
	\begin{equation*}
		D(g)^\dagger\gamma^0D(g)=\begin{pmatrix}g^\dagger&0\\0&g^{-1}\end{pmatrix}\begin{pmatrix}0&\mathbb{1}\\\mathbb{1}&0\end{pmatrix}\begin{pmatrix}g&0\\0&(g^\dagger)^{-1}\end{pmatrix}=\begin{pmatrix}0&g^\dagger(g^\dagger)^{-1}\\g^{-1}g&0\end{pmatrix}=\gamma^0.
	\end{equation*}
	Item 3 follows from $\theta|_{SU(2)}=\mathrm{id}$ and $\theta(A_p)=A_p^{-1}$, item 2 of Lemma~\ref{lem:inversion}. Item 4. For $u\in SU(2)$, $D(u)\iota_\epsilon\eta=(u\eta,\epsilon u\eta)=\iota_\epsilon(\pi(u)\eta)$, and $\gamma^0(\eta,\epsilon\eta)=(\epsilon\eta,\eta)=\epsilon(\eta,\epsilon\eta)$. The injectivity of $\iota_\epsilon$ is clear.
\end{proof}

\begin{remark}[Hypothesis H6 is a choice]\label{rem:minimality-choice}
	By item 4 of Lemma~\ref{lem:dirac-space} the triple $(D,W_{\mathrm{D}},\iota_\epsilon)$ satisfies H4 and H5 with $\dim W_{\mathrm{D}}=4$. By Proposition~\ref{prop:minimal-extension}, H6 holds exactly for the triples with $\dim W=4$, \textit{i.e.}, with $D\cong D^{(1/2,0)}\oplus D^{(0,1/2)}$, and by Lemma~\ref{lem:parity-matrix-unique} the parity matrix is then equivalent to $\gamma^0$. What remains to be determined is the rest subspace $F(\mathring{p})$ inside $W_{\mathrm{D}}$, which Proposition~\ref{prop:rest-selection} fixes up to the sign $\epsilon$ of Lemma~\ref{lem:parity-stable}.
	
	Hypothesis H6 is a requirement of economy and not a consequence of H1 to H5. Larger spaces also satisfy H4 and H5. For instance, $D^{(1/2,1)}\oplus D^{(1,1/2)}$ contains the spin $\frac12$ representation once in each summand, and the eigenspaces of the exchange of the two copies provide $M$-invariant choices of $\iota(V)$. By Lemma~\ref{lem:local-realization} such realizations describe the same states as $\mathcal{H}_{m,1/2}$, with redundant components and further constraints. The restriction to $j=\frac12$ enters the proof of Proposition~\ref{prop:minimal-extension} through the exclusion of $j_1=j_2$, which is a consequence of half-integrality. For $j=1$ the representation $D^{(1/2,1/2)}$ is exchanged with itself by $\theta$ and is smaller than $D^{(1,0)}\oplus D^{(0,1)}$.
\end{remark}

\section{The Dirac Equation from the Rest Frame}
\label{sec:dirac-equation}

The states of an elementary relativistic system of mass $m>0$ and spin $\frac12$ are the rays of $\mathcal{H}_{m,1/2}$, by Theorem~\ref{thm:wigner-mackey}. This section describes the same states by fields with values in the Dirac space. A state of momentum $p\in\mathcal{O}_m^+$ is then a spinor lying in a subspace $F(p)\subseteq\mathbb{C}^4$, and the group acts through the fixed matrices $D(g)$. The task is to determine the subspaces $F(p)$ of admissible spinors. What is given is a local covariant realization as per Definition~\ref{def:local-realization}, equipped with a parity matrix as per Definition~\ref{def:parity-realization}, so that H4 to H6 of Definition~\ref{def:hypotheses} are in force. By Proposition~\ref{prop:minimal-extension} and Lemma~\ref{lem:parity-matrix-unique}, such a realization is equivalent to one on the Dirac space $W_{\mathrm{D}}$ of Definition~\ref{def:dirac-space} with parity matrix $M=\gamma^0$. This is the form used below, with $j=\frac12$, $V=\mathbb{C}^2$ and $\pi(u)=u$ as per Definition~\ref{def:spin-data}. The realization is fixed by one object, the subspace $F(\mathring{p})$ of the states at rest, because every other $F(p)$ is the image of $F(\mathring{p})$ under the standard boost. The section obtains $F(\mathring{p})$ from the two symmetries left at rest, rotations and space inversion, and transports it along the mass shell. The transported constraint is the Dirac equation, first in momentum space and then in position space. No form of the equation, no number of components and no gamma matrices are assumed, since the Dirac space is itself an output of Proposition~\ref{prop:minimal-extension}. This is the step of the chain that yields the equation of motion. The following section needs the space of fields obtained here to define the invariant norm and to return to the projective space of states.

%--------------------------------------------------------------------
\subsection{The Rest Subspace}\label{sec:rest}

The subspace $F(\mathring{p})$ is the space of the states of the system at rest, written as spinors in $\mathbb{C}^4$. It must carry exactly one copy of the spin $\frac12$ representation of $SU(2)$, as demanded by Definition~\ref{def:local-realization}, and it must be stable under the parity matrix. The first requirement alone leaves a complex projective line of candidates, and the second selects two of them. This subsection proves both statements and records that the two candidates are exchanged by a symmetry of the Dirac space, so that the choice reduces to the sign of the intrinsic parity. The outcome is the input of the transport carried out in the next subsection.

\begin{lemma}[Invariant subspaces of the stabilizer]\label{lem:su2-subspaces}
	Identify $\mathbb{C}^4=\mathbb{C}^2\oplus\mathbb{C}^2$ with $\mathbb{C}^2_{\mathrm{spin}}\otimes\mathbb{C}^2_{\mathrm{mult}}$ through $(\xi,\zeta)\mapsto\xi\otimes e_1+\zeta\otimes e_2$, where $e_1,e_2$ is the canonical basis of $\mathbb{C}^2_{\mathrm{mult}}$.
	\begin{enumerate}
		\item $D(u)=u\otimes\mathbb{1}$ for $u\in SU(2)$, and $\gamma^0=\mathbb{1}\otimes\sigma_1$.
		\item The subspaces of $\mathbb{C}^4$ invariant under $D(SU(2))$ are the subspaces $\mathbb{C}^2_{\mathrm{spin}}\otimes L$ with $L\subseteq\mathbb{C}^2_{\mathrm{mult}}$ a subspace. Such a subspace is isomorphic to the spin $\frac12$ representation if and only if $\dim L=1$, and these subspaces form a family parametrized by $L\in\mathbb{CP}^1$.
	\end{enumerate}
\end{lemma}

\begin{proof}
	Item 1. By Lemma~\ref{lem:dirac-space}, $D(u)(\xi,\zeta)=(u\xi,u\zeta)$, which corresponds to $u\xi\otimes e_1+u\zeta\otimes e_2$. Moreover $\gamma^0(\xi,\zeta)=(\zeta,\xi)$ corresponds to $\zeta\otimes e_1+\xi\otimes e_2=(\mathbb{1}\otimes\sigma_1)(\xi\otimes e_1+\zeta\otimes e_2)$.
	
	Item 2. By item 1 the subspaces $\mathbb{C}^2\otimes L$ are invariant. Conversely, let $K$ be invariant. Since $D|_{SU(2)}$ is unitary by Lemma~\ref{lem:dirac-space}, the orthogonal projector $Q$ onto $K$ commutes with every $u\otimes\mathbb{1}$. The matrices $e^{-i\theta\sigma_j/2}=\cos\frac\theta2\,\mathbb{1}-i\sin\frac\theta2\,\sigma_j$ belong to $SU(2)$ as in Example 1.2.17, and $\theta\in\{0,\pi\}$ yields $\mathbb{1}$ and $-i\sigma_j$. These span $M_2(\mathbb{C})$, hence $Q$ commutes with $B\otimes\mathbb{1}$ for every $B\in M_2(\mathbb{C})$. Writing $Q=\sum_{k,l}E_{kl}\otimes Q_{kl}$ with matrix units $E_{kl}$, the coefficient of $E_{ab}$ in $[E_{ij}\otimes\mathbb{1},Q]=0$ is $\delta_{ai}Q_{jb}-\delta_{bj}Q_{ai}$. It follows that $Q_{jb}=0$ for $b\neq j$ and $Q_{jj}=Q_{ii}$, \textit{i.e.}, $Q=\mathbb{1}\otimes Q_0$ with $Q_0$ a projector. Then $K=\mathbb{C}^2\otimes\operatorname{ran}Q_0$. Since $u$ acts as $u\otimes\mathbb{1}$, the subspace $\mathbb{C}^2\otimes L$ is a direct sum of $\dim L$ copies of the spin $\frac12$ representation.
\end{proof}

\noindent Hypotheses used: H6, through $W=W_{\mathrm{D}}$. The lemma shows that rotations alone cannot fix $F(\mathring{p})$, since the multiplicity factor $\mathbb{C}^2_{\mathrm{mult}}$ is not touched by $D(SU(2))$. Space inversion acts on this factor as $\sigma_1$ and has simple spectrum, which is why it reduces the projective line to two points.

\begin{proposition}[Selection of the rest subspace]\label{prop:rest-selection}
	Let $(D,W_{\mathrm{D}},\iota)$ be a local covariant realization of $(d_{1/2},\mathbb{C}^2)$ whose space of fields $\mathcal{F}_D$ is invariant under $\mathcal{U}_D(\mathsf{P})$ with parity matrix $M=\gamma^0$. Then there is $\epsilon\in\{+1,-1\}$ such that
	\begin{equation}\label{eq:rest-subspace}
		F(\mathring{p})=\{\varphi\in\mathbb{C}^4\mid\gamma^0\varphi=\epsilon\varphi\}=\{(\eta,\epsilon\eta)\mid\eta\in\mathbb{C}^2\},
	\end{equation}
	$\iota$ is a multiple of $\iota_\epsilon:\eta\mapsto(\eta,\epsilon\eta)$, and for $\iota=\iota_\epsilon$ one has $J\,U^\epsilon(\mathsf{P})=\mathcal{U}_D(\mathsf{P})\,J$, with $J$ as per Lemma~\ref{lem:local-realization} and $U^\epsilon(\mathsf{P})$ as per Lemma~\ref{lem:parity-induced}. Conversely, for each $\epsilon$ the triple $(D,W_{\mathrm{D}},\iota_\epsilon)$ is a local covariant realization of $(d_{1/2},\mathbb{C}^2)$ whose space of fields is invariant under $\mathcal{U}_D(\mathsf{P})$.
\end{proposition}

\begin{proof}
	By Definition~\ref{def:local-realization}, $F(\mathring{p})=\iota(V)$ is invariant under $D(SU(2))$ and isomorphic to the spin $\frac12$ representation. Lemma~\ref{lem:su2-subspaces} yields $F(\mathring{p})=\mathbb{C}^2\otimes L$ with $\dim L=1$. Lemma~\ref{lem:parity-stable} yields $\gamma^0F(\mathring{p})=F(\mathring{p})$, \textit{i.e.}, $\mathbb{C}^2\otimes\sigma_1L=\mathbb{C}^2\otimes L$, hence $\sigma_1L=L$. The matrix $\sigma_1$ has simple spectrum $\{\pm1\}$, so $L$ is one of the lines $\mathbb{C}(e_1+\epsilon e_2)$. Then $F(\mathring{p})=\{\eta\otimes(e_1+\epsilon e_2)\}=\{(\eta,\epsilon\eta)\}$, and $\gamma^0(\eta,\epsilon\eta)=(\epsilon\eta,\eta)=\epsilon(\eta,\epsilon\eta)$ because $\epsilon^2=1$, which is Equation~\eqref{eq:rest-subspace}. The map $\iota_\epsilon^{-1}\iota:V\to V$ commutes with $\pi$, hence it is a scalar by Lemma 1.2.26.
	
	For the last identity, Lemma~\ref{lem:inversion} and Definition~\ref{def:parity-realization} yield $D(A_{\mathsf{P}p})=D(\theta(A_p))=\gamma^0D(A_p)\gamma^0$, and $\gamma^0\iota_\epsilon=\epsilon\iota_\epsilon$. It descends that
	\begin{equation*}
		(\mathcal{U}_D(\mathsf{P})J\psi)(p)=\gamma^0D(A_{\mathsf{P}p})\iota_\epsilon\psi(\mathsf{P}p)=D(A_p)\gamma^0\iota_\epsilon\psi(\mathsf{P}p)=D(A_p)\iota_\epsilon\,\epsilon\psi(\mathsf{P}p)=(J\,U^\epsilon(\mathsf{P})\psi)(p).
	\end{equation*}
	Conversely, $\iota_\epsilon(u\eta)=(u\eta,\epsilon u\eta)=D(u)\iota_\epsilon\eta$, and $\gamma^0F(\mathring{p})=F(\mathring{p})$ by Equation~\eqref{eq:rest-subspace}. Lemma~\ref{lem:parity-stable} yields the invariance of $\mathcal{F}_D$.
\end{proof}

\noindent Hypotheses used: H4, H5 and H6. The proposition fixes the states at rest up to a sign, and the sign has a meaning that the following remark identifies. It is the eigenvalue of the parity matrix on $F(\mathring{p})$, \textit{i.e.}, the intrinsic parity of Lemma~\ref{lem:parity-induced}.

\begin{remark}[Intrinsic parity and the choice of the sign]\label{rem:intrinsic-parity}
	The sign $\epsilon$ is the intrinsic parity, since $J\,U^\epsilon(\mathsf{P})=\mathcal{U}_D(\mathsf{P})J$ means that space inversion acts on the amplitude $\psi$ as $\epsilon\,\psi(\mathsf{P}p)$. The matrix $\gamma_5:=\operatorname{diag}(\mathbb{1},-\mathbb{1})$ commutes with every $D(g)$ and anticommutes with $\gamma^0$. It maps $\ker(\gamma^0-\epsilon)$ onto $\ker(\gamma^0+\epsilon)$, and it commutes with $D(A_p)$, so that $\gamma_5F^{\epsilon}(p)=F^{-\epsilon}(p)$ for every $p$. Hence $\Psi\mapsto\gamma_5\Psi$ intertwines the representations $\mathcal{U}_D(a,g)$ of the two realizations, and it exchanges the parity matrices $\gamma^0$ and $-\gamma^0$. The two choices differ by this exchange, and the constraint for $\epsilon=-1$ is $(\gamma_\mu p^\mu+m)\Psi=0$, because $\gamma_5\gamma(p)\gamma_5=-\gamma(p)$ by Lemma~\ref{lem:gamma} below. We fix $\epsilon=+1$ as a convention from now on.
\end{remark}

\begin{definition}[Rest projector]\label{def:rest-projector}
	With $\epsilon=+1$ and $\iota\eta:=(\eta,\eta)$, the \emph{rest projector} is the orthogonal projector $\mathcal{Q}(\mathring{p}):=\frac12(\mathbb{1}+\gamma^0)$ of $\mathbb{C}^4$ onto $F(\mathring{p})=\{(\eta,\eta)\mid\eta\in\mathbb{C}^2\}$.
\end{definition}

\noindent The matrix $\gamma^0$ is Hermitian with $(\gamma^0)^2=\mathbb{1}$ by Lemma~\ref{lem:dirac-space}, so that $\mathcal{Q}(\mathring{p})$ is a Hermitian idempotent with range $\ker(\gamma^0-\mathbb{1})=F(\mathring{p})$. Group theory has thus fixed $F(\mathring{p})$ except for the sign $\epsilon$, and no input on the form of the equation has entered.

%--------------------------------------------------------------------
\subsection{Covariant Transport and the Dirac Equation}\label{sec:transport}

The states at rest are fixed, and the states of momentum $p$ are obtained from them by the standard boost, $F(p)=D(A_p)F(\mathring{p})$. Three facts are needed before this transport can be written as an equation. The family $F(p)$ must not depend on the choice of the boost, it must be stable under the whole group and under space inversion, and the projector onto $F(p)$ must be the conjugate of the rest projector. The first two facts are proved in Proposition~\ref{prop:coherence}, and the third is computed in Proposition~\ref{prop:transport}. It follows that the constraint selecting the states of momentum $p$ is the eigenvalue equation of $\gamma_\mu p^\mu$, which is the Dirac equation.

\begin{proposition}[Coherence of the constraint]\label{prop:coherence}
	Let $F(\mathring{p})=\{(\eta,\eta)\mid\eta\in\mathbb{C}^2\}$ and $F(p):=D(A_p)F(\mathring{p})$ for $p\in\mathcal{O}_m^+$.
	\begin{enumerate}
		\item If $A'\in SL(2,\mathbb{C})$ satisfies $\Lambda(A')\mathring{p}=p$, then $D(A')F(\mathring{p})=F(p)$.
		\item $D(g)F(p)=F(\Lambda(g)p)$ for all $g\in SL(2,\mathbb{C})$.
		\item $\gamma^0F(p)=F(\mathsf{P}p)$.
		\item $\dim F(p)=2$.
	\end{enumerate}
\end{proposition}

\begin{proof}
	Item 1. By Proposition~\ref{prop:orbit}, $A'=A_pu$ with $u\in SU(2)$. Since $D(u)(\eta,\eta)=(u\eta,u\eta)$, one has $D(u)F(\mathring{p})=F(\mathring{p})$, and $D(A')F(\mathring{p})=D(A_p)D(u)F(\mathring{p})=F(p)$.
	
	Item 2. Put $q:=\Lambda(g)p$. Then $\Lambda(gA_p)\mathring{p}=q$, and item 1 applied to $A'=gA_p$ yields $D(g)F(p)=D(gA_p)F(\mathring{p})=F(q)$.
	
	Item 3. Lemma~\ref{lem:inversion} and Definition~\ref{def:parity-realization} yield $\gamma^0D(A_p)=D(\theta(A_p))\gamma^0=D(A_{\mathsf{P}p})\gamma^0$. Since $\gamma^0F(\mathring{p})=F(\mathring{p})$, it descends that $\gamma^0F(p)=D(A_{\mathsf{P}p})F(\mathring{p})=F(\mathsf{P}p)$.
	
	Item 4 holds because $D(A_p)$ is invertible.
\end{proof}

\noindent Hypotheses used: H4 and H5. Item 2 says that $p\mapsto F(p)$ is an $SL(2,\mathbb{C})$-invariant subbundle of the trivial bundle $\mathcal{O}_m^+\times\mathbb{C}^4$ as per Remark~\ref{rem:local-subbundle}, and item 1 says that it is determined by its fibre at the rest momentum. The rotation $u$ with $gA_p=A_qu$ is the Wigner rotation $R_W(g,q)$ of Proposition~\ref{prop:induced-action}, so that the dependence on $p$ of the induced representation is absorbed by the stability of $F(\mathring{p})$ under $D(SU(2))$.

The matrices appearing in the transport are computed next. They are the images of the Hermitian matrix model of Definition~\ref{def:hermitian-model} in the Dirac space, and they are not postulated.

\begin{lemma}[Gamma matrices]\label{lem:gamma}
	Let $\gamma^0$ be as per Definition~\ref{def:dirac-space} and $\gamma^k:=\begin{pmatrix}0&-\sigma_k\\\sigma_k&0\end{pmatrix}$ for $k=1,2,3$. For $x\in\mathbb{M}$ set $\gamma(x):=\gamma_\mu x^\mu=\gamma^0x^0-\sum_k\gamma^kx^k$. Then
	\begin{enumerate}
		\item $\gamma(x)=\begin{pmatrix}0&\hat{x}\\\tilde{x}&0\end{pmatrix}$ and $\gamma(x)^2=x^2\,\mathbb{1}$. Hence $\{\gamma^\mu,\gamma^\nu\}=2\eta^{\mu\nu}\mathbb{1}$ with $\eta=\operatorname{diag}(1,-1,-1,-1)$, and $\gamma^{\mu\dagger}=\gamma^0\gamma^\mu\gamma^0$.
		\item $D(g)\,\gamma(x)\,D(g)^{-1}=\gamma(\Lambda(g)x)$ for all $g\in SL(2,\mathbb{C})$.
		\item $\gamma^0\gamma(x)\gamma^0=\gamma(\mathsf{P}x)$.
	\end{enumerate}
\end{lemma}

\begin{proof}
	Item 1. The block form follows from $\hat{x}=x^0\mathbb{1}+\mathbf{x}\cdot\boldsymbol{\sigma}$ and $\tilde{x}=x^0\mathbb{1}-\mathbf{x}\cdot\boldsymbol{\sigma}$. By Lemma~\ref{lem:hermitian-model}, $\gamma(x)^2=\operatorname{diag}(\hat{x}\tilde{x},\tilde{x}\hat{x})=x^2\mathbb{1}$, and polarization yields the Clifford relations. Moreover $\gamma^{0\dagger}=\gamma^0$ and $\gamma^{k\dagger}=\begin{pmatrix}0&\sigma_k\\-\sigma_k&0\end{pmatrix}=\gamma^0\gamma^k\gamma^0$.
	
	Item 2. By Lemma~\ref{lem:sl2c-covering},
	\begin{equation*}
		D(g)\gamma(x)D(g)^{-1}=\begin{pmatrix}g&0\\0&(g^\dagger)^{-1}\end{pmatrix}\begin{pmatrix}0&\hat{x}\\\tilde{x}&0\end{pmatrix}\begin{pmatrix}g^{-1}&0\\0&g^\dagger\end{pmatrix}=\begin{pmatrix}0&g\hat{x}g^\dagger\\(g^\dagger)^{-1}\tilde{x}g^{-1}&0\end{pmatrix}=\gamma(\Lambda(g)x).
	\end{equation*}
	
	Item 3. $\gamma^0\gamma(x)\gamma^0=\begin{pmatrix}0&\tilde{x}\\\hat{x}&0\end{pmatrix}$, while $\widehat{\mathsf{P}x}=\tilde{x}$ and $\widetilde{\mathsf{P}x}=\hat{x}$ by Lemma~\ref{lem:inversion}.
\end{proof}

\noindent Hypotheses used: H6 only, through the form of $D$ on $W_{\mathrm{D}}$. The matrices $\gamma_\mu=\gamma(\mathbf{e}_\mu)$ are the images of the canonical basis of $\mathbb{M}$ under the linear map $x\mapsto\gamma(x)$, and item 2 states that this map is $SL(2,\mathbb{C})$-equivariant. This equivariance is what turns the transport of the rest projector into a statement on $p$.

\begin{proposition}[Transport of the rest projector]\label{prop:transport}
	For every $p\in\mathcal{O}_m^+$,
	\begin{equation}\label{eq:transport}
		D(A_p)\,\gamma^0\,D(A_p)^{-1}=\frac{\gamma_\mu p^\mu}{m}.
	\end{equation}
\end{proposition}

\begin{proof}
	By Lemma~\ref{lem:dirac-space} and Proposition~\ref{prop:orbit},
	\begin{equation*}
		\begin{pmatrix}A_p&0\\0&A_p^{-1}\end{pmatrix}\begin{pmatrix}0&\mathbb{1}\\\mathbb{1}&0\end{pmatrix}\begin{pmatrix}A_p^{-1}&0\\0&A_p\end{pmatrix}=\begin{pmatrix}0&A_p^2\\A_p^{-2}&0\end{pmatrix}=\frac1m\begin{pmatrix}0&\hat{p}\\\tilde{p}&0\end{pmatrix}=\frac{\gamma(p)}{m}.
	\end{equation*}
	The same follows from item 2 of Lemma~\ref{lem:gamma}, since $\gamma^0=\gamma(\mathring{p})/m$ and $\Lambda(A_p)\mathring{p}=p$.
\end{proof}

\noindent Hypotheses used: H6 only. The rest projector is mapped to the four-momentum contracted with the gamma matrices, and the linearity in $p$ comes from $A_p^2=\hat{p}/m$, a property of the standard boost of Proposition~\ref{prop:orbit}. The equation of the section follows.

\begin{theorem}[Dirac equation in momentum space]\label{thm:dirac}
	For $p\in\mathcal{O}_m^+$ let $\mathcal{Q}(p):=D(A_p)\,\mathcal{Q}(\mathring{p})\,D(A_p)^{-1}$. Then
	\begin{equation}\label{eq:covariant-projector}
		\mathcal{Q}(p)=\tfrac12\Big(\mathbb{1}+\frac{\gamma_\mu p^\mu}{m}\Big)
	\end{equation}
	is a projector of rank $2$ with range $F(p)$, and
	\begin{equation}\label{eq:dirac-momentum}
		F(p)=\big\{\Psi\in\mathbb{C}^4\ \big|\ (\gamma_\mu p^\mu-m)\Psi=0\big\}.
	\end{equation}
	Consequently the space of fields of Definition~\ref{def:local-realization} is
	\begin{equation*}
		\mathcal{F}_D=\big\{\Psi:\mathcal{O}_m^+\to\mathbb{C}^4\ \big|\ \Psi\ \text{Borel},\ (\gamma_\mu p^\mu-m)\Psi(p)=0\ \ \text{for all }p\big\},
	\end{equation*}
	and the map $J$ of Lemma~\ref{lem:local-realization}, with $\iota\eta=(\eta,\eta)$, is a bijection from the Borel functions $\mathcal{O}_m^+\to\mathbb{C}^2$ onto $\mathcal{F}_D$ with $J\,U_{m,1/2}(\kappa)=\mathcal{U}_D(\kappa)\,J$ for all $\kappa\in\mathcal{P}$.
\end{theorem}

\begin{proof}
	Equation~\eqref{eq:covariant-projector} is Proposition~\ref{prop:transport} applied to Definition~\ref{def:rest-projector}. The matrix $\mathcal{Q}(p)$ is conjugate to the projector $\mathcal{Q}(\mathring{p})$ of rank $2$, hence it is a projector of rank $2$, and its range is $D(A_p)\operatorname{ran}\mathcal{Q}(\mathring{p})=D(A_p)F(\mathring{p})=F(p)$. A vector $\Psi$ satisfies $\mathcal{Q}(p)\Psi=\Psi$ if and only if $\gamma(p)\Psi=m\Psi$, which yields Equation~\eqref{eq:dirac-momentum}. The statement on $J$ is Lemma~\ref{lem:local-realization}.
\end{proof}

\noindent Hypotheses used: H4, H5 and H6, with the convention $\epsilon=+1$. The theorem identifies the space of states at fixed momentum with the positive eigenspace of $\gamma_\mu p^\mu/m$, and it does so for the whole group at once, since $J$ intertwines $U_{m,1/2}$ with $\mathcal{U}_D$. The scalar product that makes $J$ unitary is the object of the following section.

\begin{remark}[What was an output]\label{rem:dirac-outputs}
	The number of components is the dimension of $W_{\mathrm{D}}$, which is fixed by Proposition~\ref{prop:minimal-extension}. The gamma matrices are the block matrices of Lemma~\ref{lem:gamma}, built from $\hat{x}$ and $\tilde{x}$. The first-order form of the equation expresses the fact that $D(A_p)\gamma^0D(A_p)^{-1}$ is linear in $p$, and the mass enters as the eigenvalue of this matrix because $\gamma^0$ has eigenvalue $1$ on $F(\mathring{p})$. The only choice left is the sign $\epsilon$ of Remark~\ref{rem:intrinsic-parity}.
\end{remark}

The equation is now written in position space, using that the plane wave $e^{-ip\cdot x/\hbar}$ satisfies $i\hbar\,\partial_\mu e^{-ip\cdot x/\hbar}=p_\mu e^{-ip\cdot x/\hbar}$. Multiplication by $\gamma_\mu p^\mu$ becomes a first-order differential operator.

\begin{corollary}[Dirac equation in position space]\label{cor:dirac-position}
	Let $\Psi\in\mathcal{F}_D$ be continuous with compact support, and set $\Psi(x):=\int e^{-ip\cdot x/\hbar}\Psi(p)\,d\mu(p)$ for $x\in\mathbb{M}$. Then $\Psi(x)$ is smooth and
	\begin{equation}\label{eq:dirac-position}
		\big(i\hbar\,\gamma^\mu\partial_\mu-m\big)\Psi(x)=0,
	\end{equation}
	\textit{i.e.}, $i\hbar\,\partial_t\Psi=\big(\boldsymbol{\alpha}\cdot(-i\hbar\nabla)+m\beta\big)\Psi$ with $\beta:=\gamma^0$ and $\alpha^k:=\gamma^0\gamma^k=\operatorname{diag}(\sigma_k,-\sigma_k)$.
\end{corollary}

\begin{proof}
	The support of $\Psi$ is compact in $\mathcal{O}_m^+\cong\mathbb{R}^3$, so differentiation under the integral is allowed by dominated convergence. It yields $i\hbar\gamma^\mu\partial_\mu\Psi(x)=\int\gamma_\mu p^\mu e^{-ip\cdot x/\hbar}\Psi(p)\,d\mu(p)$, and Equation~\eqref{eq:dirac-momentum} turns the integrand into $m\,e^{-ip\cdot x/\hbar}\Psi(p)$. Multiplying by $\gamma^0$ and using $(\gamma^0)^2=\mathbb{1}$ yields the Hamiltonian form with $p^0\to i\hbar\partial_t$ and $\mathbf{p}\to-i\hbar\nabla$. The block form of $\alpha^k$ follows from the definition of $\gamma^k$.
\end{proof}

\noindent Hypotheses used: H4, H5 and H6, with $\epsilon=+1$. The Hamiltonian form closes the loop with Theorem~\ref{thm:relativistic-hamiltonian}. The matrix $h_{\mathrm{D}}(\mathbf{p}):=\boldsymbol{\alpha}\cdot\mathbf{p}+m\beta$ is the Dirac Hamiltonian of a plane wave, and the constraint selects its positive spectral subspace.

\begin{proposition}[The constraint is the positive spectral subspace]\label{prop:dirac-hamiltonian}
	The matrix $h_{\mathrm{D}}(\mathbf{p})$ is Hermitian with $h_{\mathrm{D}}(\mathbf{p})^2=\omega(\mathbf{p})^2\,\mathbb{1}$, and $F(p)=\ker\big(h_{\mathrm{D}}(\mathbf{p})-\omega(\mathbf{p})\big)$ for $p=(\omega(\mathbf{p}),\mathbf{p})\in\mathcal{O}_m^+$. In particular the energy $p^0=\omega(\mathbf{p})$ of Theorem~\ref{thm:relativistic-hamiltonian} is the eigenvalue of $h_{\mathrm{D}}(\mathbf{p})$ on $F(p)$, and $h_{\mathrm{D}}(\mathbf{p})$ has eigenvalues $\pm\omega(\mathbf{p})$, each with multiplicity $2$.
\end{proposition}

\begin{proof}
	The matrices $\alpha^k=\operatorname{diag}(\sigma_k,-\sigma_k)$ and $\beta$ are Hermitian. Block multiplication yields $\{\alpha^j,\alpha^k\}=2\delta_{jk}\mathbb{1}$, $\{\alpha^k,\beta\}=0$ and $\beta^2=\mathbb{1}$, so that $h_{\mathrm{D}}^2=(|\mathbf{p}|^2+m^2)\mathbb{1}$. Since $\operatorname{tr}\alpha^k=\operatorname{tr}\beta=0$, the eigenvalues $\pm\omega$ have equal multiplicity. Multiplying $\gamma(p)\Psi=m\Psi$ by $\gamma^0$ yields the equivalent equation $p^0\Psi=h_{\mathrm{D}}(\mathbf{p})\Psi$, because $\gamma^0\gamma^kp^k=\boldsymbol{\alpha}\cdot\mathbf{p}$. Since $p^0=\omega(\mathbf{p})$ on $\mathcal{O}_m^+$, Equation~\eqref{eq:dirac-momentum} reads $F(p)=\ker(h_{\mathrm{D}}-\omega)$.
\end{proof}

\noindent Hypotheses used: H6, with H3 entering only through $p^0>0$ on $\mathcal{O}_m^+$ as per Proposition~\ref{prop:mass-shell}. The negative eigenspace $\ker(h_{\mathrm{D}}+\omega)$ is not part of $F(p)$, and it would correspond to the shell $\mathcal{O}_m^-$ excluded by positivity of the energy.

\begin{remark}\label{rem:Q-not-hermitian}
	The projector $\mathcal{Q}(p)$ is not Hermitian for $\mathbf{p}\neq\mathbf{0}$. Indeed $\gamma(p)^\dagger=\begin{pmatrix}0&\tilde{p}\\\hat{p}&0\end{pmatrix}$ differs from $\gamma(p)$ as soon as $\hat{p}\neq\tilde{p}$, \textit{i.e.}, as soon as $\mathbf{p}\neq\mathbf{0}$. This reflects the fact that $D(A_p)$ is Hermitian and positive definite but not unitary, so that conjugation by it turns the orthogonal projector $\mathcal{Q}(\mathring{p})$ into an oblique one. The subspace $F(p)$ is therefore not the Euclidean orthogonal complement of $\ker(\gamma(p)+m)$, and the Euclidean scalar product of $\mathbb{C}^4$ is not invariant, in agreement with Proposition~\ref{prop:no-local-unitary}. The invariant positive form on $F(p)$ is computed in Proposition~\ref{prop:norms} and used in Theorem~\ref{thm:covariant-norm}.
\end{remark}

\section{Return to the State Space}
\label{sec:probability}

The states of an elementary relativistic system are the points of the projective space $\mathbb{P}\mathcal{H}$ with the Kähler structure of Chapter 3, and the Poincar\'e group acts on them by Kähler isomorphisms. Sections 6.4 to 6.6 replaced $\mathbb{P}\mathcal{H}$ by increasingly explicit spaces of rays. The last of them consists of the rays of the solutions of the Dirac equation, on which the symmetry group acts through the fixed matrices $D(g)$. The question of this section is whether this space of rays is a model of the same state space, \textit{i.e.}, a Kähler manifold isomorphic to $\mathbb{P}\mathcal{H}$, on which both the Poincar\'e action and the time evolution take the form they have on $\mathbb{P}\mathcal{H}$. The data are the unitary operator $J$ of Lemma~\ref{lem:local-realization}, the representation $\mathcal{U}_D$ and the fibres $F(p)$ of Theorem~\ref{thm:dirac}. These fibres carry two quadratic forms, and the section compares them. It identifies the one reproducing the scalar product of $\mathcal{H}_{m,1/2}$ and shows that $J$ induces a Kähler isomorphism, so that the Dirac rays realize $(\mathbb{P}\mathcal{H},\Phi)$ in the sense of Definition~\ref{def:realization}. It then proves that the Hamiltonian flow of the mean energy on the Dirac rays is the Dirac evolution and records the conserved current. This is the last step of the chain. It supplies the premises of the proof of part (B) of Theorem~\ref{thm:main}, given at the end of the section. To avoid a clash with the parity matrix of $\mathbf{H5}$, spacetime is written $\mathbb{R}^{1,3}$ here, and $p^0$ denotes the energy on $\mathcal{O}_m^+$, while $\omega$ is the symplectic form of Chapter 3. Unless stated otherwise, $j=\tfrac12$ and $\epsilon=+1$.

\begin{proposition}[Explicit solutions]\label{prop:explicit-solutions}
	Let $(D,\mathbb{C}^4,\iota)$ be the Dirac realization of Theorem~\ref{thm:dirac}. For $p\in\mathcal{O}_m^+$,
	\begin{equation*}
		F(p)=\big\{\Psi=(A_p\eta,\ A_p^{-1}\eta)^T\ \big|\ \eta\in\mathbb{C}^2\big\},
	\end{equation*}
	and for $\psi\in\mathcal{H}_{m,1/2}$ one has $(J\psi)(p)=(A_p\psi(p),A_p^{-1}\psi(p))^T$, \textit{i.e.}, $\eta=\psi(p)$.
\end{proposition}

\begin{proof}
	By Lemma~\ref{lem:dirac-space}, $D(A_p)=\operatorname{diag}(A_p,A_p^{-1})$, and $F(\mathring{p})=\{(\eta,\eta)\}$ by Definition~\ref{def:rest-projector}. It descends that $F(p)=D(A_p)F(\mathring{p})=\{(A_p\eta,A_p^{-1}\eta)\}$, and the formula for $J$ is Lemma~\ref{lem:local-realization} with $\iota\eta=(\eta,\eta)$. These vectors solve $(\gamma_\mu p^\mu-m)\Psi=0$, since by Lemma~\ref{lem:gamma} and Equation~\eqref{eq:standard-boost}
	\begin{equation*}
		\gamma(p)\Psi=(\hat{p}A_p^{-1}\eta,\ \tilde{p}A_p\eta)=(mA_p\eta,\ mA_p^{-1}\eta)=m\Psi.
	\end{equation*}
\end{proof}

The Euclidean form $\Psi^\dagger\Phi$ is the one that $D(g)$ fails to preserve, as Proposition~\ref{prop:no-unitary} excludes unitarity. The Dirac form $\bar{\Psi}\Phi:=\Psi^\dagger\gamma^0\Phi$, with Dirac adjoint $\bar{\Psi}:=\Psi^\dagger\gamma^0$, is invariant by item 2 of Lemma~\ref{lem:dirac-space},
\begin{equation}\label{eq:dirac-form-invariance}
	\overline{D(g)\Psi}\,D(g)\Phi=\Psi^\dagger D(g)^\dagger\gamma^0D(g)\Phi=\bar{\Psi}\Phi,\qquad g\in SL(2,\mathbb{C}).
\end{equation}
The next proposition compares the two forms on $F(p)$ and shows that the invariant one is positive definite there.

\begin{proposition}[Euclidean and covariant norm]\label{prop:norms}
	Let $\Psi=(A_p\eta,A_p^{-1}\eta)$ and $\Phi=(A_p\eta',A_p^{-1}\eta')$ be in $F(p)$.
	\begin{enumerate}
		\item $\bar{\Psi}\Phi=2(\eta,\eta')$. In particular $\bar{\Psi}\Psi=2|\eta|^2$, and the Dirac form is positive definite on $F(p)$.
		\item For every $q\in\mathbb{R}^{1,3}$, $\bar{\Psi}\gamma(q)\Psi=\frac{2\,p\cdot q}{m}|\eta|^2$, \textit{i.e.}, $\bar{\Psi}\gamma^\mu\Psi=\frac{2p^\mu}{m}|\eta|^2$. In particular
		\begin{equation}\label{eq:euclidean-norm}
			\Psi^\dagger\Psi=\frac{2p^0}{m}|\eta|^2=\frac{p^0}{m}\,\bar{\Psi}\Psi.
		\end{equation}
	\end{enumerate}
\end{proposition}

\begin{proof}
	Item 1. Since $\gamma^0\Phi=(A_p^{-1}\eta',A_p\eta')$ and $A_p$ is Hermitian, $\bar{\Psi}\Phi=\eta^\dagger A_pA_p^{-1}\eta'+\eta^\dagger A_p^{-1}A_p\eta'=2(\eta,\eta')$.
	
	Item 2. By Lemma~\ref{lem:gamma}, $\gamma^0\gamma(q)=\operatorname{diag}(\tilde{q},\hat{q})$, so that $\bar{\Psi}\gamma(q)\Psi=\eta^\dagger\big(A_p\tilde{q}A_p+A_p^{-1}\hat{q}A_p^{-1}\big)\eta$. Equation~\eqref{eq:tilde-covariance} with $g=A_p^{-1}$ and Equation~(lorentz-action) with the same $g$ yield $A_p\tilde{q}A_p=\widetilde{q'}$ and $A_p^{-1}\hat{q}A_p^{-1}=\widehat{q'}$, where $q':=\Lambda(A_p)^{-1}q$. The bracket equals $\widetilde{q'}+\widehat{q'}=2q'^0\,\mathbb{1}$, and $q'^0=q'\cdot e_0=q\cdot\Lambda(A_p)e_0=q\cdot p/m$ by invariance of the form. Taking $q=e_0$, for which $\gamma(e_0)=\gamma^0$, yields Equation~\eqref{eq:euclidean-norm}.
\end{proof}

\noindent The factor $p^0/m$ is the time dilation factor of a particle of momentum $p$, because $\bar{\Psi}\Psi$ is a scalar while $\Psi^\dagger\Psi=\bar{\Psi}\gamma^0\Psi$ is the time component of the four vector $\bar{\Psi}\gamma^\mu\Psi$.

Proposition~\ref{prop:norms} identifies the form that must replace the Euclidean one. It remains to show that it makes the whole construction unitary and not only the fibres positive.

\begin{theorem}[Covariant norm and equivalence with the induced representation]\label{thm:covariant-norm}
	Let $\mathcal{F}_D$ and $\mathcal{U}_D$ be as per Theorem~\ref{thm:dirac}, which uses $\mathbf{H4}$ and $\mathbf{H5}$ through the Dirac space. On $\mathcal{F}_D$ set
	\begin{equation*}
		\langle\Psi,\Phi\rangle_D:=\tfrac12\int\bar{\Psi}(p)\Phi(p)\,d\mu(p),\qquad \mathcal{F}_D^2:=\Big\{\Psi\in\mathcal{F}_D\ \Big|\ \int\bar{\Psi}\Psi\,d\mu<\infty\Big\},
	\end{equation*}
	with fields identified when they agree $\mu$-almost everywhere. Then $\langle\cdot,\cdot\rangle_D$ is a positive definite scalar product, $\mathcal{F}_D^2$ is a separable Hilbert space, $J:\mathcal{H}_{m,1/2}\to\mathcal{F}_D^2$ is unitary, and $\mathcal{U}_D$ is a representation of $\mathcal{P}$ by unitary operators on $\mathcal{F}_D^2$ with $J\,U_{m,1/2}(a,g)=\mathcal{U}_D(a,g)\,J$. Hence $\mathcal{U}_D$ is strongly continuous, irreducible and unitarily equivalent to $U_{m,1/2}$.
\end{theorem}

\begin{proof}
	By Proposition~\ref{prop:norms}, $\tfrac12\overline{J\psi}(p)(J\phi)(p)=(\psi(p),\phi(p))_{\mathbb{C}^2}$ pointwise, so that $\langle J\psi,J\phi\rangle_D=\langle\psi,\phi\rangle_{\mathcal{H}_{m,1/2}}$ whenever one side is finite. The map $J$ is a bijection from the Borel functions onto $\mathcal{F}_D$ by Lemma~\ref{lem:local-realization}. Its inverse $\psi(p)=\iota^{-1}D(A_p)^{-1}\Psi(p)$ is Borel because $p\mapsto A_p$ is smooth, and $\|\psi(p)\|^2=\tfrac12\bar{\Psi}\Psi(p)$. Hence $J$ restricts to a bijective isometry of $\mathcal{H}_{m,1/2}$ onto $\mathcal{F}_D^2$, which is therefore a separable Hilbert space and the form is a positive definite scalar product. Equation~\eqref{eq:J-intertwines} yields $\mathcal{U}_D=JU_{m,1/2}J^{-1}$, so that $\mathcal{U}_D$ is unitary, strongly continuous and irreducible by Theorem~\ref{thm:wigner-mackey}. Directly, unitarity also follows from Equation~\eqref{eq:dirac-form-invariance}, $|e^{ia\cdot p/\hbar}|=1$ and the invariance of $\mu$,
	\begin{equation*}
		\overline{(\mathcal{U}_D(a,g)\Psi)(p)}\,(\mathcal{U}_D(a,g)\Phi)(p)=\bar{\Psi}(\Lambda(g)^{-1}p)\,\Phi(\Lambda(g)^{-1}p).
	\end{equation*}
\end{proof}

\noindent The unitary structure is not lost in passing to the non-unitary $D$. It is carried by the form $\tfrac12\bar{\Psi}\Phi$, which is indefinite on $\mathbb{C}^4$ and positive definite only on the fibres $F(p)$.

The scalar product $\langle\cdot,\cdot\rangle_D$ determines a projective space and a Kähler structure through Chapter 3. Since $J$ preserves the scalar product, the identification of this structure with the one of $\mathbb{P}\mathcal{H}$ requires only the expression of $g$ and $\omega$ through the inner product. The following proposition records the identification and its consequence for $\Phi$.

\begin{proposition}[Kähler structure on the Dirac rays]\label{prop:kahler-dirac}
	Let $\mathbb{P}\mathcal{F}_D^2$ carry the Kähler structure of Chapter 3 built from $\langle\cdot,\cdot\rangle_D$ as per Equation (3.24). Then $\mathbb{P}J:\mathbb{P}\mathcal{H}_{m,1/2}\to\mathbb{P}\mathcal{F}_D^2$, $[\psi]\mapsto[J\psi]$, is an isomorphism of Kähler manifolds, and
	\begin{equation*}
		\mathbb{P}J\circ[U_{m,1/2}(\kappa)\,\cdot\,]=[\mathcal{U}_D(\kappa)\,\cdot\,]\circ\mathbb{P}J,\qquad \kappa\in\mathcal{P}.
	\end{equation*}
	Let $\Phi$ be an elementary relativistic system of mass $m>0$ and spin $\tfrac12$ ($\mathbf{H1}$ to $\mathbf{H3}$), and let $\mathcal{W}$ be the unitary operator of Theorem~\ref{thm:wigner-mackey}. Then $F:=\mathbb{P}(J\mathcal{W}):\mathbb{P}\mathcal{H}\to\mathbb{P}\mathcal{F}_D^2$ is a realization of $\Phi$ as per Definition~\ref{def:realization}, with $F\circ\Phi_\kappa=[\mathcal{U}_D(\kappa)\,\cdot\,]\circ F$.
\end{proposition}

\begin{proof}
	By Definition 3.2.5 and Proposition 3.2.6, the chart at $[\psi]$ with unit representative $\psi$ is $\phi\mapsto\phi/(\psi,\phi)-\psi$. Since $J$ is a linear isometry, $\mathcal{F}_{[J\psi]}\circ\mathbb{P}J=J\circ\mathcal{F}_{[\psi]}$, with $J$ mapping $\psi^\perp$ onto $(J\psi)^\perp$. Hence $\mathbb{P}J$ is a biholomorphic diffeomorphism for the atlas of Proposition 3.2.7, and its differential acts as $J$ on $\psi^\perp$. By Equation (3.26) the Fubini--Study metric and the symplectic form at $[\psi]$ are expressed through the scalar product on $\psi^\perp$, which $J$ preserves. It descends that $\mathbb{P}J$ preserves the three structures. The intertwining relation is Theorem~\ref{thm:covariant-norm}. For the second statement, $\mathbb{P}\mathcal{W}$ is a Kähler isomorphism for the same reason, and $\mathcal{W}\Pi(\kappa)\mathcal{W}^{-1}=U_{m,1/2}(\kappa)$ by Theorem~\ref{thm:wigner-mackey}. Hence $F\circ\Phi_\kappa=[J\mathcal{W}\Pi(\kappa)\,\cdot\,]=[\mathcal{U}_D(\kappa)J\mathcal{W}\,\cdot\,]$.
\end{proof}

\noindent By Proposition 3.3.6 each $[\mathcal{U}_D(\kappa)\,\cdot\,]$ is a Kähler isomorphism of $\mathbb{P}\mathcal{F}_D^2$, and by Proposition~\ref{prop:realization-equivalence} the realization $F$ corresponds to the exact unitary equivalence $J\mathcal{W}\,\Pi=\mathcal{U}_D\,J\mathcal{W}$.

The energy of the system is an unbounded operator, whereas Proposition 3.3.2 treats a bounded Hamiltonian $H$ and defines the vector field through the musical isomorphism acting on $d\langle H\rangle$. For unbounded $A$ the function $\langle A\rangle$ of Definition~\ref{def:mean-value-domain} is defined only on $\mathbb{P}\mathcal{D}$, which is dense but not open, so that $d\langle A\rangle$ exists only along directions represented in $\mathcal{D}$. The Hamiltonian vector is therefore characterized by its pairing with these directions, in analogy with Lemma~\ref{lem:pointwise-formulas}.

\begin{definition}[Hamiltonian vector of a mean value function]\label{def:hamiltonian-vector}
	Let $K$ be a Hilbert space with the Kähler structure of $\mathbb{P}K$ as per Chapter 3, let $A$ be a selfadjoint operator on $K$ and let $\mathcal{D}\subseteq D(A)$ be a dense subspace with $A\mathcal{D}\subseteq\mathcal{D}$. For $[x]\in\mathbb{P}\mathcal{D}$ with $\|x\|=1$, a vector $\xi\in x^\perp\cong T_{[x]}\mathbb{P}K$ is a \emph{Hamiltonian vector} of $\langle A\rangle$ at $[x]$ if
	\begin{equation*}
		\omega_{[x]}(\xi,\zeta)=d\langle A\rangle_{[x]}(\zeta):=\frac{d}{ds}\langle A\rangle([x+s\zeta])\Big|_{s=0},\qquad \zeta\in x^\perp\cap\mathcal{D}.
	\end{equation*}
\end{definition}

\begin{proposition}[The Hamiltonian flow is the Dirac evolution]\label{prop:dirac-flow}
	Let $\Phi$, $\mathcal{W}$ and $F$ be as in Proposition~\ref{prop:kahler-dirac}, with hypotheses $\mathbf{H1}$ to $\mathbf{H3}$. Let $H^\infty_\Phi$ be the space of smooth vectors of the lift $\Pi$ of Proposition~\ref{prop:poincare-lift}, and set $H_D:=J\mathcal{W}H\mathcal{W}^{-1}J^{-1}$ and $\mathcal{F}_D^\infty:=J\mathcal{W}H^\infty_\Phi$, with $K=\mathcal{F}_D^2$ and $\mathcal{D}=\mathcal{F}_D^\infty$ in Definition~\ref{def:hamiltonian-vector}.
	\begin{enumerate}
		\item $H_D$ is the selfadjoint operator on $\mathcal{F}_D^2$ of multiplication by $p^0$. The space $\mathcal{F}_D^\infty$ is dense, invariant under every $\mathcal{U}_D(\kappa)$, contained in $D(H_D)$, and $H_D\mathcal{F}_D^\infty\subseteq\mathcal{F}_D^\infty$.
		\item Let $\Psi\in\mathcal{F}_D^\infty$ with $\|\Psi\|_D=1$. Then $\Psi_t:=e^{-itH_D/\hbar}\Psi=\mathcal{U}_D((-t,\mathbf{0}),\mathbb{1})\Psi$ belongs to $\mathcal{F}_D^\infty$, $\Psi_t(p)=e^{-itp^0/\hbar}\Psi(p)$, and $h:=\langle H_D\rangle([\Psi_t])$ does not depend on $t$.
		\item The curve $\gamma(t):=[\Psi_t]$ is differentiable in $\mathbb{P}\mathcal{F}_D^2$, with velocity represented by $\eta_t:=-\tfrac{i}{\hbar}(H_D-h)\Psi_t\in\Psi_t^\perp\cap\mathcal{F}_D^\infty$. The vector $\eta_t$ is the unique Hamiltonian vector of $\langle H_D\rangle$ at $[\Psi_t]$, so that $\gamma$ is an integral curve of the Hamiltonian vector field of $\langle H_D\rangle$.
		\item Let $y:I\to\mathcal{F}_D^\infty$ be strongly differentiable with $\|y(t)\|_D=1$, and suppose that the velocity of $t\mapsto[y(t)]$ is the Hamiltonian vector of $\langle H_D\rangle$ at $[y(t)]$ for every $t$. Then $[y(t)]=[e^{-itH_D/\hbar}y(0)]$ for $0\in I$.
		\item Let $\Psi\in\mathcal{F}_D^2$ be continuous with compact support. The field $\Psi(x)$ of Corollary~\ref{cor:dirac-position} satisfies $\Psi(t,\mathbf{x})=\int e^{i\mathbf{p}\cdot\mathbf{x}/\hbar}\Psi_t(p)\,d\mu(p)$, \textit{i.e.}, the Cauchy datum at time $t$ is the spatial Fourier transform of $\Psi_t$.
	\end{enumerate}
\end{proposition}

\begin{proof}
	Item 1. By Proposition~\ref{prop:induced-unitary} and Theorem~\ref{thm:wigner-mackey}, $\mathcal{W}H\mathcal{W}^{-1}$ is the multiplication by $p^0$ on $\mathcal{H}_{m,1/2}$, and $J$ intertwines it with $H_D$. Theorem~\ref{thm:smooth-vectors}, which holds for any strongly continuous unitary representation of a Lie group and hence for $\Pi$ by Proposition~\ref{prop:poincare-lift}, states that $H^\infty_\Phi$ is dense, invariant under $\Pi(\kappa)$ and under the operators $d\Pi(\xi)$, whose closures are the generators. Since $S(s)=\Pi(\tau_s)=e^{isH/\hbar}$ as per Definition~\ref{def:poincare-generators}, $H=-i\hbar\,d\Pi(\xi_0)$ on $H^\infty_\Phi$, which yields the remaining statements after transport through the unitary operator $J\mathcal{W}$. Since the lift is unique by Proposition~\ref{prop:poincare-lift}, no analogue of item 1 of Proposition~\ref{prop:unique-smooth} is needed, and item 2 of that proposition holds for the same reason as there.
	
	Item 2. By item 1, $\mathcal{F}_D^\infty$ is invariant under $\mathcal{U}_D((-t,\mathbf{0}),\mathbb{1})$, which acts by multiplication by $e^{-itp^0/\hbar}$ and equals $e^{-itH_D/\hbar}$ by the active convention of Remark~\ref{rem:active-convention}. The function $\langle H_D\rangle$ is constant because $e^{itH_D/\hbar}$ commutes with $H_D$ on $D(H_D)$ and is unitary.
	
	Item 3. Apply item 1 of Lemma~\ref{lem:pointwise-formulas} to $A=B=H_D$ and $\mathcal{D}=\mathcal{F}_D^\infty$ at the base point $\Psi_t$, noting that $s\mapsto[e^{-isH_D/\hbar}\Psi_t]=\gamma(t+s)$. It yields the velocity $\eta_t$. By item 2 of the same lemma, $d\langle H_D\rangle_{[\Psi_t]}(\zeta)=2\operatorname{Re}(\zeta,H_D\Psi_t)_D$ for $\zeta\in\Psi_t^\perp\cap\mathcal{F}_D^\infty$. By Equation (3.26), for $\xi,\zeta\in\Psi_t^\perp$ one has $\omega(\xi,\zeta)=2\hbar\operatorname{Im}(\xi,\zeta)_D$. Since $(\Psi_t,\zeta)_D=0$,
	\begin{equation*}
		\omega(\eta_t,\zeta)=2\hbar\operatorname{Im}\Big(\tfrac{i}{\hbar}(H_D\Psi_t,\zeta)_D\Big)=2\operatorname{Re}(H_D\Psi_t,\zeta)_D=d\langle H_D\rangle_{[\Psi_t]}(\zeta),
	\end{equation*}
	so that $\eta_t$ is a Hamiltonian vector. For uniqueness, the space $\Psi_t^\perp\cap\mathcal{F}_D^\infty$ is dense in $\Psi_t^\perp$, since $\zeta\in\Psi_t^\perp$ is the limit of $\zeta_n-(\Psi_t,\zeta_n)_D\Psi_t$ with $\zeta_n\in\mathcal{F}_D^\infty$, $\zeta_n\to\zeta$. If $\xi\in\Psi_t^\perp$ satisfies $\omega(\xi,\zeta)=0$ on this dense subspace, continuity gives $\operatorname{Im}(\xi,\zeta)_D=0$ on $\Psi_t^\perp$, and $\zeta=i\xi$ yields $\xi=0$.
	
	Item 4. The velocity of $[y(t)]$ is represented by the component of $y'(t)$ orthogonal to $y(t)$. Since $\operatorname{Re}(y,y')_D=0$, one has $y'=\eta+i\beta y$ with $\beta(t)\in\mathbb{R}$ and $\eta=-\tfrac{i}{\hbar}(H_D-h(t))y$ by item 3. Set $w(t):=e^{itH_D/\hbar}y(t)$. Strong continuity of the group and uniform boundedness yield $w'=e^{itH_D/\hbar}\big(y'+\tfrac{i}{\hbar}H_Dy\big)=i(\beta+h/\hbar)w$. Then $q(t):=(w(0),w(t))_D$ satisfies $q'=i(\beta+h/\hbar)q$, hence $|q|$ is constant and equal to $1$. The equality case of the Cauchy--Schwarz inequality yields $w(t)=q(t)w(0)$, and the claim follows.
	
	Item 5. By Corollary~\ref{cor:dirac-position} and the definition of $\Psi(x)$, $e^{-ip\cdot x/\hbar}=e^{i(\mathbf{p}\cdot\mathbf{x}-p^0t)/\hbar}$ at $x=(t,\mathbf{x})$, and the factor $e^{-ip^0t/\hbar}$ is absorbed in $\Psi_t$.
\end{proof}

\noindent Item 4 shows that the flow is determined by the vector field, so that the Dirac evolution is the Hamiltonian flow of the mean energy and not merely a curve solving the same equation.

By Proposition~\ref{prop:kahler-dirac} and Proposition~\ref{prop:dirac-flow}, $F$ carries the flow $[e^{-itH/\hbar}\,\cdot\,]$ of Chapter 5 on $\mathbb{P}H^\infty_\Phi$ onto the Dirac evolution on $\mathbb{P}\mathcal{F}_D^\infty$, because $\langle H_D\rangle\circ F=\langle H\rangle$ and $F$ preserves $\omega$. The role of the scalar product in this statement is the content of the next remark.

\begin{remark}[The Euclidean form cannot be used]\label{rem:euclidean-norm}
	Let $\|\Psi\|_E^2:=\int\Psi^\dagger\Psi\,d\mu$ be the norm defined by the Euclidean form of $\mathbb{C}^4$ on $\mathcal{F}_D$. By Equation~\eqref{eq:euclidean-norm}, Equation~\eqref{eq:dirac-form-invariance} and the invariance of $\mu$,
	\begin{equation*}
		\|\mathcal{U}_D(0,g)\Psi\|_E^2=\int\frac{(\Lambda(g)q)^0}{m}\,\bar{\Psi}(q)\Psi(q)\,d\mu(q),\qquad \|\Psi\|_E^2=\int\frac{q^0}{m}\,\bar{\Psi}(q)\Psi(q)\,d\mu(q).
	\end{equation*}
	Take $g=A_k$ with $\mathbf{k}\neq\mathbf{0}$, so that $(\Lambda(A_k)q)^0=(q^0k^0+\mathbf{q}\cdot\mathbf{k})/m$. For fields with $\bar{\Psi}\Psi=2|v|^2\chi_B$ and $B$ shrinking to a point $q_0$, the ratio $\|\mathcal{U}_D(0,g)\Psi\|_E^2/\|\Psi\|_E^2$ tends to $(\Lambda(g)q_0)^0/q_0^0$. At $q_0=\mathring{p}$ this is $k^0/m>1$. Along $\mathbf{q}=-s\mathbf{k}/|\mathbf{k}|$ with $s\to\infty$ it tends to $(k^0-|\mathbf{k}|)/m<1$, because $k^0-|\mathbf{k}|=m^2/(k^0+|\mathbf{k}|)<m$. If $[\mathcal{U}_D(0,g)\,\cdot\,]$ were a Kähler isomorphism of $\mathbb{P}\mathcal{F}_D$ with the Euclidean structure, Proposition 3.3.6 would yield a unitary $V$ with $[\mathcal{U}_D(0,g)\Psi]=[V\Psi]$ for all $\Psi$. Then $V^{-1}\mathcal{U}_D(0,g)$ maps every vector to a multiple of itself, hence is a scalar $c$, and the ratio would equal $|c|^2$ for every $\Psi$. The two limits exclude this. The same weight enters the mean energy. With $\nu:=\bar{\Psi}\Psi\,d\mu/\int\bar{\Psi}\Psi\,d\mu$, the Euclidean form yields $\int(p^0)^2d\nu/\int p^0d\nu$ and the covariant one yields $\int p^0d\nu$, which is $\langle H_D\rangle$. The Cauchy--Schwarz inequality shows that the first exceeds the second unless $p^0$ is $\nu$-almost everywhere constant.
\end{remark}

The covariant form also fixes the conserved density. The time component of the current $\bar{\Psi}\gamma^\mu\Psi$ is the Euclidean density of the position-space field, and its spatial integral is determined by the Wigner norm.

\begin{proposition}[Conservation of probability]\label{prop:continuity}
	Let $\Psi\in\mathcal{F}_D$ be smooth with compact support in $\mathbf{p}$ and let $\Psi(x)=\int e^{-ip\cdot x/\hbar}\Psi(p)\,d\mu(p)$ be the position-space field, the argument $x\in\mathbb{R}^{1,3}$ distinguishing it from the amplitude. Then $j^\mu:=\bar{\Psi}\gamma^\mu\Psi$ satisfies
	\begin{equation}\label{eq:continuity}
		\partial_\mu j^\mu=0,\qquad j^0=\Psi^\dagger\Psi\geq0,
	\end{equation}
	and for every $t\in\mathbb{R}$, with $\psi=J^{-1}\Psi$,
	\begin{equation}\label{eq:charge}
		\int_{\mathbb{R}^3}\Psi^\dagger(t,\mathbf{x})\Psi(t,\mathbf{x})\,d^3\mathbf{x}=\frac{(2\pi\hbar)^3}{m}\,\|\psi\|^2_{\mathcal{H}_{m,1/2}}.
	\end{equation}
\end{proposition}

\begin{proof}
	By Corollary~\ref{cor:dirac-position}, $i\hbar\gamma^\mu\partial_\mu\Psi=m\Psi$. Taking the adjoint, using $\gamma^{\mu\dagger}=\gamma^0\gamma^\mu\gamma^0$ and multiplying on the right by $\gamma^0$ yields $-i\hbar\,\partial_\mu\bar{\Psi}\gamma^\mu=m\bar{\Psi}$. It descends that
	\begin{equation*}
		\partial_\mu(\bar{\Psi}\gamma^\mu\Psi)=(\partial_\mu\bar{\Psi}\gamma^\mu)\Psi+\bar{\Psi}(\gamma^\mu\partial_\mu\Psi)=\frac{im}{\hbar}\bar{\Psi}\Psi-\frac{im}{\hbar}\bar{\Psi}\Psi=0,
	\end{equation*}
	and $j^0=\bar{\Psi}\gamma^0\Psi=\Psi^\dagger\Psi$. For Equation~\eqref{eq:charge}, at time $t$ one has $\Psi(t,\mathbf{x})=\int d^3\mathbf{p}\,e^{i(\mathbf{p}\cdot\mathbf{x}-p^0t)/\hbar}\Psi(p)/2p^0$, a Fourier transform multiplied by a phase of modulus one. The Plancherel theorem \cite[TODO-location]{TODO-RS03} and Equation~\eqref{eq:euclidean-norm} yield
	\begin{equation*}
		\int\Psi^\dagger\Psi\,d^3\mathbf{x}=(2\pi\hbar)^3\int\frac{\Psi^\dagger\Psi}{4(p^0)^2}\,d^3\mathbf{p}=\frac{(2\pi\hbar)^3}{m}\int|\psi|^2\frac{d^3\mathbf{p}}{2p^0},
	\end{equation*}
	which is the right-hand side of Equation~\eqref{eq:charge}.
\end{proof}

\begin{remark}[Limits of the probabilistic reading]\label{rem:probability-limits}
	The density $j^0$ is positive and its integral is a multiple of the Wigner norm, so that it is independent of $t$ and invariant under the whole group. The factor $2p^0/m$ of Equation~\eqref{eq:euclidean-norm} combines with the Jacobian $1/4(p^0)^2$ of the Fourier measure into the invariant measure $d\mu=d^3\mathbf{p}/2p^0$. This does not make $j^0(x)$ the density of a position measurement of the particle. The construction of this chapter introduces no position observable. In the Galilean case $\mathbf{K}/m$ played this role by Remark~\ref{rem:weyl}, and the corresponding relativistic question is not treated here. The positive energy hypothesis $\mathbf{H3}$ restricts the amplitudes to $\mathcal{O}_m^+$. Solutions supported on $\mathcal{O}_m^-$ would yield a second irreducible representation, which does not appear in $\mathcal{F}_D$.
\end{remark}

The results of this section close the chain of reductions. They allow the proof of the uniqueness statement in Theorem~\ref{thm:main}, which collects the minimality of the Dirac space, the uniqueness of the parity matrix and of the rest subspace, and the identifications above. A local realization is understood as in Definition~\ref{def:local-realization} with the scalar product on its fields for which $J_D:\psi\mapsto D(A_p)\iota\psi(p)$ is unitary.

\begin{proof}[Proof of Theorem~\ref{thm:main}, part (B)]
	Let $\Phi$ be an elementary relativistic system of mass $m>0$ with $j=\tfrac12$. By part (A), $\mathbb{P}\mathcal{H}\cong\mathbb{P}\mathcal{H}_{m,1/2}$ through $\mathbb{P}\mathcal{W}$. Let $(D,W,\iota)$ be a local covariant realization of $(d_{1/2},\mathbb{C}^2)$ with a parity matrix $M$ ($\mathbf{H4}$, $\mathbf{H5}$) and $\dim W$ minimal ($\mathbf{H6}$).
	
	\emph{Step 1, the representation.} By Lemma~\ref{lem:dirac-space} and the converse in Proposition~\ref{prop:rest-selection}, the Dirac space satisfies $\mathbf{H4}$ and $\mathbf{H5}$, so the minimal dimension is at most $4$. Proposition~\ref{prop:minimal-extension} yields $\dim W\geq4$ and, for $\dim W=4$, a linear isomorphism $T_0:W\to\mathbb{C}^4$ with $T_0D(g)T_0^{-1}=D_{\mathrm{Dir}}(g):=\operatorname{diag}(g,(g^\dagger)^{-1})$.
	
	\emph{Step 2, the parity matrix.} Set $M_0:=T_0MT_0^{-1}$. Equation~\eqref{eq:parity-matrix} and item 1 of Lemma~\ref{lem:dirac-space} yield $M_0D_{\mathrm{Dir}}(g)M_0^{-1}=\gamma^0D_{\mathrm{Dir}}(g)\gamma^0$, so that $\gamma^0M_0$ commutes with $D_{\mathrm{Dir}}$. By item 1 of Proposition~\ref{prop:parity-problem} the two summands are inequivalent and irreducible, and Schur's lemma yields $\gamma^0M_0=\operatorname{diag}(c_1\mathbb{1},c_2\mathbb{1})$. Then $M_0^2=\mathbb{1}$ forces $c_1c_2=1$. With $S:=\operatorname{diag}(c_2\mathbb{1},\mathbb{1})$, which commutes with $D_{\mathrm{Dir}}$, one has $S\gamma^0S^{-1}=M_0$. This is Lemma~\ref{lem:parity-matrix-unique}. The isomorphism $T:=S^{-1}T_0$ satisfies $TD(g)T^{-1}=D_{\mathrm{Dir}}(g)$ and $TMT^{-1}=\gamma^0$.
	
	\emph{Step 3, the rest subspace.} The map $T\iota$ intertwines $d_{1/2}$ with $D_{\mathrm{Dir}}|_{SU(2)}$, and $\gamma^0T\iota(\mathbb{C}^2)=TM\iota(\mathbb{C}^2)=T\iota(\mathbb{C}^2)$ by $\mathbf{H5}$. Proposition~\ref{prop:rest-selection} yields $\epsilon\in\{\pm1\}$ and $c\in\mathbb{C}^\times$ with $T\iota=c\,\iota_\epsilon$ and $TF(\mathring{p})=\ker(\gamma^0-\epsilon)$.
	
	\emph{Step 4, equivalence.} By Step 2 and Step 3, $TJ_D\psi=cJ_{\mathrm{Dir}}\psi$ and $T\mathcal{U}_D(a,g)=\mathcal{U}_{\mathrm{Dir}}(a,g)T$, so that $T$ maps the fields of the given realization onto $\mathcal{F}_D$ of Theorem~\ref{thm:dirac}, and $T/|c|$ is unitary for the scalar products making $J_D$ and $J_{\mathrm{Dir}}$ unitary. It descends that $\mathbb{P}T\circ\mathbb{P}(J_D\mathcal{W})=\mathbb{P}(J_{\mathrm{Dir}}\mathcal{W})$, so the two realizations are equivalent, and a rescaling of $\iota$ does not change the rays. The sign $\epsilon$ labels the two parity choices and is fixed to $+1$ by convention, so the realization is unique for each $\epsilon$.
	
	\emph{Step 5, the Dirac realization.} By Theorem~\ref{thm:dirac}, the states are the rays of the solutions of $(\gamma_\mu p^\mu-m)\Psi=0$ and Poincar\'e acts through $D$. By Theorem~\ref{thm:covariant-norm} and Proposition~\ref{prop:kahler-dirac}, the scalar product is $\langle\cdot,\cdot\rangle_D$ and $F=\mathbb{P}(J\mathcal{W})$ is a realization of $\Phi$. By Proposition~\ref{prop:dirac-flow} and the remark following it, $F$ carries the Hamiltonian flow of $\langle H\rangle$ onto the Dirac evolution.
\end{proof}
\end{document}